% !TEX program = xelatex
\documentclass[11pt,a4paper]{article}

% ── 字体与中文支持 ──
\usepackage{fontspec}
% 使用系统已有的中文字体；无需额外安装 xeCJK。
\setmainfont{Noto Serif CJK SC}
\setsansfont{Noto Sans CJK SC}
\setmonofont{DejaVu Sans Mono}[Scale=0.86]
\XeTeXlinebreaklocale "zh"
\XeTeXlinebreakskip=0pt plus 1pt

% ── 页面与排版 ──
\usepackage[margin=2.5cm,headheight=16pt]{geometry}
\usepackage{titlesec}
\usepackage{enumitem}
\usepackage{booktabs}
\usepackage{array}
\usepackage{tabularx}
\usepackage{colortbl}
\usepackage{longtable}
\usepackage{graphicx}
\usepackage{float}
\usepackage{listings}
\lstset{basicstyle=\ttfamily\footnotesize,columns=fullflexible,keepspaces=true,breaklines=true,breakatwhitespace=false,showstringspaces=false,aboveskip=8pt,belowskip=8pt,xleftmargin=4pt,xrightmargin=4pt}
\setcounter{tocdepth}{2}
\usepackage{needspace}
\renewcommand{\contentsname}{目录}
\renewcommand{\figurename}{图}
\renewcommand{\tablename}{表}
\renewcommand{\abstractname}{摘要}
\setlength{\parindent}{2em}
\setlength{\parskip}{3pt}
\setlength{\tabcolsep}{4.5pt}
\linespread{1.08}
\widowpenalty=10000
\clubpenalty=10000
\raggedbottom
\emergencystretch=1.5em
\usepackage{caption}
\usepackage[hidelinks]{hyperref}
\hypersetup{pdftitle={面向资产管理的多 Agent 研发系统 v1.21},pdfsubject={工作原理、总体架构与实施配置},pdfauthor={}}
\usepackage{fancyhdr}
\usepackage{xcolor}
\usepackage{tikz}
\usetikzlibrary{shapes.geometric, arrows.meta, positioning, calc, fit, backgrounds}

% ── 颜色定义 ──
\definecolor{primaryblue}{RGB}{25,84,166}
\definecolor{accentorange}{RGB}{230,126,34}
\definecolor{lightgray}{RGB}{242,242,242}
\definecolor{darkgray}{RGB}{70,70,70}
\definecolor{softgreen}{RGB}{39,174,96}
\definecolor{softblue}{RGB}{52,152,219}
\definecolor{softpurple}{RGB}{142,68,173}
\definecolor{amber}{RGB}{243,156,18}
\definecolor{softred}{RGB}{192,57,43}


% 3.3 调用图共用样式：请求、返回与事件采用不同线型。
\tikzset{
  callhead/.style={rounded corners=3pt,draw=primaryblue!45,fill=primaryblue!7,
    font=\small\bfseries\sffamily,align=center,minimum height=.63cm,inner sep=4pt},
  callstep/.style={rounded corners=3pt,draw=primaryblue!40,fill=white,
    font=\footnotesize\sffamily,align=center,inner sep=5pt,minimum height=.55cm},
  callrequest/.style={-{Stealth[length=2mm]},draw=primaryblue,line width=.85pt},
  callreturn/.style={-{Stealth[length=2mm]},draw=darkgray!85,dashed,line width=.8pt},
  callevent/.style={-{Stealth[length=2mm]},draw=accentorange!90!black,densely dotted,line width=1pt},
  calllife/.style={draw=darkgray!25,densely dashed,line width=.6pt},
  calllabel/.style={font=\footnotesize\sffamily,fill=white,inner sep=2.2pt,align=center},
  callnote/.style={font=\footnotesize\sffamily,align=left,text=darkgray}
}

% ── 标题格式 ──
\titleformat{\section}{\Large\bfseries\sffamily\color{primaryblue}}{\thesection}{1em}{}[\titlerule]
\titleformat{\subsection}{\large\bfseries\sffamily\color{darkgray}}{\thesubsection}{1em}{}
\titleformat{\subsubsection}{\normalsize\bfseries\sffamily\color{darkgray}}{\thesubsubsection}{1em}{}

% ── 页眉页脚 ──
\pagestyle{fancy}
\fancyhf{}
\fancyhead[L]{\small\sffamily\color{darkgray} 多 Agent 研发系统架构方案}
\fancyhead[R]{\small\sffamily\color{darkgray} v1.21 \textbar\ 2026-10-13}
\fancyfoot[C]{\small\sffamily\color{darkgray} \thepage}
\renewcommand{\headrulewidth}{0.4pt}

% ── 表格列类型 ──
\newcolumntype{L}{>{\raggedright\arraybackslash}X}
\newcolumntype{P}[1]{>{\raggedright\arraybackslash}p{#1}}

% ── 文档信息 ──
\title{%
  \vspace{-1.5cm}
  {\sffamily\color{primaryblue}\rule{\textwidth}{1.5pt}}\\[0.7cm]
  {\Huge\sffamily\bfseries\color{primaryblue} 面向资产管理的\\[0.25cm] 多 Agent 研发系统}\\[0.5cm]
  {\large\sffamily\color{darkgray} 工作原理、总体架构与实施配置}\\[0.5cm]
  {\sffamily\color{primaryblue}\rule{\textwidth}{1.5pt}}
}
\author{}
\date{}

\begin{document}

\maketitle
\thispagestyle{empty}

\vspace{0.8cm}
\begin{center}
\renewcommand{\arraystretch}{1.5}
\begin{tabular}{ll}
\textbf{版本} & v1.21 \\
\textbf{日期} & 2026-10-13 \\
\textbf{文档性质} & 原理、总体设计与实施配置方案\\
\textbf{适用范围} & 面向资产管理的研发系统；投研交易系统另行设计 \\
\textbf{组织依据} & 《组织架构与研发协作规范》v1.2（2026-10-08）\\
\end{tabular}
\end{center}

\vspace{1cm}
\begin{abstract}
\noindent 本方案面向个人资产管理所需的研究、算法与软件能力建设，说明多 Agent 研发系统的工作原理、总体架构和实施配置。人负责提出意图、约束与关键批准，Agent 在授权范围内承担研究、开发、测试和运维任务。系统通过任务记录、固定版本、隔离执行与验证证据支撑协作，使成果能够审阅、交接和追溯。

建议由自有平台管理项目、任务与成果，以应用自行部署的 Agents SDK 承载智能协作，通过 Codex SDK 接入开发执行；上述组件可经适配与验证后替换。GitHub 私有仓库作为程序、组织规则和项目文档的权威版本核心。Mac 承担自建平台控制与主要状态记录，dot 云电脑作为独立的轻量研发和多 Agent 协作环境，经 GitHub 审阅交接成果；RTX 4090 承担数据处理与训练实验，T640 在网络和模型适配通过验证后提供可选自托管推理，阿里云承载一般 Web 服务及轻量接口验证。过程数据与归档以本地受控存储为主，各节点保留必要的运行状态和日志。

本方案将组织职责落实为任务、组件、接口、数据管理与验收条件，并选定供实施验证的软件基线。自建系统内部的自动交接是建设目标；dot 与自建平台的自动任务接入仍须另行验证。本文是待实施的架构与配置方案，不表示软件已经安装、接口已经连通或部署已经获准。组织职责以《组织架构与研发协作规范》v1.2 为依据；具体研发项目、资产运营及实盘交易系统另行设计。
\end{abstract}

\newpage
\begingroup
\setlength{\parskip}{0pt}
\linespread{1.0}\selectfont
\tableofcontents
\endgroup
\newpage

% ════════════════════════════════════════════════
% ════════════════════════════════════════════════
\section{背景、组织依据与方案范围}
% ════════════════════════════════════════════════

\subsection{项目背景}

个人资产管理需要持续积累研究、算法、数据处理与软件运维能力。相关研发工作涉及需求澄清、方案设计、实现、测试和交付，专业分工不会因为由个人主导而消失。有限精力若长期用于重复执行和环节协调，就会压缩研究判断与能力建设的时间。

研究、开发、测试和运维有不同的工作重点。同一人可以承担多个角色，但频繁切换会增加遗漏、交接不清与自我验证的风险。引入多个 Agent 可以分担执行工作，却不会仅凭角色名称形成独立验证：仍须固定输入与候选版本，隔离执行上下文和输出目录，并由测试角色依据验收标准留存可复核的证据。本项目因此既需要可委派的执行能力，也需要支撑版本管理、职责分工和结果检查的技术系统。

\subsection{解决思路与人机分工}

人类负责人提出意图、目标与约束，作出重要取舍和关键批准，并审阅交付结果；Agent 在获准范围内承担具体工作。自建研发系统的设计目标是通过程序完成任务拆解、依赖管理、执行调度和内部交接，减少人逐项转交任务的负担。遇到授权不足、重大变更或无法自行消除的阻塞时，系统应保留状态并提交人处理。

协作按版本与证据衔接。GitHub 私有仓库作为程序、组织规则和项目文档的权威版本核心，Mac 上的自建平台维护主要任务状态与执行协调。dot 云电脑是独立的轻量研发环境，在其支持的任务内开展多 Agent 协作，以固定版本取得输入，输出候选修改与检查记录，经审阅完成交接。共享仓库不意味着共享会话或调度程序，也不表示 dot 已接入自建平台的自动任务队列。

这一分工要求系统同时具备两类能力：一是在关键节点暂停、请求批准和恢复执行，使权限边界落到具体操作；二是记录责任任务、输入版本、执行结果及验证依据，使“已完成”能够由产物与证据核实，而非仅凭对话判断。

\subsection{组织依据}

本方案依据《组织架构与研发协作规范》v1.2（2026-10-08）展开技术设计。该规范确定投研与策略部、研发部及其软件开发、算法、运维、独立测试四组的职责，并列出人类负责人、研发协调负责人、开发组负责人和各组负责人的责任边界。财务、风险管理、合规等为预留职能，不表示本方案已经建设或运行这些职能。

规范规定规划、方案、研发、集成、验收、交接、反馈七个阶段，以及跨组接收确认、缺陷归属和规格调整规则；科研任务的否定结论也可按约定作为验收结果。上述组织规则由独立规范文档维护，本方案说明如何用任务记录、Agent 执行、版本成果和验证机制支撑其落实，不替代或重新定义组织章程。

\subsection{方案范围与文档结构}

本文覆盖研发系统的 Agent 原理、协作与通信、模型和工具调用、任务与成果管理、恢复机制及部署配置。第二、三章说明单个 Agent 与多 Agent 的一般工作机制；第四章提出本项目的职责、组件和端到端流程；第五章选定建议软件基线，并给出接口、数据结构、部署、分阶段验收与费用预算方案。附录 D 提供部署形态参考，不替代第五章的项目实施基线。

第五章的选型属于本方案范围，具体版本、性能、许可、安全配置与运维条件仍须在实施中验证和锁定；文中配置与命令不构成生产部署授权。设备分工服务于这一基线：Mac 负责控制，dot 承担轻量研发，4090 承担训练实验，T640 保留后续推理扩展，阿里云用于一般 Web 服务及轻量接口验证。过程数据以本地保存和归档为主，必要的节点状态与日志按各环境的实际条件管理。

资产运营和实盘交易系统的内部实现不在本文范围内。具体研究模型、策略、数据处理管线或 Web 项目，应另立需求与方案设计文档，并按本系统的研发流程形成交付物；本文不替这些项目预设业务功能或验收指标。

\section{单个 Agent 的组成与工作原理}
\label{sec:single-agent}
% ════════════════════════════════════════════════

单个 Agent 是围绕任务运行的软件单元。模型提供理解、推理与决策能力，工具和执行环境落实操作，交互接口接收要求并交付结果，存储与检索机制留存和取用信息。执行引擎把这些能力组织为可检查、可暂停并在条件具备时继续的过程。本章先解释各部分是什么，再说明它们怎样共同完成一次任务；多 Agent 的组织、通信与协作另在第三章展开。

\subsection{从模型到 Agent}
\label{sec:single-agent-concepts}

\subsubsection{大语言模型与 Agent 的区别}

大语言模型接收上下文并生成输出，可以解释资料、比较方案、编写代码，也可以提出结构化的工具调用请求。一次模型调用主要完成从输入到输出的计算。模型写出分析代码、文件路径或“已经完成”，并不说明代码实际运行、文件确实存在或外部系统已经发生改变；这些事实需要由真实执行结果确认。

Agent 在模型之外加入任务、工具、状态和反馈，使一次工作可以跨越多次调用。运行程序把当前任务与必要信息送入模型，接收其提出的行动或候选结论，在适用检查通过后执行，再把观察结果用于后续判断。简单任务可能一次调用即可结束，复杂任务则可能经历读取、计算、检查和调整。持续推进依赖这套运行机制，不能仅靠为模型指定一个岗位名称实现。

这里的“单个”指一份具有明确任务归属和运行状态的执行实例，并不规定它必须独占一个模型、进程或计算节点。模型服务、工具程序和数据存储可以位于不同位置；只要执行引擎能够按约定调用并记录它们，仍可构成同一个 Agent 的工作过程。

\subsubsection{脑、手、口耳与记忆的功能全景}

本方案沿用“脑、手、口耳与记忆”四个称呼，以便从功能上理解 Agent。“脑”以模型为核心，结合任务与证据判断问题的含义、可采用的方法和下一步动作；“手”由工具、脚本和计算程序构成，读取资料、处理数据或修改获准的对象。二者通过调用请求与观察结果相连：判断提出动作，执行提供事实。

“口耳”是接收与传递信息的接口，承接任务、补充说明、进展通知和最终交付；它可以是对话界面、API 或其他受支持的入口。“记忆”是保存与再次取用信息的机制，保留相关事实、结论和历史记录，并在需要时将它们加入当前上下文。消息已送达、信息已保存与模型已经理解，属于不同环节，均需要具体程序支持。

四个称呼并不要求四项独立服务。同一个文件工具既可作为“手”读取数据，也可服务于“记忆”的检索；同一交互接口既能接收新任务，也能接收授权决定。执行引擎负责把功能接起来，并落实权限、状态和资源限制。图~\ref{fig:agent-function} 表达的是这种功能关系，而非本地或云端的部署位置。

\begin{figure}[H]
\centering
% 图一为可编辑矢量概念图；图标与边界均由 TikZ 原生路径构成。
\begin{tikzpicture}[
  x=1cm,y=1cm,
  every node/.style={font=\small\sffamily,text=darkgray,inner sep=2pt},
  outline/.style={line width=1.1pt,line cap=round,line join=round},
  flow/.style={-{Stealth[length=5pt,width=4pt]},line width=1.1pt},
  capflow/.style={{Stealth[length=4pt,width=3.5pt]}-{Stealth[length=4pt,width=3.5pt]},line width=.9pt,draw=darkgray!65},
  labelmain/.style={font=\small\bfseries\sffamily},
  labelsub/.style={font=\footnotesize\sffamily,text=darkgray!85}
]
\path[use as bounding box] (-7.25,-5.15) rectangle (7.25,3.6);

% 脑徽记只表示推理能力，不表示云端或本地。
\draw[outline,draw=primaryblue!30,fill=primaryblue!5] (0,2.44) circle (1.04);

% 脑：左右脑轮廓、脑沟与神经连接。
\begin{scope}[shift={(0,2.44)},scale=.90]
  \path[outline,draw=primaryblue,fill=white]
    (0,.60)
    .. controls (-.23,.88) and (-.65,.72) .. (-.66,.44)
    .. controls (-1,.44) and (-1.10,.02) .. (-.84,-.16)
    .. controls (-1.03,-.50) and (-.71,-.78) .. (-.45,-.62)
    .. controls (-.24,-.86) and (0,-.65) .. (0,-.43)
    .. controls (0,-.65) and (.24,-.86) .. (.45,-.62)
    .. controls (.71,-.78) and (1.03,-.50) .. (.84,-.16)
    .. controls (1.10,.02) and (1,.44) .. (.66,.44)
    .. controls (.65,.72) and (.23,.88) .. (0,.60) -- cycle;
  \draw[outline,draw=primaryblue] (0,.60)--(0,-.43);
  \draw[draw=primaryblue,line width=.8pt,line cap=round]
    (-.62,.39).. controls (-.29,.41) and (-.36,.11)..(-.43,-.04)
    (-.83,-.15)..controls(-.64,.05)and(-.28,-.21)..(-.34,-.45)
    (.62,.39)..controls(.29,.41)and(.36,.11)..(.43,-.04)
    (.83,-.15)..controls(.64,.05)and(.28,-.21)..(.34,-.45);
  \foreach \x/\y in {-.43/-.04,-.34/-.45,.43/-.04,.34/-.45}
    \fill[primaryblue] (\x,\y) circle (.045);
\end{scope}
\node[labelmain,text=primaryblue,anchor=west] at (1.36,2.65) {脑};
\node[labelsub,anchor=west] at (1.36,2.20) {大语言模型};

% 推理与执行形成闭环；工具、存储与通信由执行环境按授权调动。
\draw[flow,draw=softgreen!85!black] (-.70,-1.05) -- (-.70,1.63);
\node[align=right,anchor=east,labelsub,text=softgreen!70!black] at (-.95,.60) {上下文\\执行结果};
\draw[flow,draw=primaryblue] (.70,1.63) -- (.70,-1.05);
\node[align=left,anchor=west,labelsub,text=primaryblue] at (.95,.60) {结论\\调用请求};

% 通用控制器：齿轮表示编排，勾选表示执行前的授权校验。
\path[outline,draw=primaryblue,fill=primaryblue!5,rounded corners=8pt]
  (-1.62,-1.08) rectangle (1.62,-2.38);

\begin{scope}[shift={(-.84,-1.72)}]
  \foreach \a in {0,45,...,315}
    \path[fill=primaryblue,rotate=\a] (-.09,.30) rectangle (.09,.48);
  \fill[primaryblue] (0,0) circle (.36);
  \fill[white] (0,0) circle (.17);
  \fill[white] (.40,-.30) circle (.24);
  \draw[outline,draw=softgreen!85!black] (.26,-.30)--(.37,-.40)--(.58,-.16);
\end{scope}
\node[align=center,labelmain,text=primaryblue] at (.68,-1.72) {编排与\\执行环境};

% 各能力连回控制器；没有绕过授权校验的路径。
\draw[capflow] (-1.74,-1.81) -- (-3.50,-1.81);
\draw[capflow] (1.74,-1.81) -- (3.50,-1.81);
\draw[capflow] (0,-2.55) -- (0,-3.31);

% 手：终端与扳手，象征真正落地的工具操作。
\begin{scope}[shift={(-4.60,-1.75)}]
  \path[outline,draw=softgreen!70!black,fill=white,rounded corners=3pt]
    (-.97,-.52) rectangle (.65,.61);
  \draw[draw=softgreen!70!black,line width=.8pt] (-.97,.31)--(.65,.31);
  \foreach \x in {-.79,-.62,-.45} \fill[softgreen!70!black] (\x,.46) circle (.025);
  \draw[outline,draw=softgreen!75!black] (-.67,.08)--(-.43,-.10)--(-.67,-.28) (-.26,-.28)--(.03,-.28);
  \begin{scope}[shift={(.54,-.21)},rotate=-38]
    \path[outline,draw=softgreen!70!black,fill=softgreen!12]
      (-.13,-.72)--(.13,-.72)--(.13,.10)
      ..controls(.47,.23)and(.42,.64)..(.20,.77)
      --(.20,.43)--(-.20,.43)--(-.20,.77)
      ..controls(-.42,.64)and(-.47,.23)..(-.13,.10)--cycle;
    \fill[softgreen!70!black] (0,-.55) circle (.035);
  \end{scope}
\end{scope}
\node[labelmain,text=softgreen!70!black] at (-4.60,-2.89) {手};
\node[labelsub] at (-4.60,-3.29) {工具与脚本};

% 口耳：相向的对话气泡表示双向消息收发。
\begin{scope}[shift={(4.60,-1.75)}]
  \path[outline,draw=softpurple,fill=softpurple!9,rounded corners=3pt]
    (-.09,.22)--(.95,.22)--(.95,-.45)--(.74,-.45)
    --(.76,-.71)--(.44,-.45)--(-.09,-.45)--cycle;
  \path[outline,draw=softpurple,fill=white,rounded corners=3pt]
    (-.96,.66)--(.43,.66)--(.43,-.08)--(-.35,-.08)
    --(-.72,-.38)--(-.68,-.08)--(-.96,-.08)--cycle;
  \foreach \x in {-.64,-.29,.06} \fill[softpurple] (\x,.30) circle (.046);
  \draw[draw=softpurple,line cap=round,line width=.8pt] (.31,-.21)--(.69,-.21);
\end{scope}
\node[labelmain,text=softpurple] at (4.60,-2.89) {口耳};
\node[labelsub] at (4.60,-3.29) {消息收发};

% 记忆：数据库圆柱与书页，兼有留存记录和积累认知的含义。
\begin{scope}[shift={(0,-3.97)}]
  \path[outline,draw=accentorange!90!black,fill=amber!10]
    (-.77,.35)--(-.77,-.37)
    ..controls(-.77,-.67)and(.77,-.67)..(.77,-.37)
    --(.77,.35)--cycle;
  \draw[outline,draw=accentorange!90!black,fill=white] (0,.35) ellipse (.77 and .22);
  \draw[draw=accentorange!90!black,line width=.8pt]
    (-.77,.01)..controls(-.77,-.28)and(.77,-.28)..(.77,.01);
  \path[draw=accentorange!90!black,line width=.9pt,fill=white,line join=round]
    (.22,-.18)..controls(.42,-.12)and(.63,-.18)..(.77,-.28)
    ..controls(.91,-.18)and(1.12,-.12)..(1.32,-.18)
    --(1.32,-.67)..controls(1.10,-.62)and(.91,-.69)..(.77,-.78)
    ..controls(.63,-.69)and(.42,-.62)..(.22,-.67)--cycle;
  \draw[draw=accentorange!90!black,line width=.75pt] (.77,-.28)--(.77,-.78);
\end{scope}
\node[labelmain,text=accentorange!90!black,anchor=east] at (-1.15,-3.95) {记忆};
\node[labelsub,anchor=west] at (1.58,-3.95) {事实、历史与记录};

\end{tikzpicture}
\caption{Agent 的功能构成与执行环境}
\label{fig:agent-function}
\end{figure}

例如，一项研发任务要求检查某批实验数据的异常，并给出有依据的说明。“口耳”接收检查范围与交付要求，“脑”确定应查看哪些字段和统计量，“手”读取数据并运行分析，“记忆”保存输入版本、发现与结果位置。统计结果返回后，“脑”再判断异常是否得到解释，并通过“口耳”交付结论。以下各节用这一例子说明机制，不预设专用数据库、固定接口或特定框架。

\subsubsection{角色、任务、运行实例与执行引擎}

功能需要落实为程序能够读取、执行和维护的对象。本方案用 \texttt{AgentDefinition} 表示角色定义，\texttt{Task} 表示任务，\texttt{AgentRun} 表示运行实例，\texttt{Runtime} 表示执行引擎。这些名称便于说明，不是所有产品共同遵循的标准，也不要求实现时采用同名类。

\textbf{角色定义}描述长期职责、工作要求和能力配置，包括使用何种模型、可以调用哪些工具、能够读取哪些记忆，以及适用的权限和预算。例如，数据分析角色可以读取指定资料、运行统计工具并提交带证据的说明。提供给模型的指令与程序强制执行的边界应分别落实；“只读”要求还需由工具接口、账户权限或沙箱约束。角色配置宜保留版本，以便确认某次工作依据了什么要求；调整配置时，应明确其生效范围并留下记录，不能悄然改变正在运行任务的依据。

\textbf{任务}描述本次目标、输入及其版本、范围、约束、期望产物和验收条件。实验数据检查可以要求覆盖指定批次，列出异常判定依据，保存可复核的统计结果，并说明仍无法解释的部分。验收条件应能指导实际检查，而不只写“给出高质量分析”。研究也可能得到否定结论或证据不足的结论；只要满足约定的方法、证据和披露要求，便不应为得到肯定答案而无限重试。

\textbf{运行实例}把角色与本次任务关联起来，持有上下文、进度、待处理调用、消耗及产物引用。同一角色可以反复执行不同任务，各次运行仍需分别记录。状态应能区分运行中、等待工具、等待信息或授权，以及完成、失败和取消；否则系统难以判断某一步尚未开始、正在进行，还是已经产生结果。逻辑上的任务隔离也不天然等同于进程或机器隔离。

\textbf{执行引擎}组织上下文、发起模型调用、解析返回、调度工具并更新状态。它可以从一段受控循环程序起步，在等待、恢复等要求增加时再引入相应基础设施。模型接口承接“脑”，工具接口承接“手”，任务与消息接口承接“口耳”，存储与检索接口支撑“记忆”。现有分析函数通常可以通过包装接口接入，无需为接入 Agent 重写全部业务逻辑。

\begin{figure}[H]
\centering
\begin{tikzpicture}[
 x=1cm,y=1cm, every node/.style={font=\small\sffamily,text=darkgray,align=center},
 box/.style={draw=primaryblue!55,fill=primaryblue!4,rounded corners=5pt,line width=.85pt,minimum height=1.2cm},
 flow/.style={-{Stealth[length=2mm]},line width=.95pt,draw=primaryblue},
 both/.style={{Stealth[length=1.8mm]}-{Stealth[length=1.8mm]},line width=.9pt,draw=primaryblue},
 tinylabel/.style={font=\footnotesize\sffamily,fill=white,inner sep=2pt}
]
\path[use as bounding box] (-7.6,-6.65) rectangle (7.6,1.02);
\node[box,text width=5.8cm] (def) at (-3.9,.28) {\textbf{AgentDefinition · 角色定义}\\{\footnotesize 职责、模型策略、工具、记忆范围与限制}};
\node[box,text width=5.8cm] (task) at (3.9,.28) {\textbf{Task · 本次任务}\\{\footnotesize 目标、输入、约束、产物与验收条件}};
\node[box,text width=7.4cm,fill=softpurple!5,draw=softpurple!60] (run) at (0,-1.50) {\textbf{AgentRun · 运行实例}\\{\footnotesize 关联角色与任务，持有本次上下文、状态及记录}};
\draw[flow] (def.south)--(-3.9,-.66)--(-1.8,-.66)--(run.north -| -1.8,0);
\draw[flow] (task.south)--(3.9,-.66)--(1.8,-.66)--(run.north -| 1.8,0);
\node[box,text width=7.4cm,fill=softgreen!5,draw=softgreen!65!black] (runtime) at (0,-3.26) {\textbf{Runtime · 执行引擎}\\{\footnotesize 组织上下文、检查权限、调用能力、提交状态}};
\draw[both] (run.south)--(runtime.north);
\draw[draw=primaryblue,line width=.95pt] (runtime.south)--(0,-4.16) (-5.55,-4.16)--(5.55,-4.16);
\foreach \x in {-5.55,-1.85,1.85,5.55} {\draw[both] (\x,-4.16)--(\x,-4.70);}
% 四个原生矢量图标对应功能称呼。
\begin{scope}[shift={(-5.55,-5.15)},draw=primaryblue,line width=1pt]
 \path[fill=primaryblue!5] (0,.30)..controls(-.30,.55)and(-.70,.25)..(-.47,-.04)..controls(-.58,-.42)and(-.05,-.46)..(0,-.21)..controls(.05,-.46)and(.58,-.42)..(.47,-.04)..controls(.70,.25)and(.30,.55)..(0,.30)--cycle;
 \draw (0,.30)--(0,-.21) (-.43,.13)..controls(-.14,.14)and(-.16,-.10)..(-.28,-.18) (.43,.13)..controls(.14,.14)and(.16,-.10)..(.28,-.18);
\end{scope}
\begin{scope}[shift={(-1.85,-5.15)},draw=softgreen!75!black,line width=1pt]
 \draw[fill=white,rounded corners=2pt] (-.51,-.34) rectangle (.51,.34);
 \draw (-.51,.15)--(.51,.15) (-.32,.02)--(-.16,-.09)--(-.32,-.20) (-.02,-.20)--(.28,-.20);
\end{scope}
\begin{scope}[shift={(1.85,-5.15)},draw=softpurple,line width=1pt]
 \path[fill=softpurple!7,rounded corners=2pt] (-.50,.34)--(.50,.34)--(.50,-.15)--(-.08,-.15)--(-.32,-.37)--(-.29,-.15)--(-.50,-.15)--cycle;
 \foreach \x in {-.28,0,.28}{\fill[softpurple] (\x,.10) circle (.032);}
\end{scope}
\begin{scope}[shift={(5.55,-5.15)},draw=accentorange!90!black,line width=1pt]
 \path[fill=amber!10] (-.46,.22)--(-.46,-.25)..controls(-.46,-.47)and(.46,-.47)..(.46,-.25)--(.46,.22)--cycle;
 \draw[fill=white] (0,.22) ellipse (.46 and .15);
 \draw (-.46,-.03)..controls(-.46,-.23)and(.46,-.23)..(.46,-.03);
\end{scope}
\node at (-5.55,-5.94) {\textbf{脑} · 模型接口};
\node at (-1.85,-5.94) {\textbf{手} · 工具接口};
\node at (1.85,-5.94) {\textbf{口耳} · 消息接口};
\node at (5.55,-5.94) {\textbf{记忆} · 存储检索};
\node[font=\footnotesize\sffamily,text=darkgray!85] at (0,-6.49) {图中是软件对象与调用关系；一个对象不必对应一项独立部署的服务。};
\end{tikzpicture}
\caption{从功能比喻到软件对象的实现组成}
\label{fig:agent-implementation}
\end{figure}

图~\ref{fig:agent-implementation} 将能力配置、本次目标与实际运行分开，便于分别调整。更换模型不会自动改变任务的验收条件，增加一个工具也不会自动授予使用它的权限。组件替换后仍应验证输入输出、错误处理及任务效果，不能仅凭接口名称兼容便认定行为相同。

\subsection{理解、推理与决策：脑的作用}
\label{sec:agent-reasoning}

“脑”的作用在于把要求、资料和观察结果转化为可检验的判断。模型参与理解与选择，执行引擎则决定这些选择如何进入实际流程。下面区分模型已有能力、本次输入和反馈调整，避免把全部推理过程简化为一次 API 传输，也避免把模型的判断当作无需验证的事实。

\subsubsection{模型能力的来源与使用边界}

模型的基础能力来自训练。预训练通常利用大量数据学习语言结构、知识关联和常见问题模式；指令微调、偏好优化或强化学习等后训练方法，进一步影响模型遵循要求、解决问题和生成回答的方式。不同模型采用的数据与训练方法不同，能力覆盖、稳定性和适用范围也不同。InstructGPT 的研究说明人类示范与反馈可以改善指令遵循，同时保留了模型仍会犯错的边界\cite{instructgpt-training}。

在使用阶段，模型把输入转换为内部表示，并依据已训练的参数和当前上下文生成输出。它能够在上下文中利用任务说明、示例和新提供的资料，而不必为每个任务重新训练；GPT-3 的少样本研究展示了不更新参数、仅通过文本任务与示例进行适配的方式\cite{gpt3-fewshot}。普通调用中加入一份实验记录，改变的是本次判断可用的信息，不会因此自动把记录写进模型权重。检索记忆与重新训练也是两类不同机制。

训练所得知识适合提供背景和候选解释，却未必覆盖最新事实、私有数据或当前环境。对某批实验数据的实际分布、某次程序的返回值和某个文件的内容，应依靠本次提供的资料与真实工具结果。精确计算、大规模数据处理和可重复测试宜交给相应工具；模型负责选择方法、组织检查并解释结果。具体任务能达到什么水平，需要用代表性输入与验收标准验证，不能只按模型名称推断。

\subsubsection{指令、目标与上下文的作用}

指令说明怎样工作，目标说明本次希望完成什么，上下文提供当前判断所能使用的材料。角色指令可以要求引用来源、保留不确定性或遵循既定方法；任务目标限定对象、范围和交付；上下文还可包含输入片段、历史结论、工具说明、当前状态及最新观察。将这些内容区分清楚，有助于模型识别哪些是约束、哪些是证据、哪些只是待验证的假设。

上下文的选择会直接影响判断。缺少字段定义时，模型可能误解实验数据的含义；只看到汇总均值时，可能忽略局部异常；混入其他批次的旧结论时，又可能把不适用的经验当作依据。运行程序应围绕当前问题选取相关信息，保留来源、时间或版本及必要的不确定性，并在后续取得新证据时更新输入。由于可接收长度和内容类型受接口限制，长资料可通过筛选、分段读取或摘要组织，但关键判断仍应能追溯到原始证据。

指令也不能代替执行控制。提示词中的预算、权限或数据外发要求，应由相应程序和运行环境落实。网页、文件或工具返回中可能包含看似指令的文字，它们首先是任务资料，不能自行扩大授权。来源和约束冲突时，应按系统既定规则处理；模型无法确认的目标或重要取舍，可通过交互接口澄清。目标并非越长越好，清楚的边界和可检查的完成条件更有助于稳定执行。

\subsubsection{需求理解、任务规划与行动选择}

需求理解首先要把自然语言中的意图落到可检查的问题上。对“找出这批数据为什么异常”，模型需要辨明异常所指的指标、比较基准、数据范围和用户希望得到的产物，再区分已知条件与未证实假设。已有资料足够时可以继续；缺少某个可自行读取的定义时先补资料；只有会改变目标、重要取舍或授权且无法从现有信息确认的缺口，才需要向人提问。这一转换是带有不确定性的解释过程，并非确定性编译。

规划把目标分解为有依赖关系的检查或操作。例如，先确认数据结构和单位，才能选择合适的统计方法；先定位异常集中在哪些样本，再决定是否需要更细的比较。计划应表达关键步骤、依赖和预期检查依据，并结合工具可用性、时间与资源约束选择范围。简单任务可以直接处理，复杂任务可以逐步展开；并非每次都必须生成长计划，也不必在获得新信息之前确定所有后续细节。

行动选择是在当前证据和可用能力之间作取舍。模型可以选择直接回答、检索缺失资料、运行统计工具、读取更细的样本，或请求澄清。选择时应考虑某个动作能消除什么不确定性、是否依赖尚未完成的步骤，以及成本和副作用是否适合当前任务。工具名称与参数说明帮助模型提出请求，真正的可执行性和授权仍须由执行侧核实。把某个工具写进计划，不代表它已经可用或已获许可。

工程上需要的是可检查的计划、关键假设、选择依据与产物，而不是取得模型不可见的内部思维过程。对重要选择，可以要求模型简要说明采用什么方法、依据哪些事实、还存在哪些限制，并将这些外显信息与实际执行记录关联。这样既能支持复核，也能在方向错误时明确修正任务或方法，而不依赖模型对自身内部过程的完整叙述。

\subsubsection{结果解释、方案调整与判断局限}

工具返回之后，“脑”需要判断结果与任务目标的关系。一次统计返回高缺失率，只能支持相应样本存在缺失现象；是否由采集失败、字段不适用或清洗逻辑造成，还需其他证据。模型应区分观测事实、由事实支持的推断和尚待验证的解释，避免将相关性直接写成原因，也避免把预期结果代替实际结果。

新证据可能要求局部修正计划。若发现单位不一致，可以先核对口径；若工具报告输入格式错误，应修正输入或选择适配工具；若原假设与结果不符，应保留该结果并调整解释。调整应说明新增证据和改变的部分，以免在反复尝试中悄悄改变目标或验收标准。需要超出原任务范围、增加显著成本或取得新权限时，应回到相应决定与授权环节。

模型的流畅表达、详细计划和自我检查都不能保证正确。同一模型反复审阅自己的输出可能保留相同盲点，较长推理也可能围绕错误前提继续展开。因此，应尽可能把重要判断落到原始资料、确定性计算、实际测试或人工决策等适当依据上。证据不足时应报告缺口；方法存在适用限制时应说明边界。模型擅长组织和解释这些依据，但不能凭一句“检查通过”替代它们。

\subsection{行动、交互与信息留存}
\label{sec:single-agent-capabilities}

“脑”的判断要依靠“手”取得新事实，通过“口耳”与外部交换信息，并利用“记忆”保持必要的连续性。本节说明这些组成部分的职责与边界；它们在一次任务中的先后衔接见第~\ref{sec:single-agent-cycle} 节。

\subsubsection{工具与执行环境}

工具是 Agent 能够调用的具体能力，例如文件读取、数据库查询、代码运行、数据统计或外部服务访问。工具声明通常给出名称、用途、参数和返回约定；实际实现可以是函数、命令行程序、子进程或远程接口。声明帮助模型选择和组织请求，实现负责产生真实效果，二者需要相互对应。

执行环境决定工具在哪里运行、拥有什么依赖和权限，以及能够访问哪些文件、网络或计算资源。一个数据分析工具可以在本地处理整份数据，仅返回统计结果；一个远程查询工具则需要相应服务、账户与可达网络。工具接口存在并不保证运行条件已经具备，生成的代码也只有在适用检查通过并由运行器执行后，才会产生输出。

工具的结果应区分内容、状态与副作用。成功读取可以返回记录，统计程序可以返回指标和产物位置，错误应带有可用于定位问题的信息。修改文件、提交外部操作等动作还会改变系统状态，其超时和重试不能按纯读取操作处理。允许访问的对象、资源上限、取消方式和错误语义，都由具体工具与执行环境确定，框架不能自动替任意工具补足这些保证。

\subsubsection{输入输出与交互接口}

“口耳”使 Agent 能够接收任务和补充信息，也能提交进展、问题和交付结果。入口可以是聊天、网页、API 或命令行，内容可以是文字、文件以及接口支持的图像、音频等。界面接收到了附件，并不表示其内容已进入模型上下文；程序仍需确认文件可读取、内容如何解析，以及需要把哪一部分提供给当前判断。

交互应围绕当前任务保持关联。新任务、对原任务的补充、取消指令和对某项待决问题的回复具有不同含义，运行程序需把它们交给正确实例处理。例如，对某项操作的授权应对应明确对象和范围；与该问题无关的一条消息，不能被当作放行。消息收发和状态关联由程序落实，模型帮助解释内容并生成合适的回应。

输出可以是自然语言说明、结构化结果或带引用的产物。进展应反映已经观察到的事实，交付应让接收者取得所需文件并理解结论、依据和限制。消息发送成功不等于接收者已阅读，给出路径也不等于接收者能够访问。需要澄清、决定或新授权时，接口承担提问与接收回复；依赖该回复的工作进入等待，而不能把提出问题本身当成取得许可。

\subsubsection{外部数据、上下文与记忆}
\label{sec:single-agent-data}

Agent 并不规定业务数据必须采用某种格式或集中保存于某个数据库。文档、代码、表格、数据库记录、图片、音频和计算产物都可能成为输入或输出；它们可保存在文件系统、关系数据库、对象存储或外部服务中。框架或产品可以提供知识库与存储功能，但这些是具体实现选择。能否使用某类数据，还取决于工具支持、权限、网络可达性和模型接口的内容限制。表~\ref{tab:single-data-access} 列出常见的访问方式与约束。

\begin{table}[H]
\centering\small\renewcommand{\arraystretch}{1.24}
\begin{tabularx}{\textwidth}{P{2.75cm} P{4.2cm} L}
\toprule
\textbf{存储或服务} & \textbf{常见访问方式} & \textbf{实际约束示例}\\
\midrule
本地或远程文件 & 读取、写入、搜索、下载等工具 & 路径含义、文件权限、挂载范围和可用格式\\
\rowcolor{lightgray}
数据库 & 查询工具或业务服务接口 & 账户、表列权限、查询限制及数据库事务语义\\
对象存储 & 存取对象或解析对象引用 & 对象权限、访问凭据、版本和传输限制\\
\rowcolor{lightgray}
外部 API & 按接口提交请求并解析返回 & 认证、接口字段、限流、服务范围及数据用途\\
\bottomrule
\end{tabularx}
\caption{数据可以分布保存，访问边界来自具体接口、权限与运行环境}
\label{tab:single-data-access}
\end{table}

需要明确区分三件事：\textbf{外部数据存在，工具能够读取，模型实际收到。}宿主机上有文件，或者任务中出现文件路径，都不会使模型自动知道内容。工具可以返回全文、片段、统计摘要或引用；只有进入本次请求，或通过模型服务明确支持的机制实际加载的内容，才成为模型可利用的上下文。引用用于定位内容，其自身既不等于内容，也不授予读取权限。

在实验数据检查中，大数据表可以留在工具侧，由统计程序处理后向模型提供指标和必要样本。这减少上下文占用，但摘要可能遗漏细节，模型应能够在需要时请求补充读取。若模型部署在云端，将本地统计摘要或文件片段送入请求，也会产生相应数据传输；是否外发要检查实际字段、附件和检索结果，不能仅凭原文件保存在本地作判断。

“记忆”强调信息跨调用或跨任务的留存与再次取用。历史消息、经整理的结论和长期资料，可以通过显式读取、检索或摘要进入当前上下文；保存后是否可见、何时检索、如何处理过期信息，取决于应用策略与具体能力。记忆应保留必要来源、适用条件和更新时间，避免旧结论在条件变化后继续被使用。它提供可取用的信息，不会自动更新模型权重，也不意味着每次调用都会装入全部历史。

\subsubsection{运行状态、历史记录与恢复}
\label{sec:single-agent-recovery}

信息留存还要区分内容与控制用途。消息历史记录对话和工具往返，业务产物记录任务处理的文件、代码和结果；执行状态则说明当前步骤、变量、等待原因及待处理调用。检查点在实现支持并配置时保存某些可恢复状态。这几类信息可以共用基础设施，但保存聊天历史不能替代保存执行状态，也不能证明业务产物已经可靠落盘。

恢复能力取决于实际保存的范围。以内存保存的状态可能随进程结束而丢失，持久检查点也未必包含工具子进程、外部数据库或远程服务的变化。LangGraph 的 checkpointer 保存线程范围内的图状态，store 可用于跨线程应用数据；二者的用途及后端持久性需要分别理解\cite{langgraph-persist}。这是具体实现的例子，并不要求其他框架使用相同对象，也不能据此把检查点视为完整进程快照。

恢复前应核实原任务、输入版本、待处理调用和已产生的结果。若某项写入或外部提交已成功，而记录在中断前未更新，简单重放可能重复执行；若计算仍在运行，重新提交又可能浪费资源。某些暂停机制恢复时会重新执行节点的前部代码，LangGraph 的 interrupt 即要求注意此类副作用\cite{langgraph-interrupt}。因此，恢复需要结合实际结果查询、去重或其他适用保护，并重新确认引用可达、环境可用和授权有效。

运行记录服务于这一核对过程。它应在需要的范围内关联任务、模型请求、工具调用、返回结果和状态变化，使问题能够回到实际发生的环节。记录什么、保留多久以及是否包含敏感信息，应按应用要求配置；拥有日志不自动意味着能够恢复，能够恢复也不自动证明原结论正确。

\subsection{单个 Agent 的完整执行循环}
\label{sec:single-agent-cycle}

上述组成在运行时形成反馈循环：“口耳”取得要求，“记忆”提供相关信息，“脑”选择下一步，“手”执行获准动作，真实观察再进入后续判断。图~\ref{fig:single-agent-loop} 给出这一过程的概览。它表达职责与典型路径，并非要求每项任务都调用工具或执行固定轮数；等待、失败和取消也可能发生在中间步骤。

\begin{figure}[H]
\centering
\begin{tikzpicture}[
 x=1cm,y=1cm, every node/.style={font=\small\sffamily,text=darkgray,align=center},
 model/.style={draw=primaryblue!65,fill=primaryblue!6,rounded corners=5pt,text width=5.4cm,minimum height=1.10cm,line width=.9pt},
 engine/.style={draw=softgreen!70!black,fill=softgreen!5,rounded corners=5pt,text width=5.4cm,minimum height=1.10cm,line width=.9pt},
 flow/.style={-{Stealth[length=2mm]},line width=.95pt,draw=darkgray!80},
 note/.style={font=\footnotesize\sffamily,fill=white,inner sep=2pt}
]
\path[use as bounding box] (-7.6,-6.75) rectangle (7.6,1.20);
% 模型与执行两侧均有可编辑矢量图标。
\begin{scope}[shift={(-5.24,.86)},scale=.53,draw=primaryblue,line width=1.1pt]
 \path[fill=primaryblue!5] (0,.30)..controls(-.30,.55)and(-.70,.25)..(-.47,-.04)..controls(-.58,-.42)and(-.05,-.46)..(0,-.21)..controls(.05,-.46)and(.58,-.42)..(.47,-.04)..controls(.70,.25)and(.30,.55)..(0,.30)--cycle;
 \draw (0,.30)--(0,-.21) (-.43,.13)..controls(-.14,.14)and(-.16,-.10)..(-.28,-.18) (.43,.13)..controls(.14,.14)and(.16,-.10)..(.28,-.18);
\end{scope}
\begin{scope}[shift={(1.74,.86)},scale=.58]
 \foreach \a in {0,45,...,315}{\path[fill=softgreen!70!black,rotate=\a] (-.07,.25) rectangle (.07,.40);}
 \fill[softgreen!70!black] (0,0) circle (.31);
 \fill[white] (0,0) circle (.14);
\end{scope}
\node[font=\small\bfseries\sffamily,text=primaryblue] at (-3.9,.86) {模型判断};
\node[font=\small\bfseries\sffamily,text=softgreen!60!black] at (3.9,.86) {执行环境与真实工具};
\node[engine] (prep) at (3.9,-.15) {\textbf{1 接收与准备}\\{\footnotesize 任务、角色、上下文与运行状态}};
\node[model] (reason) at (-3.9,-.15) {\textbf{2 判断下一步}\\{\footnotesize 提出动作请求或候选结果}};
\node[engine] (act) at (3.9,-2.07) {\textbf{3 校验并执行动作}\\{\footnotesize 允许才执行；拒绝或等待需记录}};
\node[model] (observe) at (-3.9,-2.07) {\textbf{4 解读观察结果}\\{\footnotesize 根据证据调整，更新下一步判断}};
\node[model] (candidate) at (-3.9,-4.18) {\textbf{5 提交候选结果}\\{\footnotesize 结论、产物、证据及未解决事项}};
\node[engine] (verify) at (3.9,-4.18) {\textbf{6 检查结束条件}\\{\footnotesize 验收证据、状态、限制与待决事项}};
\draw[flow] (prep.west)--node[note,above]{上下文}(reason.east);
\draw[flow] (reason.south east)--node[note,above,sloped]{动作请求}(act.north west);
\draw[flow] (act.west)--node[note,above]{真实结果}(observe.east);
\draw[flow,draw=primaryblue] (observe.north)--node[note,left]{继续}(reason.south);
\draw[flow] (observe.south)--node[note,left]{可交付}(candidate.north);
\draw[flow,draw=primaryblue] (reason.west)--(-6.85,-.15)--(-6.85,-4.18)--(candidate.west);
\node[note,rotate=90,text=primaryblue] at (-6.85,-2.10) {无需工具时提交候选结果};
\draw[flow] (candidate.east)--node[note,above]{候选产物}(verify.west);
\draw[flow,draw=accentorange!90!black] (verify.east)--(7.45,-4.18)--(7.45,-5.39)--(-7.45,-5.39)--(-7.45,-.15)--(reason.west);
\node[note,text=accentorange!90!black] at (0,-5.39) {可修复且在限制内：带问题与证据进入下一轮};
\draw[flow] (verify.south)--(3.9,-5.91);
\node[engine,minimum height=.62cm,text width=7.1cm] at (2.8,-6.24) {完成 \quad / \quad 等待信息或授权 \quad / \quad 失败或取消};
\end{tikzpicture}
\caption{单个 Agent 的执行循环：模型判断与程序执行各负其责}
\label{fig:single-agent-loop}
\end{figure}

\subsubsection{接收任务与准备上下文}

执行引擎先把输入关联到任务和角色配置，建立本次运行状态，并核对目标、范围、验收条件和已有授权。对于实验数据检查，应确认所指批次与数据位置、需要解释的异常和交付要求。可在已授权范围内检查的数据与说明，先通过现有能力补齐；会影响目标或重要取舍而仍无法确认的内容，再通过交互接口提问。

随后，程序为下一次模型调用准备必要上下文，包括角色指令、任务目标、当前进度、相关资料或记忆、可用工具说明和已有观察。资料应携带必要来源与版本，长记录可整理为摘要和可追溯引用。这里的目标是让模型得到足以判断下一步的信息，而不是把所有文件、历史和无关任务复制进请求。

准备阶段还应落实实际输入边界，例如附件是否能解析、工具是否可用、哪些内容允许传给模型服务，以及本次运行的时间或资源限制。缺失信息可以形成明确的待补充项，外部依赖未满足可以形成等待或失败状态。后续步骤取得新资料时再更新上下文，无需在起点一次性收集所有可能用到的信息。

\subsubsection{调用模型与解释返回结果}

运行程序或适配层把准备好的输入转换为所选模型接口支持的请求，并按配置调用服务。输入可以包含文字、图像、音频等受支持内容，也可声明可用工具和输出格式。适配层承担格式转换、响应解析和可用的错误或用量信息处理，但不能消除不同模型服务在能力、流式事件、结构化输出及状态管理上的差异。

\begin{figure}[H]
\centering
\begin{tikzpicture}[x=1cm,y=1cm,font=\footnotesize\sffamily,text=darkgray]
\path[use as bounding box] (0,-7.38) rectangle (15.4,.38);
\foreach \x in {2.0,7.7,13.4} {\draw[calllife] (\x,-.38)--(\x,-6.80);}
\node[callhead,minimum width=3.4cm] at (2,0) {Agent 运行程序};
\node[callhead,minimum width=4.25cm] at (7.7,0) {上下文 / 模型适配器};
\node[callhead,minimum width=3.4cm] at (13.4,0) {模型服务};
\draw[callrequest] (2,-1.0)--node[calllabel,above]{目标、约束、已有结果}(7.7,-1.0);
\node[callstep,text width=4.35cm] at (7.7,-1.79)
  {\mbox{组装上下文与工具 schema}\\\mbox{按配置检查外发范围与限额}};
\draw[callrequest] (7.7,-2.88)--node[calllabel,above]{转换为供应商 / 自部署推理请求}(13.4,-2.88);
\node[callstep,text width=2.6cm] at (13.4,-3.63) {执行模型推理};
\draw[callreturn] (13.4,-4.45)--node[calllabel,above]{内容或工具调用请求；也可返回错误}(7.7,-4.45);
\node[callstep,text width=4.35cm] at (7.7,-5.19)
  {\mbox{解析与校验返回结构}\\\mbox{处理用量与错误（依接口）}};
\draw[callreturn] (7.7,-6.09)--node[calllabel,above,pos=.72]{模型输出或错误类别}(2,-6.09);
\node[callstep,text width=3.55cm] at (2,-6.81) {\mbox{解析下一步：继续 / 结束}\\\mbox{或进入工具执行控制}};
\node[callnote,text width=8.8cm,anchor=west] at (5.6,-6.81)
  {工具请求在此仍未执行；模型输出不等于成功证据。\\失败、超时与重试由调用方按实际配置处理。};
\end{tikzpicture}
\caption{模型服务调用示意：运行程序或适配层组织输入、解析返回；本例中的自定义工具仍由执行侧运行，见图~\ref{fig:tool-execution-control}。}
\label{fig:model-service-call}
\end{figure}

图~\ref{fig:model-service-call} 中，模型响应可能包含自然语言内容、结构化结果、一个或多个工具请求，也可能以错误或不完整响应结束。运行程序需要识别返回类型、检查必要结构，并决定进入工具执行、补充信息、候选结果检查或错误处理。形式有效只表示满足接口约定，尚不说明内容正确、参数对应正确对象或目标已经完成。

若返回自定义函数调用请求，此时尚未发生工具动作。OpenAI 函数调用文档展示了模型提出调用、应用执行函数、应用回传结果的职责分工\cite{openai-function}。有些服务也提供由提供方执行的托管工具，应另外核实其执行位置、数据接收方和返回语义。模型接入中的超时、调用次数与严格费用上限，同样需要程序实际落实，不能只依赖提示词要求。

\subsubsection{执行工具与回传观察结果}

执行侧接到请求后，解析工具名称和参数，核对目标对象、运行身份、任务范围及适用限制，再决定允许、等待授权或拒绝。参数符合声明格式，只完成了其中一项检查。检查可以位于应用代码、框架钩子、工具服务、沙箱或数据库，并不要求统一网关；但凡需要强制约束，都应在真正执行的路径上生效。

\begin{figure}[!htb]
\centering
\begin{tikzpicture}[x=1cm,y=1cm,font=\footnotesize\sffamily,text=darkgray]
\path[use as bounding box] (0,-10.59) rectangle (15.4,.38);
\foreach \x in {2.0,5.8,9.6,13.4} {\draw[calllife] (\x,-.38)--(\x,-9.65);}
\node[callhead,minimum width=3.35cm] at (2,0) {运行程序 / 上下文};
\node[callhead,minimum width=3.35cm] at (5.8,0) {执行侧检查与调度};
\node[callhead,minimum width=3.35cm] at (9.6,0) {工具执行环境};
\node[callhead,minimum width=3.35cm] at (13.4,0) {数据 / 服务 / 存储};
\draw[callrequest] (2,-.90)--node[calllabel,above]{请求、身份与上下文}(5.8,-.90);
\node[callstep,text width=3.2cm] at (5.8,-1.49)
  {核对参数、身份与范围\\落实适用执行规则};
% 三个互斥分支分别成带；数据访问只发生在获准的实际执行路径上。
\draw[primaryblue!35,rounded corners=3pt] (.15,-2.02) rectangle (15.25,-5.65);
\node[calllabel,anchor=north west,text=primaryblue,font=\footnotesize\bfseries\sffamily] at (.25,-2.07) {允许};
\draw[callrequest] (5.8,-2.58)--node[calllabel,above]{执行点复核后调用}(9.6,-2.58);
\draw[callrequest] (9.6,-3.08)--node[calllabel,above]{按权限读取输入}(13.4,-3.08);
\draw[callreturn] (13.4,-3.59)--node[calllabel,above]{数据片段 / 引用}(9.6,-3.59);
\draw[callrequest] (9.6,-4.10)--node[calllabel,above]{按任务需要写入产物}(13.4,-4.10);
\draw[callreturn] (13.4,-4.61)--node[calllabel,above]{写入回执 / 产物引用}(9.6,-4.61);
\draw[callreturn] (9.6,-5.27)--node[calllabel,above]{真实结果、错误、日志与引用；关联本次调用}(2,-5.27);
\node[callnote,text width=6.6cm,align=left,anchor=west] at (.55,-3.91)
  {工具按获准范围访问数据或服务。\\只读工具无需写入；不涉及外部数据的工具\\可直接返回执行结果。};
\draw[accentorange!65,rounded corners=3pt] (.15,-5.89) rectangle (15.25,-8.58);
\node[calllabel,anchor=north west,text=accentorange!90!black,font=\footnotesize\bfseries\sffamily] at (.25,-5.94) {待审批};
\draw[callreturn] (5.8,-6.46)--node[calllabel,above]{暂停并提出待决事项}(2,-6.46);
\draw[callrequest] (2,-7.06)--node[calllabel,above]{取得决定后恢复调用}(5.8,-7.06);
\node[callstep,text width=2.9cm,draw=accentorange!65] at (5.8,-7.63)
  {复核授权与操作对象\\核实当前状态};
\draw[callrequest] (5.8,-8.31)--node[calllabel,above,pos=.8]{复核通过才执行}(9.6,-8.31);
\node[callnote,text width=5.8cm,anchor=west] at (8.6,-7.06)
  {依赖决定的动作保持未执行。\\否决或复核失败，转入拒绝分支。};
\draw[softred!45,rounded corners=3pt] (.15,-8.82) rectangle (15.25,-9.75);
\node[calllabel,anchor=north west,text=softred,font=\footnotesize\bfseries\sffamily] at (.25,-8.87) {拒绝};
\draw[callreturn,draw=softred] (5.8,-9.40)--node[calllabel,above]{返回拒绝原因；未执行}(2,-9.40);
\node[callnote,anchor=west,text width=6.0cm] at (8.6,-9.30) {检查拒绝 / 审批否决 / 复核失败};
\node[callnote,anchor=west,text width=14.5cm] at (.35,-10.22)
  {运行程序保留调用与结果关联，只把所需且允许外发的观察纳入下一次模型输入。\\模型实际得到的是回传的数据片段、摘要或引用；持有引用不等于已读取完整数据。};
\end{tikzpicture}
\caption{工具执行与数据交互示例：数据由获准的工具读取或按需写入，观察经运行程序回传。人工等待为可选分支，检查与恢复须由实际实现支持。}
\label{fig:tool-execution-control}
\end{figure}

在图~\ref{fig:tool-execution-control} 的允许分支，执行环境实际调用函数、启动程序或访问外部服务，收集返回内容、错误、日志和产物引用。对实验数据检查，模型提出统计请求后，必须等统计程序产生真实指标，才能据此解释异常。运行程序保留调用与结果之间的关联，并按模型接口要求回传观察；结果过大时可提供摘要和引用，必要时再读取细节。

图中数据或服务由获准工具访问，读写按任务选择。模型实际看到的是按外发边界回传的内容；工具请求不等于直接访问数据库，持有引用也不等于已读取完整数据。

待授权分支应先保存必要状态并提出具体问题，依赖该决定的动作保持未执行。获得决定后还要复核操作对象、范围和当前状态，通过才继续；否决或复核失败则进入拒绝分支。若已有明确授权且仍在有效范围内，可直接按允许路径处理。等待机制只有在应用、框架或运行环境实际支持并配置时才成立，不是所有工具或平台自动具备的能力。

错误、拒绝和等待都应作为真实观察或状态记录，不能改写为成功。一次调用失去响应时，外部动作可能仍在运行或已经完成，应按工具能力查询结果再决定后续处理。较长计算也可能先返回作业标识，再由程序查询或接收完成事件；收到“已提交”只说明作业被接受，不说明结果已就绪。后续调用应使用实际观察，不得用模型预期的输出代替。

\subsubsection{检查、调整、重试与结束}

模型取得观察后，依据第~\ref{sec:agent-reasoning} 节所述方法解释结果，判断是否补充读取、调整方法或提交候选产物。执行引擎将候选结果与验收条件对照：程序可检查文件、结构或测试，模型可辅助核对论证与遗漏，必要的重要判断交由人处理。检查结论只覆盖实际检查的内容，格式合格或模型自评不能证明业务结论成立。

图~\ref{fig:single-agent-verification} 展开检查与后续分支。验证工具和人均按需要参与，三个分支并非固定步骤；等待、重试与取消能力须由实际实现落实，并非各框架自动具备。

\begin{figure}[!htb]
\centering
\begin{tikzpicture}[x=1cm,y=1cm,font=\footnotesize\sffamily,text=darkgray]
\path[use as bounding box] (0,-10.02) rectangle (15.4,.38);
\foreach \x in {2.0,5.8,9.6,13.4} {\draw[calllife] (\x,-.38)--(\x,-9.9);}
\node[callhead,minimum width=3.35cm] at (2,0) {运行程序 / 控制};
\node[callhead,minimum width=3.35cm] at (5.8,0) {模型判断};
\node[callhead,minimum width=3.35cm] at (9.6,0) {验证工具 / 测试};
\node[callhead,minimum width=3.35cm,draw=accentorange!65,fill=amber!6] at (13.4,0) {人类（按需参与）};
\draw[callreturn] (5.8,-1.0)--node[calllabel,above]{候选产物、依据与缺口}(2,-1.0);
\draw[callrequest] (2,-1.62)--node[calllabel,above]{可程序化的检查请求（按任务需要）}(9.6,-1.62);
\node[callstep,text width=2.9cm] at (9.6,-2.26)
  {确定性校验 / 测试\\记录实际检查与结果};
\draw[callreturn] (9.6,-2.95)--node[calllabel,above]{检查结果、失败项与可核查证据}(2,-2.95);
\node[callstep,text width=2.9cm] (gate) at (2,-3.66)
  {对照验收条件与限制\\核对证据覆盖范围};
\node[callnote,text width=8.4cm,anchor=west] at (5.8,-3.66)
  {模型自评不等于验收通过；结论仍需相应证据。\\检查通过仅覆盖已检查的对象与条件。};
% 下面三带表示条件分支，不是每次运行都要顺序经过的固定阶段。
\draw[primaryblue!35,rounded corners=3pt] (.15,-4.29) rectangle (15.25,-6.12);
\node[calllabel,anchor=north west,text=primaryblue,font=\footnotesize\bfseries\sffamily] at (.75,-4.34) {可修复};
\draw[callrequest] (2,-5.08)--node[calllabel,above]{带缺口与证据调整}(5.8,-5.08);
\draw[callreturn] (5.8,-5.71)--node[calllabel,above]{修订候选 / 动作建议}(2,-5.71);
\draw[callrequest] (2,-5.71)--(.27,-5.71)--(.27,-1.62)--(2,-1.62);
\node[callnote,text width=6.0cm,anchor=west] at (8.45,-5.25)
  {运行程序按 2.4.3 受控执行，\\取得新观察后按需重新验证；\\返工受原目标、授权与预算限制。};
\draw[accentorange!65,rounded corners=3pt] (.15,-6.37) rectangle (15.25,-8.09);
\node[calllabel,anchor=north west,text=accentorange!90!black,font=\footnotesize\bfseries\sffamily] at (.25,-6.42) {缺少关键输入 / 需要新授权};
\draw[callevent] (2,-7.11)--node[calllabel,above]{暂停依赖步骤，提出具体问题或授权请求}(13.4,-7.11);
\draw[callreturn] (13.4,-7.79)--node[calllabel,above]{补充信息或作出决定；复核继续条件后才恢复}(2,-7.79);
\draw[softgreen!45,rounded corners=3pt] (.15,-8.33) rectangle (15.25,-9.96);
\node[calllabel,anchor=north west,text=softgreen!60!black,font=\footnotesize\bfseries\sffamily] at (.25,-8.38) {收尾与结果区分};
\node[callnote,text width=13.8cm,anchor=west] at (.75,-9.33)
  {验收通过：交付成果与证据；有证据的否定发现：是否通过取决于约定。\\证据不足 / 确认无法完成 / 达到限制：说明实际进展、缺口及停止原因。\\取消：停止新派发，并请求终止在途作业；核实终止条件后才能确认取消完成。};
\end{tikzpicture}
\caption{单个 Agent 的检查与结束泳道示例：证据核验、受限返工与按需等待。}
\label{fig:single-agent-verification}
\end{figure}

若问题能在现有范围和限制内修复，程序把缺口与证据送入下一轮。瞬时网络错误可能适合重试，参数错误通常需要修正输入，缺少关键资料则需要补充信息。具有副作用且结果不明的操作，应先核实是否已生效；重试装饰器或循环结构本身不会判断是否可安全重复。调整应有新依据，并受尝试次数、时间、费用和资源约束，避免无信息增益的反复执行。

结束可来自验收通过、确认无法完成、达到限制或取消处理完成。收到取消请求后，停止派发相关新动作，并按能力请求终止在途工具或作业；核实适用的终止条件满足后，才能报告取消完成。仍在运行或结果不明时，应保留待确认状态并说明情况，不能把取消请求已记录等同于外部执行已终止。

对于实验数据检查，既可以交付得到支持的异常解释，也可以交付已确认的现象、排除的方法和仍缺少的证据；后者是否满足任务，应以原先约定的验收要求判断。有证据支持的否定发现与证据不足也应分别表述，不能将未能证明解释成立等同于已证明它不成立。最终说明应区分已完成工作、实际产物、证据和限制，并保留后续继续所需的有效引用。

澄清或授权未定时，相关步骤等待并说明待决事项与继续条件；已获授权且不依赖该决定的工作可继续。收尾须核对在途工具和外部作业、保存状态，通过“口耳”交付可访问的结果。“脑、手、口耳与记忆”由此闭环，但循环结束与结论正确仍需分别确认。



\section{多 Agent 系统的组织与协作机制}

第二章说明单个 Agent 如何把理解、判断、操作、通信和记忆组织成执行循环。本章进一步讨论：多个这样的循环如何分工、交换信息并利用共享服务。多 Agent 系统的增加值来自有意义的专业分工、并行处理与独立检查；实例数量增加，也会带来上下文传递、重复工作和协调成本。

以下图示表达逻辑职责与交互关系，不规定组织岗位、软件组合或部署拓扑。调度、权限、持久化、恢复和取消等能力，有些由框架或平台提供，有些由应用实现，须按实际产品、配置和运行证据判断。研发任务的交付可以是研究结论、算法与软件；这些机制本身不授予后续生产运行或交易操作的权限。

\subsection{从单个 Agent 到协作系统}
\label{sec:architecture}

\subsubsection{Hadoop 的职责分工与有限类比}

Agent 框架首先仍是软件框架：它把任务输入、模型调用、工具操作、状态与反馈组织成可执行过程，并按实现提供编排、调度或委派等能力。引入大语言模型，是在软件流程中增加理解与判断的方式，并没有离开传统软件系统的范畴；真正的执行仍依靠程序、接口和运行环境。框架也不只是调度器，它还可能提供上下文、工具与会话等抽象。为理解其中“组织工作”与“承载执行”的区别，可以先看 Hadoop 的职责划分，再作有限类比。

\textbf{先独立理解 Hadoop。}Hadoop 提供分布式存储与计算基础设施。以数据清洗为例，开发者先把处理规则写成程序，提交代码、输入和运行配置，再由基础设施安排执行；业务目标和处理逻辑由人及应用程序提供。

YARN 管理计算资源与应用运行。ResourceManager（RM）在集群范围管理和分配资源；ApplicationMaster（AM）面向一个应用申请资源、组织和监控任务；NodeManager（NM）在各节点启动和管理容器、监测资源使用。容器表示分配给执行任务的资源及其运行承载，与管理它的 NM 不是同一个对象。应用可以包含单个作业，也可以包含多个有关联的作业\cite{hadoop-yarn}。

HDFS 负责跨机器持久保存和共享文件，侧重大数据集的高吞吐访问与容错；它与 YARN 的计算管理分属不同职责。两者都可用于自有机器组成的环境，不以公有云为前提\cite{hadoop-hdfs}。

\textbf{再比较 Agent 系统中的部分职责。}Agent 系统也需要安排执行、监测过程和保存数据，但其组成不能与 Hadoop 逐项等同。下表仅比较工作内容。

\begin{center}
\small\renewcommand{\arraystretch}{1.16}
\begin{tabularx}{\textwidth}{P{4.4cm}L}
\toprule
\textbf{Hadoop 中的职责} & \textbf{Agent 系统中可类比的部分与限度}\\
\midrule
RM：跨应用分配资源 & 运行程序的资源调度与并发控制有相似职责；协调 Agent 的语义判断不等于 RM。\\
\addlinespace[3pt]
AM：组织一个应用的执行 & 任务编排也要推进步骤、跟踪状态；相似之处是执行生命周期管理。\\
\addlinespace[3pt]
NM 与容器：管理与承载执行 & 运行器负责启动和监测，进程或沙箱承载操作；不能把整个执行 Agent 对应为 NM 或容器。\\
\addlinespace[3pt]
HDFS：持久保存共享文件 & 数据与产物服务也提供可共享内容；这里只类比作用，不要求相同存储机制。\\
\bottomrule
\end{tabularx}
\end{center}

模型承担的需求理解、任务分析、动作选择和结果解释，更接近传统流程中由人或应用逻辑承担的语义判断。模型提出行动后，仍由普通程序落实权限检查、工具执行和状态记录。因此，这一类比\textbf{不表示组件一一对应、协议兼容或可以直接替换，也不构成对 Hadoop 的依赖}。多个 Agent 可以在一个本地程序中运行；多轮推理、多步骤或多个计算作业，都不等于需要多套集群。

\subsubsection{多 Agent 系统的组成与交互全景}

多 Agent 系统包含多个处理任务的运行实例，以及支持它们推理、操作、通信和保存状态的程序。每个实例都可以有自己的任务、上下文与执行循环。一个协调实例可组织其他实例；也可由确定性工作流派发，或由实例按条件移交，并非所有系统都必须设中央协调 Agent。

图~\ref{fig:multiagent-concept} 采用一种典型的协调式安排：人向居中的 A 提交目标，A 分别与 B、C 交换任务、消息和结果。三个脑图标分别表示三个实例在各自上下文中进行理解与判断。\textbf{多个脑表示多个 Agent 推理与行动循环，不表示每个实例独占一个模型、一台服务器或一份模型权重。}A 与 B 共享推理服务 X，C 示例使用另一服务 Y；三者也可以调用同一服务，或在不同步骤选用不同模型。

\begin{figure}[H]
\centering
\begin{tikzpicture}[x=1cm,y=1cm,font=\small\sffamily,text=darkgray,
  flow/.style={-{Stealth[length=2mm]},draw=primaryblue,line width=.9pt},
  exchange/.style={{Stealth[length=1.7mm]}-{Stealth[length=1.7mm]},draw=primaryblue,line width=.9pt},
  modelcall/.style={{Stealth[length=1.7mm]}-{Stealth[length=1.7mm]},draw=softpurple,line width=.9pt},
  capability/.style={draw=accentorange!85!black,dashed,line width=.85pt},
  title/.style={font=\small\bfseries\sffamily},
  detail/.style={font=\footnotesize\sffamily,align=center},
  panel/.style={rounded corners=4pt,draw=primaryblue!40,fill=white,align=center,inner sep=5pt}]
\path[use as bounding box] (0,-10.75) rectangle (15.4,.66);
\path[draw=primaryblue!35,fill=primaryblue!2,rounded corners=8pt] (.05,-1.1) rectangle (15.35,-9.78);
\node[title,fill=white,inner xsep=5pt,text=primaryblue] at (4.38,-1.1) {典型示例：A 协调 B、C};
% 人在执行系统之外；入口连线只是一种示例。
\begin{scope}[shift={(6.05,.1)},draw=darkgray,line width=1pt]
\draw[fill=white](0,.22) circle(.17);
\draw[fill=white](-.34,-.35)..controls(-.34,.06)and(.34,.06)..(.34,-.35)--cycle;
\end{scope}
\node[title,anchor=west] at (6.60,.27) {人类};
\node[detail,anchor=west] at (6.60,-.18) {目标、约束与必要决策；接收结果};
\draw[exchange] (7.70,-.54)--(7.70,-1.79);
\node[detail,fill=white,anchor=west] at (7.93,-.78) {目标 / 结果};
% 三个运行实例均有显著脑图标，不能误读为共用一个中央脑。
\foreach \x/\col/\letter in {2.4/softblue/B,7.7/primaryblue/A,13.0/softgreen/C}{
  \path[draw=\col!65,fill=white,rounded corners=5pt,line width=1pt] (\x-2.12,-1.8) rectangle (\x+2.12,-5.4);
  \node[title,text=\col!85!black] at (\x,-2.19) {Agent \letter};
  \draw[draw=\col!30,fill=\col!5,line width=.85pt](\x,-3.13) circle(.73);
  \begin{scope}[shift={(\x,-3.13)},scale=.65,draw=\col!90!black,line width=1.15pt,line cap=round,line join=round]
    \path[draw,fill=white]
      (0,.60)..controls(-.23,.88)and(-.65,.72)..(-.66,.44)
      ..controls(-1,.44)and(-1.10,.02)..(-.84,-.16)
      ..controls(-1.03,-.50)and(-.71,-.78)..(-.45,-.62)
      ..controls(-.24,-.86)and(0,-.65)..(0,-.43)
      ..controls(0,-.65)and(.24,-.86)..(.45,-.62)
      ..controls(.71,-.78)and(1.03,-.50)..(.84,-.16)
      ..controls(1.10,.02)and(1,.44)..(.66,.44)
      ..controls(.65,.72)and(.23,.88)..(0,.60)--cycle;
    \draw(0,.60)--(0,-.43);
    \draw(-.65,.40)..controls(-.30,.44)and(-.19,.22)..(-.35,.04);
    \draw(-.82,-.16)..controls(-.52,.04)and(-.26,-.10)..(-.42,-.37);
    \draw(.65,.40)..controls(.30,.44)and(.19,.22)..(.35,.04);
    \draw(.82,-.16)..controls(.52,.04)and(.26,-.10)..(.42,-.37);
  \end{scope}
  \node[detail,text=\col!80!black] at (\x,-4.04) {上下文 \letter\quad 独立执行循环};
  % 小手、对话与存储图标只是功能提示，不表示独立服务。
  \begin{scope}[shift={(\x-1.27,-4.62)},draw=accentorange!90!black,line width=.85pt,line cap=round,line join=round]
    \path[draw,fill=amber!7](-.15,-.14)--(-.22,.02)..controls(-.25,.09)and(-.16,.13)..(-.10,.04)
      --(-.08,.09)--(-.08,.31)..controls(-.08,.39)and(.01,.39)..(.01,.31)
      --(.01,.12)--(.01,.39)..controls(.01,.47)and(.10,.47)..(.10,.39)
      --(.10,.12)--(.10,.33)..controls(.10,.41)and(.19,.41)..(.19,.33)
      --(.19,.10)--(.19,.22)..controls(.19,.30)and(.28,.30)..(.28,.22)
      --(.28,-.06)--(.18,-.22)--(-.08,-.22)--cycle;
  \end{scope}
  \begin{scope}[shift={(\x,-4.59)},draw=softblue!90!black,line width=.85pt,line join=round]
    \path[draw,fill=softblue!5,rounded corners=1pt](-.28,-.07)--(-.28,.26)--(.25,.26)--(.25,-.07)--(-.06,-.07)--(-.18,-.21)--(-.16,-.07)--cycle;
    \draw(-.17,.15)--(.14,.15);\draw(-.17,.05)--(.05,.05);
  \end{scope}
  \begin{scope}[shift={(\x+1.28,-4.55)},draw=softpurple!80!black,line width=.85pt]
    \path[draw,fill=softpurple!5](-.23,.20)--(-.23,-.16)arc(180:360:.23 and .08)--(.23,.20);
    \draw[fill=white](0,.20)ellipse(.23 and .08);
    \draw(-.23,.04)arc(180:360:.23 and .08);
  \end{scope}
  \node[detail] at (\x-1.27,-5.06) {手};
  \node[detail] at (\x,-5.06) {口耳};
  \node[detail] at (\x+1.28,-5.06) {记忆};
}
\draw[exchange](4.54,-3.13)--(5.55,-3.13);
\node[detail,fill=primaryblue!2,inner sep=1pt] at (5.05,-2.66) {任务\\消息 / 结果};
\draw[exchange](9.85,-3.13)--(10.85,-3.13);
\node[detail,fill=primaryblue!2,inner sep=1pt] at (10.35,-2.66) {任务\\消息 / 结果};
% A/B 使用同一推理服务，C 可使用其他服务；模型端不是单一中央脑。
\node[panel,draw=softpurple!55,fill=softpurple!4,text width=6.35cm,minimum height=1.0cm] (mx) at (3.95,-6.66)
 {\textbf{共享模型服务 X}\\{\footnotesize A、B 分别提交自己的上下文并取得响应}};
\node[panel,draw=softpurple!55,fill=softpurple!4,text width=6.35cm,minimum height=1.0cm] (my) at (11.43,-6.66)
 {\textbf{另一模型服务 Y（也可与 X 相同）}\\{\footnotesize C 独立调用；分别传入上下文并接收响应}};
\draw[modelcall](2.4,-5.41)--(2.4,-6.06);
\draw[modelcall](7.7,-5.41)--(7.7,-5.95)--(5.65,-5.95)--(5.65,-6.06);
\draw[modelcall](13,-5.41)--(13,-6.06);
\node[detail,fill=primaryblue!2,text=softpurple] at (10.14,-5.93) {模型请求 / 响应};
\node[detail,text=darkgray!90] at (7.7,-7.48) {模型服务位置仅表示调用关系，不表示本地、云端或独占部署};
% 共享能力与运行功能：分别列示，不要求经单一中央服务。
\node[panel,draw=accentorange!50,fill=amber!5,text width=6.35cm,minimum height=.92cm] (tools) at (3.95,-8.35)
 {\textbf{共享工具、接口与计算能力}\\{\footnotesize 可按任务分别授权；短调用或外部长作业}};
\node[panel,draw=softblue!45,fill=softblue!4,text width=6.35cm,minimum height=.92cm] (data) at (11.43,-8.35)
 {\textbf{数据与状态保存（按实际需要）}\\{\footnotesize 本地或远程；显式读取，不自动同步上下文}};


\node[detail,text width=14.4cm] at (7.7,-9.27)
 {运行功能：调度、上下文组装、模型与工具接入、状态管理。\\这些能力可由各实例、框架、应用或运行环境提供。};
\node[detail,text width=14.6cm,anchor=north] at (7.7,-10.03)
 {本图只画一种典型拓扑。获授权且接口支持时，B 与 C 也可直接交互；并非必须由 A 转发。};
\end{tikzpicture}
\caption{多 Agent 系统的组成与交互全景：以 A 为协调者的一种典型逻辑安排。}
\label{fig:multiagent-concept}
\end{figure}

图中未画出的 B--C 连线不表示禁止直接交互；是否允许直接协作，取决于授权、接口和应用规则。协调也可由确定性工作流承担，不要求任何系统都采用图中的中心安排。

人位于自动执行系统之外，提交意图、目标和约束，接收结果，并在需要澄清、重要取舍或新授权时参与。系统内的任务、消息和结果由程序传递，不要求人逐步转交。图中的运行功能可分散在各实例、框架库、应用程序或托管环境中；工具、模型和存储也可以分别提供，图中的连线不规定部署数量与物理位置。

\subsubsection{角色、运行实例与执行节点的映射}

角色定义回答长期负责什么及可使用什么能力；运行实例回答某次任务由谁处理、已进行到哪里以及掌握哪些信息；执行节点则提供运行程序或作业所需的环境。第二章使用的对象名称是解释方式，不要求所有框架采用相同的类、字段或生命周期划分。

\begin{table}[H]
\centering\small\renewcommand{\arraystretch}{1.18}
\begin{tabularx}{\textwidth}{P{2.6cm}P{4.5cm}L}
\toprule
\textbf{对象} & \textbf{示例} & \textbf{映射关系与边界}\\
\midrule
角色定义 & 一份资料分析角色配置 & 可产生多个实例，不等于常驻进程。\\
\rowcolor{lightgray}
运行实例 & 分别研究资料甲、乙的两次运行 & 可复用角色与模型，各有上下文和执行状态。\\
执行节点 & 本地进程、服务器或托管环境 & 一个环境可承载多个实例；进程、容器与物理机是不同层级。\\
\rowcolor{lightgray}
模型服务 & 两个实例调用同一推理 API & 只因此共享推理能力，不自动共享任务、记忆和操作权限。\\
\bottomrule
\end{tabularx}
\caption{角色、实例、执行节点与模型服务不存在固定的一一对应关系}
\end{table}

一个角色可以同时运行多个实例，一个执行环境也可以承载不同角色。工作流图中的“节点”通常表示处理步骤，它可以是普通函数、模型请求或 Agent 循环，不必是一台计算机。逻辑上的多 Agent 因而不天然等于物理上的分布式系统；是否分布到多台机器，应由资源、隔离和运行要求决定。

\subsubsection{框架、应用与运行环境的职责}

框架提供可复用的执行与组织方式，例如模型和工具抽象、循环、状态图、会话或委派接口；应用负责把这些能力接到具体目标、资料、工具和业务规则上；运行环境承载进程、账户、网络、存储和算力。三者可以以库、服务或一体化平台交付，功能有时重叠，但责任仍需在实际实现中明确。

例如，OpenAI Agents SDK 提供 Agent、Runner、工具、handoff 和会话等抽象；LangGraph 用图状态、处理节点与边组织执行，并可配置持久化机制\cite{agents-sdk,langgraph-overview,langgraph-persist}。这些是两种具体实现，不能据此推导出所有框架都有相同调度器、消息协议、状态库或检查点。

框架有一个调用入口，不等于应用已配置费用上限、业务审批或可靠恢复。可执行范围还取决于工具实现、运行账户、数据库权限、沙箱和网络策略。角色指令可以说明边界，但只有相应程序与权限机制才能在执行侧落实约束。

\subsection{任务分工与协作控制}
\label{sec:orchestration}

\subsubsection{任务拆分、派发与执行状态}

任务可来自用户请求、预定流程或外部事件。拆分时先识别目标、输入、依赖和可检查的结果，再决定哪些工作可独立处理、哪些需要等待前序结果。对资料作并行研究、对算法作实现和独立测试，属于不同拆分维度；拆得过细会使交接成本超过收益，重复分配同一问题也不一定提高可靠性。

派发把工作交给具体实例或处理步骤。输入可以是自然语言、函数参数、消息或结构化对象，通常需带上目标、必要背景、范围和期望结果，不要求所有系统使用统一任务表。路由可由规则或模型提出，但创建运行、传入输入和启动执行的是程序。关键目标或授权缺失时，可暂停依赖它的步骤；不受影响的工作是否继续，取决于流程和实际运行能力。

执行状态用于区分未开始、处理中、等待、完成、失败或取消等情况，状态名称与转换由实现定义。\textbf{任务已派发、实例已结束与目标已完成，是不同层次的事实。}例如，实例提交了长作业后可能进入等待，而不是完成；某个分支返回成功，也不表示其他分支或最终验收已结束。并发数、排队、预算与超时控制须由框架、平台或应用实际实现，不能从流程图中的并行箭头直接推定。


图~\ref{fig:task-decomposition} 把“确定要做哪些工作”与“交给谁执行”分开表示。甲、乙没有相互依赖时可成为并行候选；丙需要前两项结果，因此仍须等待。任务描述、处理任务的实例，以及工具提交的计算作业，是三个不同对象。

\begin{figure}[htbp]
\centering
\begin{tikzpicture}[x=1cm,y=1cm,font=\footnotesize\sffamily,text=darkgray,
  stage/.style={rounded corners=4pt,draw=primaryblue!50,fill=primaryblue!3,minimum width=2.55cm,minimum height=3.65cm},
  title/.style={font=\small\bfseries\sffamily,text=primaryblue,align=center},
  tinybox/.style={rounded corners=2pt,draw=primaryblue!70,fill=white,minimum width=.66cm,minimum height=.40cm,inner sep=2pt},
  txt/.style={align=center,text width=2.3cm},
  flow/.style={-{Stealth[length=1.8mm]},draw=primaryblue,line width=.9pt}]
\path[use as bounding box] (0,-5.25) rectangle (15.4,.20);
\foreach \x/\title in {1.38/目标,4.54/子任务与依赖,7.70/分配给实例,10.86/执行状态,14.02/汇总与交付}{
  \node[stage] at (\x,-1.65) {};
  \node[title] at (\x,-.18) {\title};
}
\foreach \x in {2.72,5.88,9.04,12.20}{\draw[flow](\x,-1.55)--++(.43,0);}
% 目标靶心。
\begin{scope}[shift={(1.38,-1.15)},draw=primaryblue,line width=1pt]
\draw[fill=white](0,0)circle(.47);\draw(0,0)circle(.29);\draw[fill=primaryblue](0,0)circle(.085);
\draw[-{Stealth[length=1.5mm]},draw=accentorange,line width=1.4pt](.62,.53)--(.05,.03);
\end{scope}
\node[txt] at (1.38,-2.38) {明确问题、范围\\与验收要求};
\node[txt,text=primaryblue] at (1.38,-3.16) {拆分：做什么};
% 子任务依赖图：甲、乙可独立，丙依赖二者。
\node[tinybox] (ta) at (4.00,-.95) {甲};
\node[tinybox] (tb) at (5.08,-.95) {乙};
\node[tinybox] (tc) at (4.54,-1.95) {丙};
\draw[flow](ta.south)--(tc.north west);
\draw[flow](tb.south)--(tc.north east);
\node[txt] at (4.54,-2.79) {甲、乙可独立\\丙等待甲、乙};
% 两个脑徽记分别表示独立实例与上下文。
\foreach \x/\letter in {7.16/A,8.24/B}{
  \begin{scope}[shift={(\x,-1.12)},scale=.33,draw=primaryblue,line width=1.4pt,line cap=round]
    \path[draw,fill=white](0,.60)..controls(-.23,.88)and(-.65,.72)..(-.66,.44)
    ..controls(-1,.44)and(-1.10,.02)..(-.84,-.16)
    ..controls(-1.03,-.50)and(-.71,-.78)..(-.45,-.62)
    ..controls(-.24,-.86)and(0,-.65)..(0,-.43)
    ..controls(0,-.65)and(.24,-.86)..(.45,-.62)
    ..controls(.71,-.78)and(1.03,-.50)..(.84,-.16)
    ..controls(1.10,.02)and(1,.44)..(.66,.44)
    ..controls(.65,.72)and(.23,.88)..(0,.60)--cycle;
    \draw(0,.60)--(0,-.43);
    \draw(-.64,.35)..controls(-.2,.40)and(-.2,.05)..(-.43,-.12);
    \draw(.64,.35)..controls(.2,.40)and(.2,.05)..(.43,-.12);
  \end{scope}
  \node at (\x,-1.72) {实例 \letter};
}
\node[txt] at (7.70,-2.45) {各有上下文\\可接不同子任务};
\node[txt,text=primaryblue] at (7.70,-3.16) {派发：谁来做};
% 状态提示不是统一字段规范。
\foreach \y/\state/\col in {-.88/待执行/darkgray,-1.45/处理中/primaryblue,-2.02/等待/amber}{
  \draw[fill=\col!85,draw=\col] (10.03,\y) circle(.065);
  \node[anchor=west] at (10.18,\y) {\state};
}
\node[txt] at (10.86,-2.77) {完成、失败\\或取消};
% 产物与检查标记。
\begin{scope}[shift={(14.02,-1.18)},draw=softgreen!80!black,line width=1pt]
\draw[fill=white](-.35,-.42)rectangle(.33,.46);
\draw(-.21,.23)--(.17,.23);\draw(-.21,.07)--(.17,.07);
\draw[line width=1.5pt](-.15,-.18)--(-.01,-.31)--(.27,-.02);
\end{scope}
\node[txt] at (14.02,-2.41) {核对必要结果\\并对照原始目标};
\node[txt,text=softgreen!65!black] at (14.02,-3.16) {交付或报缺项};
\node[callnote,text width=14.8cm,anchor=north west] at (.25,-3.92)
 {独立任务只是并行候选，实际并发还取决于调度、资源与权限。一个实例可处理多个任务，一个任务也可跨实例推进；实例结束或计算作业成功，都不能代替对任务目标的验收。};
\end{tikzpicture}
\caption{从目标到协作执行：先拆分任务与依赖，再派发实例并跟踪状态。}
\label{fig:task-decomposition}
\end{figure}

\subsubsection{串行、并行与结果汇总}

串行执行适合存在明确依赖的步骤，后一步根据前一步的有效结果继续；并行执行适合相对独立、且资源和权限允许同时开展的任务。汇总并非简单拼接文本，而要对照原目标确认覆盖范围、共同假设、矛盾与尚缺证据。格式转换、数据合并或确定性检查通常可由普通程序完成，不必每一步都调用模型。

\begin{figure}[htbp]
\centering
\begin{tikzpicture}[x=1cm,y=1cm,font=\footnotesize\sffamily,text=darkgray,
  panel/.style={rounded corners=5pt,draw=primaryblue!35,fill=primaryblue!2},
  work/.style={callstep,text width=1.25cm,minimum height=.68cm},
  serialhead/.style={callhead,minimum width=1.6cm,font=\footnotesize\bfseries\sffamily}]
\path[use as bounding box] (0,-7.60) rectangle (15.4,.45);
\foreach \x in {0,7.9}{\path[panel](\x, .38) rectangle (\x+7.5,-6.56);}
\node[font=\small\bfseries\sffamily,text=primaryblue] at (3.75,.04) {串行：乙依赖甲的有效结果};
\node[font=\small\bfseries\sffamily,text=primaryblue] at (11.65,.04) {并行候选：甲、乙彼此独立};
\foreach \offset in {0,7.9}{
  \foreach \x/\name in {1.10/协调者,3.70/实例 A,6.35/实例 B}{
    \node[serialhead] at (\offset+\x,-.66) {\name};
    \draw[calllife](\offset+\x,-1.02)--(\offset+\x,-5.38);
  }
}
% 左：后续派发必须等待甲的有效结果。
\draw[callrequest](1.1,-1.33)--node[calllabel,above]{派发甲}(3.7,-1.33);
\node[work,fill=softblue!7] at (3.7,-2.00) {处理甲};
\draw[callreturn](3.7,-2.66)--node[calllabel,above]{甲的有效结果}(1.1,-2.66);
\draw[callrequest](1.1,-3.26)--node[calllabel,above,pos=.63]{据甲的结果派发乙}(6.35,-3.26);
\node[work,fill=softgreen!7] at (6.35,-3.99) {处理乙};
\draw[callreturn](6.35,-4.72)--node[calllabel,above]{乙的结果}(1.1,-4.72);
\node[callstep,text width=5.6cm,fill=primaryblue!6] at (3.75,-5.63) {核对依赖结果，再汇总与交付};
\draw[callrequest](1.1,-4.84)--(1.1,-5.28);
% 右：派发时刻与完成时刻可以不同；两条支线不保证同时运行。
\draw[callrequest](9.00,-1.33)--node[calllabel,above]{派发甲}(11.60,-1.33);
\draw[callrequest](9.00,-1.98)--node[calllabel,above,pos=.69]{派发乙}(14.25,-1.98);
\node[work,fill=softblue!7] at (11.6,-2.93) {处理甲};
\node[work,fill=softgreen!7] at (14.25,-2.93) {处理乙};
\draw[callreturn](11.60,-3.95)--node[calllabel,above]{甲的结果}(9.00,-3.95);
\draw[callreturn](14.25,-4.72)--node[calllabel,above,pos=.30]{乙的结果}(9.00,-4.72);
\draw[draw=accentorange,line width=2.3pt](9,-3.95)--(9,-5.23);
\node[calllabel,anchor=west,text=accentorange!85!black] at (9.15,-4.38) {等待必要结果};
\node[callstep,text width=5.6cm,fill=primaryblue!6] at (11.65,-5.63) {汇合：必要分支齐备后检查与汇总};
\draw[callrequest](9,-4.87)--(9,-5.28);
\node[calllabel,text width=6.5cm] at (3.75,-6.18) {时间向下；甲完成是乙开始的前提};
\node[calllabel,text width=6.5cm] at (11.65,-6.18) {图示为允许重叠运行时的示例};
\node[callnote,text width=14.8cm,anchor=north west] at (.25,-6.88)
 {箭头表示派发与返回顺序，不代表实际线程或进程已经并发。汇合条件由任务规定；若允许部分结果交付，应明确缺项、限制以及其余分支的处理方式。};
\end{tikzpicture}
\caption{串行与并行协作的时序比较：显式依赖决定等待，汇合条件决定何时继续。}
\label{fig:multiagent-flow}
\end{figure}

汇合点可以等待全部必要分支，也可以在满足预定条件时继续，但必须区分必要与可选结果。若允许超时后交付部分结果，应明确哪些内容已完成、哪些结论仍受缺失输入影响；其余分支是否继续或请求取消也需有明确处理。不能因为某个分支先返回，就默认为整体任务已经满足要求。

\subsubsection{协调、移交与 Agent 工具化}

协调模式中，一个 Agent 或确定性程序保持总体任务，派发子任务、处理冲突并汇总结果。协调者可以参与推理，也可以只是按规则运行；并非每种系统都需要单独的协调角色。

移交把后续处理责任交给另一个实例。交接内容通常包括当前任务、必要背景、已完成部分与待解决问题；接收方按照自己的上下文和规则继续。移交时是否携带历史、怎样筛选历史，以及原实例是否继续参与，都依框架接口和应用安排而定。

Agent 工具化则把一个 Agent 封装成可调用能力：调用方提交子问题，被调用实例在自己的循环内处理并返回结果，调用方通常继续掌握整体流程。OpenAI Agents SDK 的 handoff 与 Agent-as-tool 体现了这种控制流差别\cite{agents-sdk}。两种模式都可能用函数或网络实现，不能只凭使用了哪种传输方式判断协作责任是否已经转移。


图~\ref{fig:coordination-control} 关注后续处理责任由谁掌握，而不是通信采用函数还是网络。协调者保留总体目标；移交后由接收者继续后续处理；工具化调用返回后，由原调用方继续其流程。图中的移交采用原实例退出当前处理的典型情形，其他安排应以实际接口约定为准。

\begin{figure}[htbp]
\centering
\begin{tikzpicture}[x=1cm,y=1cm,font=\footnotesize\sffamily,text=darkgray,
  panel/.style={rounded corners=5pt,draw=primaryblue!40,fill=primaryblue!2},
  head/.style={callhead,minimum width=1.7cm,font=\footnotesize\bfseries\sffamily},
  controlstep/.style={callstep,text width=1.68cm,minimum height=.80cm}]
\path[use as bounding box](0,-6.05)rectangle(15.4,.43);
\foreach \offset/\title in {0/协调,5.2/移交（handoff）,10.4/Agent 工具化}{
  \path[panel](\offset,.36)rectangle(\offset+5,-5.28);
  \node[font=\small\bfseries\sffamily,text=primaryblue] at (\offset+2.5,0) {\title};
  \foreach \x in {1.03,3.97}{\draw[calllife](\offset+\x,-1.00)--(\offset+\x,-4.67);}
}
\node[head] at (1.03,-.67) {协调者};
\node[head] at (3.97,-.67) {执行实例};
\draw[callrequest](1.03,-1.43)--node[calllabel,above]{派发子任务}(3.97,-1.43);
\node[controlstep,fill=softblue!7] at (3.97,-2.16) {处理各项\\子任务};
\draw[callreturn](3.97,-3.02)--node[calllabel,above]{结果与问题}(1.03,-3.02);
\node[controlstep,fill=primaryblue!7] at (1.03,-3.98) {汇总协调\\继续执行};
\draw[callrequest](1.03,-3.18)--(1.03,-3.58);
\node[calllabel,text width=4.4cm] at (2.5,-4.82) {总体目标仍由协调者维护};
\node[head] at (6.23,-.67) {原实例 A};
\node[head] at (9.17,-.67) {接收实例 B};
\draw[callrequest](6.23,-1.63)--node[calllabel,above,align=center]{后续责任与\\必要上下文}(9.17,-1.63);
\draw[draw=darkgray,line width=1pt](6.05,-2.14)--(6.41,-2.50);\draw[draw=darkgray,line width=1pt](6.05,-2.50)--(6.41,-2.14);
\node[calllabel,text width=1.55cm] at (6.23,-2.98) {结束当前\\处理};
\node[controlstep,fill=softgreen!7] at (9.17,-2.63) {接手任务\\自行继续};
\draw[callrequest](9.17,-3.09)--(9.17,-3.56);
\node[controlstep,fill=softgreen!7] at (9.17,-3.98) {后续处理\\完成交付};
\node[calllabel,text width=4.4cm] at (7.7,-4.82) {后续处理转到接收者};
\node[head] at (11.43,-.67) {调用方 A};
\node[head] at (14.37,-.67) {工具 Agent B};
\draw[callrequest](11.43,-1.43)--node[calllabel,above]{请求子问题}(14.37,-1.43);
\node[controlstep,fill=softblue!7] at (14.37,-2.16) {独立循环\\处理请求};
\draw[callreturn](14.37,-3.02)--node[calllabel,above]{返回结果}(11.43,-3.02);
\node[controlstep,fill=primaryblue!7] at (11.43,-3.98) {使用结果\\继续执行};
\draw[callrequest](11.43,-3.18)--(11.43,-3.58);
\node[calllabel,text width=4.4cm] at (12.9,-4.82) {调用返回后，原调用方继续};
\node[callnote,text width=14.8cm,anchor=north west] at (.25,-5.56)
 {实线为请求或控制推进，虚线为结果返回，时间向下。协调可包含工具化调用；三种模式可以组合，但不能混淆它们的责任与返回约定。};
\end{tikzpicture}
\caption{协调、移交与 Agent 工具化：用后续处理责任和返回路径区分协作方式。}
\label{fig:coordination-control}
\end{figure}

\subsubsection{协作结果检查、返工与终止}
\label{sec:normalflow}

结果检查可由确定性程序、另一个 Agent 或人承担。程序可核验文件、格式和测试结果；模型可解释证据、识别遗漏，但也可能判断错误。若要求独立验证，应向验证者提供固定的规格、候选成果与必要数据，使其自行检查并记录依据，不能仅把执行者的总结换一种说法。独立性来自职责、证据与检查过程，而不是角色名称或多个实例的一致表态。

返工应针对具体缺项、失败证据或方法问题。短暂网络故障可能适合重试，输入不足需要补充资料，方法失败则可能需要调整方案。重试前还应确认操作是否可安全重复；超时但外部动作是否完成未知时，先查询真实状态，避免重复提交或写入。重试策略本身不能替应用判断副作用。

结束条件可以是通过检查、形成有依据的否定结论、确认无法继续，或达到约定时限、预算与轮次上限。收到取消后，应区分“取消请求已发出”和“相关工作已停止”；已有结果、仍在运行的部分和未知状态需分别说明。部分结果可以有价值，但应保留限制；循环结束、作业成功和成果满足验收条件也不能混为一谈。

图~\ref{fig:multiagent-review} 展示跨实例的独立检查与协作收尾。它在执行 Agent 自身循环之外增加独立检查者，协调者依据检查证据决定交付、有界返工或停止；返工沿用原目标和验收要求，目标变更则需要另行澄清或授权。

\begin{figure}[htbp]
\centering
\begin{tikzpicture}[x=1cm,y=1cm,font=\footnotesize\sffamily,text=darkgray,
  rework/.style={-{Stealth[length=2mm]},draw=softred,line width=.95pt},
  rlabel/.style={calllabel,text=softred}]
\path[use as bounding box](0,-8.08)rectangle(15.4,.48);
\foreach \x in {2.4,7.7,13.0}{\draw[calllife](\x,-.49)--(\x,-6.78);}
\node[callhead,minimum width=3.65cm] at (2.4,0) {执行 Agent};
\node[callhead,minimum width=3.95cm] at (7.7,0) {独立检查者\\{\footnotesize 程序 / 另一个 Agent / 人}};
\node[callhead,minimum width=3.65cm] at (13.0,0) {协调者 / 流程程序};
\draw[callrequest](13,-1.00)--node[calllabel,above]{固定目标、验收要求与执行边界}(2.4,-1.00);
\draw[callrequest](2.4,-1.88)--node[calllabel,above]{候选成果、规格与必要数据}(7.7,-1.88);
\node[callstep,text width=3.6cm,fill=softgreen!6] at (7.7,-2.70) {自行核验并记录依据\\不只复述执行者总结};
\draw[callreturn](7.7,-3.52)--node[calllabel,above]{检查结论、证据与限制}(13,-3.52);
\node[callstep,text width=3.5cm,fill=primaryblue!6] at (13,-4.33) {核对验收与结束条件};
\draw[rework](13,-5.13)--node[rlabel,above]{未通过且在限额内：依据失败证据返工，目标不变}(2.4,-5.13);
\node[callstep,text width=3.45cm,draw=softred!65,fill=softred!3] at (2.4,-6.02) {补充处理\\再次送交独立检查};
\draw[rework](2.4,-5.21)--(2.4,-5.58);
\draw[rework](.60,-6.02)--(.22,-6.02)--(.22,-1.88)--(2.28,-1.88);
\node[calllabel,text width=3.75cm] at (7.7,-5.95) {返工受轮次、预算\\与时限约束};
\node[callstep,text width=3.60cm,fill=primaryblue!6] at (13,-6.12) {通过：交付有证据的结果\\停止：说明原因与实际状态};
\draw[callrequest](13,-4.78)--(13,-5.62);
\node[calllabel,anchor=west,text=primaryblue,align=center] at (13.15,-5.08) {通过或停止};
\node[callnote,text width=14.6cm,anchor=north west] at (.4,-7.03)
 {部分结果、超时或取消请求都不自动等于验收通过。达到限额应停止相应推进；取消后还需核实在途操作与外部作业，区分已停止、仍运行和状态未知。};
\end{tikzpicture}
\caption{跨实例的独立检查、有界返工与终止：协调决定必须以证据和真实状态为依据。}
\label{fig:multiagent-review}
\end{figure}



\subsection{实例间通信与数据交换}
\label{sec:governance}

\subsubsection{任务、消息、结果与产物引用}

任务说明做什么及如何判断结束；消息传递进展、问题、协调意见或交接；结果回答完成情况与结论；产物引用则指向代码、文档、数据或计算输出。它们可以合并在同一调用返回中，也可以分别交换，不要求一种统一的消息结构。对需要追溯的任务，引用还应足以确定所使用的版本或具体内容。

交接要让接收方能够理解并继续工作，不必复制发送方全部历史。发送者应带上必要上下文、适用限制与结果依据；接收者核对其完整性并按需读取来源。任务送达、处理完成和结果合格是分别需要判断的事实。下图保留委派与返回的基本时序，具体形式可以是调用、消息或事件。

\begin{figure}[H]
\centering
\begin{tikzpicture}[x=1cm,y=1cm,font=\footnotesize\sffamily,text=darkgray]
\path[use as bounding box] (0,-5.98) rectangle (15.4,.4);
\foreach \x in {2,7.7,13.4}{\draw[calllife](\x,-.38)--(\x,-5.1);}
\node[callhead,minimum width=3.4cm] at (2,0) {Agent A 的运行程序};
\node[callhead,minimum width=4.35cm] at (7.7,0) {调用或消息传递机制};
\node[callhead,minimum width=3.4cm] at (13.4,0) {Agent B 的运行程序};
\draw[callrequest](2,-.9)--node[calllabel,above]{提交任务与必要背景}(7.7,-.9);
\draw[callrequest](7.7,-1.6)--node[calllabel,above]{传入内容或可解析引用}(13.4,-1.6);
\node[callstep,text width=3.2cm] at (13.4,-2.35) {读取获准输入并处理任务};
\draw[callreturn](13.4,-3.15)--node[calllabel,above]{结果 / 问题 / 错误 / 引用}(7.7,-3.15);
\draw[callreturn](7.7,-3.88)--node[calllabel,above]{转换为调用方可用的返回}(2,-3.88);
\node[callstep,text width=3.2cm] at (2,-4.66) {根据结果继续或结束};
\node[callnote,anchor=north west,text width=14.7cm] at (.3,-5.35)
 {中间泳道表示传递功能，可是函数调用、框架内部消息或网络协议，不要求独立中转服务。\\任务已送达、处理已结束与结果满足要求，应分别判断。};
\end{tikzpicture}
\caption{一次委派与结果返回的逻辑顺序。实例之间不需要持续的模型聊天连接。}
\label{fig:agent-delegation}
\end{figure}

引用不等于内容，也不授予访问权限。接收方需要可解析该引用的工具、实际可达的位置和有效授权；一台机器上的绝对路径，通常不能直接供另一台机器使用。文件标识、URL 或数据库主键的含义由提供方定义。大文件可保留在原存储中，通过引用交接，再由接收方读取必要片段，避免无差别复制。

\subsubsection{进程内调用、远程接口与消息传递}

同一进程中可用函数参数和返回对象交互；跨进程可使用进程通信或队列；跨机器可通过网络接口、远程调用或消息服务连接。同步调用使调用方等待返回，异步传递允许发送和处理分离。这些属于传输与调用方式；任务由谁负责、何时完成、结果如何验收，则属于协作语义，两者需要分别设计。


图~\ref{fig:communication-boundaries} 按进程与机器边界比较几种传递方式。图中的连线只表示内容如何到达，不规定任务责任，也不表示每次调用都经过网络或独立中转服务。

\begin{figure}[H]
\centering
\begin{tikzpicture}[x=1cm,y=1cm,font=\footnotesize\sffamily,text=darkgray,
  endpoint/.style={rounded corners=3pt,draw=primaryblue!55,fill=primaryblue!7,
    text width=2.7cm,minimum height=.68cm,align=center,inner sep=4pt},
  transport/.style={-{Stealth[length=2mm]},draw=cyan!55!blue,line width=.9pt},
  reply/.style={transport,dashed},
  event/.style={<->,>=Stealth,draw=cyan!55!blue,densely dotted,line width=1pt},
  labeltext/.style={fill=white,inner sep=2pt,align=center},
  panel/.style={rounded corners=4pt,draw=cyan!45!blue!35,fill=cyan!3},
  paneltitle/.style={anchor=west,font=\footnotesize\bfseries\sffamily,text=cyan!40!blue}]
\path[use as bounding box] (0,-7.85) rectangle (15.4,.05);

\draw[panel] (0,0) rectangle (15.4,-1.65);
\node[paneltitle] at (.25,-.27) {同一进程：直接调用，不跨网络};
\node[endpoint] (in-a) at (2.3,-.97) {运行程序 A};
\node[endpoint] (in-b) at (13.1,-.97) {处理步骤 B};
\draw[transport] ([yshift=5pt]in-a.east)--node[labeltext,above]{函数调用：传入参数}([yshift=5pt]in-b.west);
\draw[reply] ([yshift=-5pt]in-b.west)--node[labeltext,below]{返回对象或错误}([yshift=-5pt]in-a.east);

\draw[panel] (0,-1.87) rectangle (15.4,-3.52);
\node[paneltitle] at (.25,-2.14) {同一机器、不同进程：通过进程间通道};
\node[endpoint] (ipc-a) at (2.3,-2.84) {进程 A};
\node[endpoint] (ipc-b) at (13.1,-2.84) {进程 B};
\draw[transport] ([yshift=5pt]ipc-a.east)--node[labeltext,above]{IPC / 队列：请求或消息}([yshift=5pt]ipc-b.west);
\draw[reply] ([yshift=-5pt]ipc-b.west)--node[labeltext,below]{返回通道：结果或响应消息}([yshift=-5pt]ipc-a.east);

\draw[panel] (0,-3.74) rectangle (15.4,-6.74);
\node[paneltitle] at (.25,-4.02) {跨机器：可以直接调用，也可以异步传递};
\node[endpoint,minimum height=.8cm] (net-a) at (2.3,-4.91) {节点 A\\运行程序};
\node[endpoint,minimum height=.8cm] (net-b) at (13.1,-4.91) {节点 B\\服务或运行程序};
\draw[transport] ([yshift=5pt]net-a.east)--node[labeltext,above]{RPC / API：请求}([yshift=5pt]net-b.west);
\draw[reply] ([yshift=-5pt]net-b.west)--node[labeltext,below]{响应或结果引用}([yshift=-5pt]net-a.east);
\node[rounded corners=3pt,draw=cyan!55!blue,fill=cyan!9,text width=3.5cm,
  minimum height=.62cm,align=center,inner sep=4pt] (broker) at (7.7,-6.03) {消息服务 / 队列（可选）};
\draw[event] (net-a.south)--(2.3,-6.03)--node[labeltext,above]{投递 / 接收}(broker.west);
\draw[event] (broker.east)--node[labeltext,above]{投递 / 接收}(13.1,-6.03)--(net-b.south);
\node[align=center] at (7.7,-6.53) {异步消息可携带任务、通知或结果；发送与处理可分离};

\node[anchor=north west,align=left,text width=15.0cm] at (.2,-6.99)
 {实线为请求，虚线为返回，点线为可选的异步消息通道。各行是可选实现示例，非必经步骤。\\模型 API 连接推理服务，工具接口 / MCP 连接工具与资源，A2A 连接 Agent 服务；不必同时采用。};
\end{tikzpicture}
\caption{逻辑调用与物理传输的对应示例。传输方式解决“如何传递”，协作语义还需定义任务责任、完成条件与验收。}
\label{fig:communication-boundaries}
\end{figure}

\begin{table}[H]
\centering\small\renewcommand{\arraystretch}{1.17}
\begin{tabularx}{\textwidth}{P{2.7cm}P{4.4cm}L}
\toprule
\textbf{接口类型} & \textbf{主要连接对象} & \textbf{不能据此推定的能力}\\
\midrule
模型 API & 运行程序与推理服务 & 不等于 Agent 间的任务交接协议。\\
\rowcolor{lightgray}
工具接口或 MCP & 运行程序与工具、资源服务 & 不自动定义多 Agent 分工、接收确认或验收。\\
框架内部调用或消息 & 应用内的实例与处理步骤 & 不保证不同框架理解相同状态与错误含义。\\
\rowcolor{lightgray}
A2A 等互操作协议 & 独立 Agent 服务之间 & 仍需实现兼容、身份认证与业务权限。\\
\bottomrule
\end{tabularx}
\caption{传输连接与协作语义需要区分，三类对外接口也各有用途}
\end{table}

MCP 与 A2A 分别提供工具连接和 Agent 互操作的相关规范，不要求每个多 Agent 应用同时采用\cite{mcp-transports,a2a-spec}。一个外部服务即使能收发 JSON，也不代表已经具备兼容的 Agent 协作接口；应核对操作、消息、状态、产物和错误的实际含义。

网络断连、重复投递和顺序变化还会影响协作。需要可靠交接时，可结合请求关联、去重、确认与重试处理；适用机制依传输和任务风险选择。收到消息不等于完成任务，队列确认也未必等于业务成功；这些保障不能仅凭“支持多 Agent”推定存在。来自其他实例或外部工具的内容仍是待判断的信息，不能自动提升为更高权限的指令。

\subsubsection{共享数据与独立上下文}

单个 Agent 如何访问外部数据、组织上下文和取用记忆，见第~\ref{sec:single-agent-data} 节。多实例场景增加的问题是：哪些内容允许共享，接收方何时读取，以及多个实例是否依据同一版本工作。业务数据仍可按实际需要保存在本地或远程，不因协作而必须采用统一数据类型、存储形式或专用节点；约束来自实际接口、权限与运行环境。




图~\ref{fig:shared-data-contexts} 展示一次有依赖的更新交接：写入、通知、读取和纳入上下文是不同动作；共享存储只提供可访问的外部内容，不直接改写另一个实例正在使用的上下文。

\begin{figure}[H]
\centering
\begin{tikzpicture}[x=1cm,y=1cm,font=\footnotesize\sffamily,text=darkgray,
  agentframe/.style={rounded corners=5pt,draw=primaryblue!45,fill=primaryblue!3},
  context/.style={rounded corners=3pt,draw=primaryblue!65,fill=primaryblue!8,
    text width=4.75cm,minimum height=1.1cm,align=center,inner sep=5pt},
  tool/.style={rounded corners=3pt,draw=accentorange!75!black,fill=accentorange!10,
    text width=4.5cm,minimum height=.66cm,align=center,inner sep=4pt},
  dataflow/.style={-{Stealth[length=2mm]},draw=softpurple,line width=.95pt},
  communication/.style={-{Stealth[length=2mm]},draw=cyan!55!blue,line width=.9pt},
  localflow/.style={-{Stealth[length=2mm]},draw=primaryblue,line width=.9pt},
  labeltext/.style={fill=white,inner sep=2pt,align=center}]
\path[use as bounding box] (0,-7.42) rectangle (15.4,.05);
\draw[agentframe] (0,0) rectangle (6.1,-3.83);
\draw[agentframe] (9.3,0) rectangle (15.4,-3.83);
\node[font=\small\bfseries\sffamily,text=primaryblue] at (3.05,-.38) {Agent A：独立实例};
\node[font=\small\bfseries\sffamily,text=primaryblue] at (12.35,-.38) {Agent B：独立实例};
\node[context] (ctx-a) at (3.05,-1.3) {A 的本次上下文\\任务、已选资料与当前判断};
\node[context] (ctx-b) at (12.35,-1.3) {B 的本次上下文\\已有内容不会因共享写入而自动更新};
\node[align=center,text width=2.7cm,text=darkgray] at (7.7,-1.3) {上下文彼此独立\\没有自动同步};
\node[tool] (tool-a) at (3.05,-3.25) {A 的运行程序与获准工具};
\node[tool] (tool-b) at (12.35,-3.25) {B 的运行程序与获准工具};
\draw[localflow] (ctx-a.south)--node[labeltext,right]{形成可共享内容}(tool-a.north);
\draw[localflow] (tool-b.north)--node[labeltext,left]{④ 筛选后纳入}(ctx-b.south);
\draw[communication] (tool-a.east)--node[labeltext,above,text width=2.7cm]{② 通知 / 引用}node[labeltext,below,text width=2.7cm]{不等于内容或授权}(tool-b.west);

\node[rounded corners=4pt,draw=softpurple!80,fill=softpurple!9,
  text width=13.2cm,minimum height=1.1cm,align=center,inner sep=6pt]
  (store) at (7.7,-5.46) {\textbf{共享外部数据、记录与产物}\\可保存明确版本或可定位的内容；存储内容不等于任何实例的当前上下文};
\draw[dataflow] (tool-a.south)--node[labeltext,right,text width=4.05cm]{① 写入新内容\\保留可定位的版本或引用}(3.05,-4.89);
\draw[communication] ([xshift=-.67cm]tool-b.south)--node[labeltext,left,text width=3.5cm]{③ 显式读取\\核对可达性与权限}(11.68,-4.89);
\draw[dataflow] (13.02,-4.89)--node[labeltext,right,text width=1.55cm]{返回内容}([xshift=.67cm]tool-b.south);

\node[anchor=north west,align=left,text width=15.0cm] at (.2,-6.27)
 {每个实例分别受其账户、工具和数据权限约束；共用服务不意味着相同授权。\\存储可在本地或远程。若选取内容进入云端模型请求或远程工具，仍须检查该次外发的实际内容。};
\end{tikzpicture}
\caption{共享外部数据与独立上下文的边界。图示是一种更新交接顺序，不规定存储位置、数据表结构或统一版本机制。}
\label{fig:shared-data-contexts}
\end{figure}

共享存储不自动同步上下文。例如，实例 A 写入新结论后，B 当前使用的上下文仍可能包含旧信息。若后续判断依赖更新，应通过消息、事件或读取规则让 B 获知变化，并实际取得相关内容。共享记忆也需要明确可读范围、来源、适用时间与更新方式；某个实例写下的推测，不宜未经区分就成为所有实例共同采用的事实。


对并行任务，还要区分允许共同修改的数据与应固定的检查依据。多个实例读写同一资源时，一致性、事务和冲突处理仍遵循原存储的语义。需要固定输入时，可按任务选择快照、版本或不可变文件；需要协同修改时，则要避免后写结果悄然覆盖其他实例的有效修改。这些是工程选择，不是框架自动提供的统一制度。

独立上下文有助于控制信息范围，也可能使实例缺少他人已掌握的重要背景。交接因此应提供足够的任务信息、已确认事实和产物引用，并明确哪些仍待核实。共享工具或数据服务只表示复用能力，不表示每个实例都取得同等授权；引用可达性与数据权限也须分别检查。

跨实例与跨服务传递仍需检查出站边界。若本地资料经某个实例读取，再进入云端模型请求或远程工具，相应内容便传给该接收方；摘要也可能包含敏感信息。判断是否外发，应核对实际字段、附件、检索结果与日志，不能仅凭共享存储位于本地作结论。



\subsubsection{状态持久化与中断恢复的边界}
\label{sec:recovery}

单实例的状态保存、检查点与副作用恢复边界见第~\ref{sec:single-agent-recovery} 节。多 Agent 恢复还要重建任务之间的关系：哪些分支已完成，哪些等待输入或外部作业，哪些结果尚未送达，以及汇总方已经接收了什么。聊天记录不一定包含这些事实；多个实例各自有检查点，也不等于取得一致的系统快照。

\textbf{恢复不能把所有分支都当作未执行。}某个实例中断时，其他实例或长作业可能仍在工作。恢复协调方应先核实可获得的运行记录与外部状态，再决定继续等待、读取已有结果、恢复未完成步骤或重试。若状态未知，应保留未知，而不是直接标为失败并重新提交；否则可能重复计算、重复写入或造成其他副作用。

\textbf{返回结果需要关联到正确的工作。}重试、改派或修改输入后，原执行仍可能迟到返回。接收方需能够辨认结果对应的任务、输入或执行尝试，再判断是否纳入当前汇总；使用什么标识或版本机制由应用选择，不要求统一字段。重复消息不应被误算为多个完成结果，旧输入的成果也不能悄然覆盖已更新的结论。

\textbf{部分失败与取消需要协调。}某个分支失败后，可以保留已完成且仍有效的结果，按依赖决定哪些后续步骤暂停、返工或继续。已发送取消请求的分支，恢复时仍需核实实际终态；尚未停止的外部工作不能因为协调方重启而视作已结束。向用户交付部分结果时，应说明缺失内容及其对结论的影响。

共享数据库、消息队列或检查点各自提供的保障，不会自动组成跨系统事务；框架检查点也不覆盖全部外部作业与工具副作用。迁移到另一执行节点，还需满足依赖环境、数据引用可达性和访问权限。端到端验证应包含一次分支中断、迟到或重复结果，以及恢复后的汇总检查，确认系统能够继续原任务而非从头制造另一份工作。

\subsection{执行节点与共享服务}

\subsubsection{本地、公有云与托管运行环境}

执行循环、工具和计算作业可以放在本地，也可以部署到自有服务器、公有云机器或托管运行环境；不同部分还可以分开放置。本地执行加云端模型 API 是一种常见组合：程序与数据处理留在本地，必要上下文发送给模型服务。它不要求自建云服务器，也不表示本地必须部署大模型。

\textbf{模型 API 与托管 Agent 环境是两类能力。}前者主要接收推理输入并返回输出；后者还可能承载执行循环、会话、工具、沙箱或状态管理，具体范围须查产品接口。选择云端模型并不会把本地运行程序转移到云端；使用一个带 Agent 界面的应用，也不能据此认定它已开放可供其他系统调用的 Agent 执行接口。

把现有平台纳入协作，需确认它究竟提供推理、工具、作业还是完整 Agent 服务，以及身份认证、任务状态和结果交接是否可用。若只有交互界面或未核实的外部接口，就不能把它画成已接通的执行节点。部署位置影响持续在线、资源、隔离、网络和数据边界。第四章给出本项目的逻辑职责划分，附录 D 比较部署形态；本项目的节点分工及配置见第五章。

图~\ref{fig:execution-deployment} 将 Agent 执行位置与模型推理服务分开表示。这些环境是可选择、也可组合的逻辑部署形态；跨环境协作需另行具备可验证的接口与授权。

\begin{figure}[H]
\centering
\begin{tikzpicture}[x=1cm,y=1cm,font=\footnotesize\sffamily,text=darkgray,
  depoutline/.style={line width=.9pt,line cap=round,line join=round},
  deptitle/.style={font=\small\bfseries\sffamily,align=center},
  depdetail/.style={font=\footnotesize\sffamily,align=center},
  depmodel/.style={{Stealth[length=1.6mm]}-{Stealth[length=1.6mm]},draw=softpurple,line width=.9pt},
  depcoop/.style={{Stealth[length=1.6mm]}-{Stealth[length=1.6mm]},draw=darkgray!70,dashed,line width=.8pt}]
\path[use as bounding box] (0,0) rectangle (15.7,7.1);
% 三栏表示执行环境；虚线只表达接口可用且获授权时的可选协作。
\draw[depcoop] (2.5,6.23)--(2.5,6.63)--(7.85,6.63)--(7.85,6.23);
\draw[depcoop] (7.85,6.23)--(7.85,6.63)--(13.2,6.63)--(13.2,6.23);
\node[depdetail,fill=white,inner sep=2pt] at (7.85,6.97)
  {可选跨环境协作：核实接口、身份与授权，再交接任务、状态和结果};
\foreach \x/\col in {2.5/primaryblue,7.85/softgreen,13.2/accentorange}{
  \path[depoutline,draw=\col!65,fill=\col!3,rounded corners=5pt]
    (\x-2.4,2.0) rectangle (\x+2.4,6.23);
  \path[draw=\col!35,fill=white,rounded corners=4pt]
    (\x-2.12,3.27) rectangle (\x+2.12,5.15);
  \node[depdetail] at (\x,4.82) {一个或多个 Agent 实例};
  % 小循环图标强调程序在该环境中实际运行。
  \begin{scope}[shift={(\x-1.48,4.18)},draw=\col!85!black,line width=1pt]
    \draw[-{Stealth[length=1.3mm]}] (.20,.03) arc(8:280:.21);
    \draw[-{Stealth[length=1.3mm]}] (-.20,-.03) arc(188:100:.21);
  \end{scope}
  \node[deptitle,anchor=west,text=\col!85!black] at (\x-1.03,4.21) {Agent 执行循环};
}
% 本地设备图标。
\begin{scope}[shift={(.84,5.67)},draw=primaryblue,depoutline]
  \draw[fill=white,rounded corners=1pt] (-.36,-.17) rectangle (.36,.29);
  \draw (-.46,-.29)--(.46,-.29)--(.36,-.17)--(-.36,-.17)--cycle;
  \draw (-.12,-.23)--(.12,-.23);
\end{scope}
\node[deptitle,anchor=west,text=primaryblue] at (1.43,5.68) {本地设备};
\node[depdetail] at (2.5,3.62) {工具 · 上下文 · 运行状态};
\node[depdetail,text width=4.15cm] at (2.5,2.62)
  {程序与数据处理留在本机\\调用云端模型不改变执行位置};
% 自管公有云服务器图标。
\begin{scope}[shift={(6.15,5.68)},draw=softgreen!85!black,depoutline]
  \foreach \y in {-.20,.08}{
    \draw[fill=white,rounded corners=1pt] (-.34,\y) rectangle (.34,\y+.21);
    \fill (-.23,\y+.10) circle(.025);
    \draw (.02,\y+.10)--(.24,\y+.10);
  }
\end{scope}
\node[deptitle,anchor=west,text=softgreen!85!black] at (6.76,5.68) {自管公有云};
\node[depdetail] at (7.85,3.62) {工具 · 上下文 · 运行状态};
\node[depdetail,text width=4.15cm] at (7.85,2.62)
  {程序运行在租用的云机器上\\使用方维护应用与运行环境};
% 托管运行环境图标，不表示独占虚拟机或已开放外部 API。
\begin{scope}[shift={(11.48,5.68)},draw=accentorange!90!black,depoutline]
  \path[draw,fill=white] (-.31,-.20)
    ..controls(-.63,-.20)and(-.59,.20)..(-.31,.20)
    ..controls(-.23,.52)and(.26,.49)..(.32,.20)
    ..controls(.65,.20)and(.61,-.20)..(.35,-.20)--cycle;
  \draw(-.17,-.02)--(.20,-.02);
\end{scope}
\node[deptitle,anchor=west,text=accentorange!90!black] at (12.14,5.68) {提供方托管};
\node[depdetail] at (13.2,3.62) {会话、工具与状态按产品提供};
\node[depdetail,text width=4.15cm] at (13.2,2.62)
  {执行循环实际运行在托管端\\能力与外部接口须逐项核实};
% 推理输入与输出；模型服务位置是与 Agent 执行位置分开的部署轴。
\foreach \x in {2.5,7.85,13.2}{
  \draw[depmodel] (\x,1.98)--(\x,1.18);
  \node[depdetail,anchor=west,text=softpurple,inner sep=2pt] at (\x+.14,1.57) {推理输入 / 输出};
}
\path[depoutline,draw=softpurple!55,fill=softpurple!4,rounded corners=5pt]
  (.1,.03) rectangle (15.6,1.16);
% 模型图标表示推理能力，并非某一台物理服务器。
\begin{scope}[shift={(.88,.61)},scale=.40,draw=softpurple,depoutline]
  \path[draw,fill=white]
    (0,.60)..controls(-.23,.88)and(-.65,.72)..(-.66,.44)
    ..controls(-1,.44)and(-1.10,.02)..(-.84,-.16)
    ..controls(-1.03,-.50)and(-.71,-.78)..(-.45,-.62)
    ..controls(-.24,-.86)and(0,-.65)..(0,-.43)
    ..controls(0,-.65)and(.24,-.86)..(.45,-.62)
    ..controls(.71,-.78)and(1.03,-.50)..(.84,-.16)
    ..controls(1.10,.02)and(1,.44)..(.66,.44)
    ..controls(.65,.72)and(.23,.88)..(0,.60)--cycle;
  \draw (0,.60)--(0,-.43);
  \draw (-.65,.40)..controls(-.30,.44)and(-.19,.22)..(-.35,.04);
  \draw (.65,.40)..controls(.30,.44)and(.19,.22)..(.35,.04);
\end{scope}
\node[deptitle,text=softpurple,anchor=west] at (1.55,.80)
  {模型推理服务：各环境可分别选用，也可共同调用};
\node[depdetail,anchor=west] at (1.55,.35)
  {部署位置另定：供应商云端，或自托管于本地 / 云端；不要求与 Agent 同地或独占};
\end{tikzpicture}
\caption{Agent 执行位置与模型推理服务的关系示意。实线为推理调用，虚线为有条件的可选协作；下方服务框表示逻辑能力，不限定物理位置或单一模型。}
\label{fig:execution-deployment}
\end{figure}



\subsubsection{共享模型服务与模型适配}

模型服务支撑各实例的“脑”。多个实例可调用同一模型服务，各次请求仍携带各自的角色、任务与上下文；相同服务并不会使它们合并为一个 Agent。反之，一个实例也可按任务阶段使用不同模型，不因此变成多个运行实例。实例数量、模型种类与承载模型的硬件数量是不同维度。

共享服务还会带来共同的限流、费用、延迟和故障依赖。应用可按需求控制并发、安排重试或选择其他模型，但切换后应重新确认能力与结果质量。模型适配不仅是替换地址和名称，还涉及输入输出、工具调用、结构化结果、多模态、流式事件与错误语义。具体 SDK 或适配器的存在，不保证不同模型功能完全等价；生态比较见第~\ref{sec:agent-ecosystem} 节。

\subsubsection{共享工具服务与访问边界}

工具服务落实“手”的操作，任务和消息接口支持“口耳”的交流。多个实例可复用文件搜索、数据查询、构建测试或远程服务，不必为每个角色重建一套实现；但共享实现不等于共享授权，也不意味着所有调用可以同时修改同一份数据。

工具可在实例所在进程运行，也可由受限子进程、沙箱或远程服务执行。它能访问什么，由运行身份、输入限制、目标服务权限和执行环境共同决定。需要只读、隔离工作区或限制外发时，应落实为可检查的执行边界；一个描述为“审阅者”的角色，并不会因此自动失去写入权限。

共享工具有副作用时，还需考虑并发修改、资源争用和失败后重试。调用方至少应能判断返回的是最终结果、已受理状态还是错误，并保留足够的证据。外部工具或服务返回的说明不自动改变用户目标和原有授权；模型是否建议执行，与运行程序是否允许执行仍是两层判断。

\subsubsection{长计算作业与异步结果交接}

数据处理、构建、测试或算法训练可能超出一次短调用适合等待的时间。工具既可同步等待，也可先提交作业、返回关联标识，随后查询状态或接收事件。图~\ref{fig:async-compute-job} 展示后一种方式；它不意味着所有 Agent 框架都附带作业调度器或 GPU 管理能力。

\begin{figure}[H]
\centering
\begin{tikzpicture}[x=1cm,y=1cm,font=\footnotesize\sffamily,text=darkgray]
\path[use as bounding box] (0,-9.46) rectangle (15.4,.38);
\foreach \x in {1.75,5.72,9.70,13.67} {\draw[calllife] (\x,-.38)--(\x,-8.43);}
\node[callhead,minimum width=3.25cm] at (1.75,0) {Agent 运行程序};
\node[callhead,minimum width=3.25cm] at (5.72,0) {工具 / 作业服务};
\node[callhead,minimum width=3.25cm] at (9.70,0) {CPU / GPU 执行节点};
\node[callhead,minimum width=3.25cm] at (13.67,0) {产物存储};
\draw[callrequest] (1.75,-.88)--node[calllabel,above]{提交输入或引用与资源请求}(5.72,-.88);
\draw[callreturn] (5.72,-1.48)--node[calllabel,above]{先返回：已受理 + 作业标识}(1.75,-1.48);
\node[callstep,text width=3.0cm] at (5.72,-2.08) {\mbox{排队 / 准备}};
\draw[callrequest] (5.72,-2.76)--node[calllabel,above,pos=.76]{满足运行条件后启动}(9.70,-2.76);
\draw[callreturn] (9.70,-3.36)--node[calllabel,above]{运行状态 / 进度 / 日志引用}(5.72,-3.36);
\draw[callevent] (5.72,-3.97)--node[calllabel,above,text=accentorange!90!black]{可选状态事件：进度 / 日志}(1.75,-3.97);
\draw[callrequest] (9.70,-4.70)--node[calllabel,above]{写入结果文件或其他产物}(13.67,-4.70);
\draw[callreturn] (13.67,-5.31)--node[calllabel,above]{确认产物可读并返回引用}(9.70,-5.31);
\draw[callreturn] (9.70,-5.95)--node[calllabel,above]{终态：退出码 + 可读产物引用}(5.72,-5.95);
\draw[callevent] (5.72,-6.58)--node[calllabel,above,text=accentorange!90!black]{确认交付就绪：终态 + 引用}(1.75,-6.58);
\draw[callrequest] (1.75,-7.30)--node[calllabel,above]{按访问权限读取结果；作业标识本身不授予数据访问权}(13.67,-7.30);
\draw[callreturn] (13.67,-7.98)--node[calllabel,above]{返回获准内容或所需片段}(1.75,-7.98);
\node[callnote,anchor=north west,text width=14.7cm] at (.3,-8.53)
  {异常路径：失败保留退出码与日志；取消先请求，再停止实际执行，确认后才报告已取消。\\
   失联或超时且结果未知：先核实进程与产物，不盲目重复提交。};
\end{tikzpicture}
\caption{长计算作业的一种异步时序。示意采用状态事件；实际也可轮询，作业服务、存储及取消能力取决于工具和运行环境。}
\label{fig:async-compute-job}
\end{figure}

作业标识只是关联提交、状态与结果的一种常见约定，不要求统一字段名。“已受理”表示提交成功，“运行中”表示开始执行，“已完成”则应有实际终态与结果依据。图中先写入并确认可读取的产物，再发布完成事件；如果产物尚未可用，应继续标明发布中或未完成，不能仅凭进程退出就宣称交付已准备好。

运行作业的是计算环境，模型推理请求通常已经结束；结果到达后，程序才按需再次调用模型。CPU/GPU 数量与 Agent 数量没有一一关系，计算可位于本机、其他节点或外部服务。完整日志、大文件和中间结果可以留在工具侧，按需把指标、错误与获准片段送入模型上下文。

取消请求受理后，应继续核实实际执行是否终止；部分作业可能仍需清理或无法立即停止。失联或超时而状态不明时，先核对进程、服务记录与产物，再决定是否重交。框架检查点、计算检查点和外部作业状态分别覆盖不同范围，任何一方都不能单独证明全部工作已经恢复或安全结束。

\subsection{框架与平台生态}
\label{sec:agent-ecosystem}

\subsubsection{模型、代码框架、应用平台与托管平台}

现有产品可按主要职责分为四类：\textbf{模型与 API} 提供推理；\textbf{代码框架}组织模型、工具、上下文与执行流程；\textbf{应用平台}把配置、知识库、流程编辑和发布做成可操作产品；\textbf{托管运行平台}承载部署、运行资源、隔离与观测。这是理解生态的视角，不是互斥分类：框架可能附带运行时，平台也可能同时提供代码 SDK 和模型服务。

选模型、选框架与选部署位置，是有关联但不同的决定。安装框架只获得开发组件；具体角色、任务、工具、数据、权限和检查规则仍需组合成应用。以下依据截至 2026-10-12 的官方资料，说明代表性路线；适用场景与取舍是一般工程判断，不是产品排名或预先确定的软件组合。

\subsubsection{代码框架的主要路线与适用场景}

显式状态图适合需要清楚依赖、等待和恢复的任务；角色与团队抽象便于快速组织协作；轻量执行循环便于与现有程序组合。它们并非互斥路线，不少框架同时提供多类能力。下表用同一组维度比较，名称相近或都支持多 Agent，不代表执行语义与恢复范围相同。

\begingroup
\small
\renewcommand{\arraystretch}{1.16}
\setlength{\tabcolsep}{4pt}
\begin{longtable}{P{2.7cm}P{3.7cm}P{3.15cm}P{5.15cm}}
\caption{代表性 Agent 代码框架的定位与取舍}\label{tab:agent-frameworks}\\
\toprule
\rowcolor{primaryblue}\color{white}\textbf{框架} & \color{white}\textbf{核心抽象} & \color{white}\textbf{适合重点试验的任务} & \color{white}\textbf{需要核实与补齐的部分}\\
\midrule
\endfirsthead
\toprule
\rowcolor{primaryblue}\color{white}\textbf{框架} & \color{white}\textbf{核心抽象} & \color{white}\textbf{适合重点试验的任务} & \color{white}\textbf{需要核实与补齐的部分}\\
\midrule
\endhead
\bottomrule
\endfoot
LangGraph\newline\cite{langgraph-persist,langgraph-overview} & 状态图、节点、条件边及 checkpoint；显式组合确定性步骤和模型步骤 & 有并行、返工、人工等待及恢复要求的长任务 & 控制细，但需自行设计图、业务状态、幂等性和工具权限；checkpoint 不等于业务事务\\
\rowcolor{lightgray}
CrewAI 与 Flows\newline\cite{crewai-flows} & Agent、Task、Crew 表达角色协作；Flows 组织事件、状态和路由 & 快速试做研究员、开发者、审阅者等角色分工 & 角色配置不能替代交付契约；复杂分支、恢复与流程约束需在 Flows 和业务代码中明确\\
OpenAI Agents SDK\newline\cite{agents-sdk} & Agent、Runner、工具、handoff、session 与追踪 & 以 Python 快速组织工具循环、专家委派与小型协作 & 抽象较轻；业务任务调度、跨进程恢复和副作用控制仍需设计，session 历史不等于完整业务恢复\\
\rowcolor{lightgray}
Microsoft Agent Framework\newline\cite{microsoft-agent} & Agent、workflow、会话状态及模型提供方抽象 & 已有 Python / .NET 系统，需要把 Agent 与应用流程整合 & 需核对所用语言、提供方和部署方式的能力差异；既有 AutoGen / Semantic Kernel 代码需评估迁移\\
Google ADK\newline\cite{google-adk} & Agent、工具及会话；Python 2.0 提供图式 workflow runtime 与 Task API & 需要显式流程、Agent 间任务委派，或接入 Google 生态 & 不能把旧版教程等同当前接口；第三方模型与具体运行后端需逐项验证\\
\rowcolor{lightgray}
AgentScope\newline\cite{agentscope} & Agent 循环、模型、工具、上下文、权限与团队协作；2.0 集成服务运行能力 & 多模型、多工具协作，以及从本地原型逐步扩展服务 & 可接多家模型，不限通义；业务持久化、隔离边界与团队交接仍需配置和验证\\
\mbox{Qwen-Agent}\newline\cite{qwen-agent} & 围绕 Qwen 的 Agent、函数工具、MCP、RAG 与代码解释器组件 & 以通义模型为主的工具调用与资料处理原型 & 官方路径为 DashScope 或自建兼容接口的 Qwen 服务；不能据此承诺任意模型等价适配或自动具备生产隔离\\
\end{longtable}
\endgroup

\textbf{同一生态也会提供更高层封装。}Deep Agents 建立在 LangChain / LangGraph 之上，把文件工具、上下文管理、子 Agent 和可选规划等能力组合成 agent harness，可用于试做多步骤研究或编码助手；工具权限、工作区与验收条件仍需明确\cite{deepagents}。Microsoft Agent Framework 是 AutoGen 与 Semantic Kernel 的直接后继；AutoGen 处于维护模式，新项目需按现行路线评估\cite{microsoft-agent,autogen-status}。AgentScope 2.0 已整合原 AgentScope Runtime 的能力，不宜沿用旧教程将二者视为必装的两套组件\cite{agentscope-runtime}。

\subsubsection{应用平台、预设流程与托管能力}

\textbf{Dify、n8n 与百炼}更靠近应用交付。Dify 通过工作流节点组织模型与工具，适合知识问答和业务助手原型；其新版 Agent 节点当前仍标为 beta，且仅用于 Workflow 应用\cite{dify-agent}。n8n 用自动化工作流连接应用、事件和 AI 节点，适合资料收集、定时触发和结果分发等外围流程\cite{n8n-ai}。阿里云百炼提供模型接入，以及智能体、可视化工作流和高代码应用路径，可在其平台内构建和发布应用\cite{bailian-apps}。

预设流程把步骤和依赖事先配置，Agent 节点则可在允许范围内动态选择工具或处理路径，二者可以组合。可视化编排能够降低配置成本，但不能仅凭界面推定已经覆盖长任务恢复、版本化产物或独立测试；复杂状态、并发与外部副作用仍需逐项核验。

\textbf{阿里云 AgentRun}是托管运行产品名，与第二章用于解释运行实例的对象名称含义不同。它主要提供 Agent 部署、沙箱、模型治理与观测等基础设施能力，可与代码框架组合，也包含部分低代码入口\cite{agentrun}。是否采用托管服务，取决于持续在线、隔离、弹性和运维要求；不能因为使用某个框架，就推导出必须使用其关联云平台。地域、模型与功能可用性还需按实际账户核实。

\subsubsection{模型兼容性与框架评估维度}

\textbf{框架厂商通常不等于唯一模型厂商。}OpenAI Agents SDK 提供非 OpenAI 模型的接入点；Microsoft Agent Framework、Google ADK 和 AgentScope 也有多提供方路径\cite{agents-models,microsoft-providers,adk-models,agentscope}。但接口兼容不等于功能等价，函数调用、结构化输出、多模态、流式事件、托管工具和错误语义，都可能随模型与适配层变化。Agents SDK 的 LiteLLM / Any-LLM 集成目前标为 best-effort beta，适配器列出供应商并不保证全部能力可用\cite{agents-models}。

框架评估宜同时检查任务完成质量、流程控制、上下文管理、模型替换成本、工具权限、持久化与恢复范围、观测能力、部署要求、调用费用和维护负担。代码框架可承载核心任务，应用平台可提供入口与自动化，托管平台可提供运行环境；这些是可组合的选择，而不是固定分层或必装清单。

可使用同一小型研发任务比较候选方案：并行完成资料研究与算法实现，由独立测试发现缺陷并返工，中途重启后继续，最后追溯到代码、数据、测试证据与产物。试验还应确认部分失败、重复调用和取消的实际行为。安装成功、演示跑通或多个 Agent 得出一致结论，都不能替代这些检查；比较结果才是后续部署与选型的依据。


% ════════════════════════════════════════════════
% ════════════════════════════════════════════════
\section{面向资产管理的研发系统总体设计}
\label{sec:rd-system-design}
% ════════════════════════════════════════════════

第二、三章解释单个与多个 Agent 的一般机制。本章据此提出本项目的研发系统设计：由人确定目标和重要取舍，研发协调角色负责版本推进，开发、算法、独立测试与运维角色在程序支持下完成研究、实现、验证和交付。本章中的组件、职责与数据约定是拟采用的设计，不表示系统已经建成；具体软件、安装方式与部署组合在第五章展开。

\subsection{建设目标与系统范围}

图~\ref{fig:rd-goal-concept} 先说明本项目要建设什么：把人的研究意图与约束转化为可查证、可复现的研发成果，并明确资产运营与实盘交易不在本系统内。

\begin{figure}[H]
\centering
\begin{tikzpicture}[x=1cm,y=1cm,font=\footnotesize\sffamily,text=darkgray,
  gcbox/.style={rounded corners=4pt,draw=primaryblue!45,fill=white,align=center,inner sep=5pt},
  gcflow/.style={-{Stealth[length=1.8mm]},draw=primaryblue,line width=.85pt},
  gclabel/.style={font=\scriptsize\sffamily,fill=white,inner sep=2pt,align=center}]
\path[use as bounding box](0,0)rectangle(15.4,7.75);
% 人的目标位于系统边界外。
\draw[primaryblue,line width=1.1pt] (3.25,7.30) circle(.16);
\draw[primaryblue,line width=1.1pt] (3.25,7.13)--(3.25,6.72) (2.97,6.98)--(3.53,6.98) (3.25,6.72)--(3.02,6.45) (3.25,6.72)--(3.48,6.45);
\node[gcbox,text width=7.6cm,minimum height=.84cm] (human) at (8.15,6.98) {\textbf{人类负责人提出目标与约束}\\研究问题、交付预期、资源边界、关键取舍};
\path[rounded corners=5pt,draw=primaryblue!55,fill=primaryblue!3] (.25,2.26)rectangle(15.15,5.92);
\node[anchor=north west,font=\small\bfseries\sffamily,text=primaryblue] at (.47,5.80) {拟建设的研发系统};
\node[gcbox,fill=primaryblue!7,text width=9.6cm,minimum height=.64cm] (coord) at (7.70,5.03) {理解需求与版本组织：明确任务、责任、依赖与验收依据};
\draw[gcflow](human.south)--node[gclabel,right]{目标 / 约束}(human.south |- coord.north);
\draw[gcflow,dashed](13.95,5.05)--(13.95,6.98)--node[gclabel,above]{待决事项 / 结果}(human.east);
\foreach \x/\id/\title/\detail in {2.87/research/研究与知识组织/来源、定义与可复核结论,7.70/build/数据算法与软件研发/加工、实验、实现与集成,12.53/check/验证与研发运行/独立检查、交接与反馈}{
  \node[gcbox,text width=3.62cm,minimum height=1.03cm] (\id) at (\x,3.79) {\textbf{\title}\\\detail};
  \draw[gcflow](coord.south)--(7.70,4.42)-|(\id.north);
}
\draw[gcflow](research.south)--(2.87,2.99)--(7.70,2.99)--(7.70,1.95);
\draw[gcflow](build.south)--(7.70,2.99);
\draw[gcflow](check.south)--(12.53,2.99)--(7.70,2.99);
\node[font=\scriptsize\sffamily,align=center,fill=primaryblue!3,inner sep=3pt] at (7.70,2.67) {全过程保留任务、来源、版本与验证证据；模型提出行动，程序落实权限与真实执行};
\draw[draw=softgreen!70!black,line width=.85pt](2.0,1.95)--(13.40,1.95);
\foreach \x/\id/\title/\detail in {2.0/report/研究报告/结论可查证,5.8/data/数据流程/输入与处理可复现,9.6/algo/算法结论/实验条件与局限明确,13.4/software/软件版本/功能与接口经验证}{
  \node[gcbox,draw=softgreen!60,fill=softgreen!6,text width=3.10cm,minimum height=.88cm] (\id) at (\x,1.22) {\textbf{\title}\\\detail};
  \draw[gcflow,draw=softgreen!70!black](\x,1.95)--(\id.north);
}
\node[font=\scriptsize\sffamily,text=softred,align=center] at (7.70,.22) {系统边界外：账户资产运营、资金划转、投资决策执行与实盘交易，均需另行设计与授权};
\end{tikzpicture}
\caption{本项目的目标与范围：人的意图进入研发系统，形成有版本和证据的研究、数据、算法与软件成果。}
\label{fig:rd-goal-concept}
\end{figure}

\subsubsection{服务目标与预期产出}

系统服务于资产管理所需能力的研发。建设目标是把一个待研究的问题或待实现的需求，转化为可以查证、复现和继续使用的成果，并使人能够掌握当前进展、关键风险与待决事项。例如，“能否从指定资料中整理企业之间的关系”需要逐步形成资料范围、关系定义、提取方法、验证样本和可运行原型，而不只是生成一份看似完整的关系清单。

预期产出包括有来源的研究报告、可重复的数据处理流程、带实验条件的算法结论、经过验证的软件版本，以及相应的操作说明和运行记录。不同产出适用不同的完成标准：资料研究需要覆盖问题并区分事实与推断，软件需要满足接口和功能规格，算法实验则需要说明假设、数据、比较方法与证据。研究得出“不宜采用当前方法”也可以是有效成果，前提是证据足以支持这一结论。

系统建设本身也需要可验证：人能否从一个需求定位其任务、候选成果和验收依据；中断后能否识别实际状态并继续；跨角色交接是否遗漏必要信息。这些是后续试点应检查的能力，不能用 Agent 数量或自动生成的文字量替代。

\subsubsection{系统承担的工作与外部边界}

本系统承担资料研究、知识组织、数据加工、算法实验、软件开发、测试验证和研发环境运维。它可以为后续资产管理应用形成研究方法、数据能力和软件组件，但不承担账户资产运营、资金划转、投资决策执行或实盘交易。将一个研发成果用于这些业务，需要另行设计业务系统、使用条件与授权边界。

外部输入包括人提交的目标、获准使用的资料与数据、已有代码和运行约束；外部输出包括交付包、证据报告、运行反馈与需要人决策的问题。外部模型服务提供推理能力，工具与计算环境完成真实操作，资料提供方保有其数据使用条件。系统不能因为能够读取某项资料，就默认具有复制、外传或长期留存的权利。

\subsubsection{人机协作方式与设计约束}

人通过统一入口提出意图、查看证据、回答澄清问题并决定关键取舍。系统负责把意图整理为研发任务、安排执行、核对交接和返回结果；正常情况下不要求人把各 Agent 的消息逐条转发。遇到目标冲突、预算或权限不足、范围变化等问题，协调角色提交简明的决策请求，列出选项、影响与建议，获确认后再推进依赖该决定的工作。

本设计遵守三个约束。首先，职责与权限分开配置，模型提出行动，运行程序及实际账户权限决定能否执行。其次，成果以可定位的版本和证据交付，聊天总结不是唯一记录。最后，组件先按逻辑职责划分，不预设每个组件都是独立服务或必须分布在不同机器上。本地、云端及混合部署均可承载这套设计。

人可使用 VS Code、Codex 以及 dot 云电脑辅助编写、审阅、测试和多 Agent 任务协作。这些是建设与使用研发能力的工作环境，其中 dot 云电脑优先承担文档、LaTeX、简单 Web 和轻量代码任务。它们与 Mac 上自行部署的 Agents SDK 控制程序分开，通过 GitHub 中经审阅的代码与文档修订交接；不能据此假定 dot 自动注册为平台 Worker、始终在线，或提供可由本平台调用的外部任务 API。人工或托管环境产生的修改同样须进入项目版本与验收记录。

\subsection{研发任务与使用场景}

图~\ref{fig:rd-scenarios} 从使用场景展开研发系统。各类任务共享同一套平台与受控资源，可以单独进行，也可以按成果依赖组合，不要求每个项目都走完所有场景。

\begin{figure}[H]
\centering
\begin{tikzpicture}[x=1cm,y=1cm,font=\footnotesize\sffamily,text=darkgray,
  scbox/.style={rounded corners=4pt,draw=primaryblue!45,fill=primaryblue!3,align=center,inner sep=5pt},
  scflow/.style={-{Stealth[length=1.7mm]},draw=primaryblue,line width=.85pt},
  sclabel/.style={font=\scriptsize\sffamily,fill=white,inner sep=2pt,align=center}]
\path[use as bounding box](0,0)rectangle(15.4,7.70);
\node[font=\small\bfseries\sffamily,text=primaryblue] at (7.70,7.35) {按目标选择场景，可单独启动，也可按需组合};
\foreach \x/\id in {1.85/research,5.75/algo,9.65/dev,13.55/ops}{
  \node[scbox,text width=3.12cm,minimum height=2.90cm] (\id) at (\x,5.35) {};
}
% 文档、数据、代码与检查图标帮助识别场景。
\draw[primaryblue,line width=1pt,fill=white] (1.54,6.02)rectangle(2.16,6.62);
\foreach \yy in {6.18,6.31,6.44}{\draw[primaryblue!75] (1.66,\yy)--(2.05,\yy);}
\draw[primaryblue,line width=1pt,fill=white] (5.42,6.15)--(5.42,6.56)arc(180:0:.33 and .10)--(6.08,6.15)arc(0:-180:.33 and .10);
\draw[primaryblue,line width=1pt] (5.42,6.56)arc(180:360:.33 and .10) (5.42,6.35)arc(180:360:.33 and .10);
\draw[primaryblue,line width=1pt] (9.29,6.16)rectangle(10.01,6.63) (9.15,6.07)--(10.15,6.07);
\node[font=\small\bfseries\ttfamily,text=primaryblue] at (9.65,6.39) {< >};
\draw[primaryblue,line width=1pt] (13.22,6.05)rectangle(13.88,6.65);
\draw[softgreen!75!black,line width=1.2pt] (13.32,6.34)--(13.49,6.19)--(13.79,6.52);
\foreach \x/\title in {1.85/资料研究与知识组织,5.75/数据处理与算法实验,9.65/软件开发与能力集成,13.55/测试验证与运行维护}{\node[font=\footnotesize\bfseries\sffamily,align=center] at (\x,5.82) {\title};}
\node[align=center,text width=3.12cm] at (1.85,4.79) {问题 + 获准资料\\$\downarrow$\\报告、索引、\\关系证据};
\node[align=center,text width=3.12cm] at (5.75,4.79) {数据 + 假设 + 口径\\$\downarrow$\\数据集、方法\\与实验记录};
\node[align=center,text width=3.12cm] at (9.65,4.79) {需求 + 接口 + 方法\\$\downarrow$\\代码、组件\\与候选版本};
\node[align=center,text width=3.12cm] at (13.55,4.79) {候选成果 + 运行条件\\$\downarrow$\\验收、环境\\与运行证据};
\foreach \left/\right in {research/algo,algo/dev,dev/ops}{\draw[scflow,<->](\left.east)--(\right.west);}
\node[scbox,fill=primaryblue!8,text width=14.25cm,minimum height=.87cm] (platform) at (7.70,2.91) {\textbf{共享研发平台}\\统一入口、版本与任务、角色协作、接收确认、成果登记与证据引用};
\foreach \id in {research,algo,dev,ops}{\draw[scflow,<->](\id.south)--(\id.south |- platform.north);}
\node[scbox,draw=softgreen!55,fill=softgreen!5,text width=14.25cm,minimum height=.92cm] (resource) at (7.70,1.40) {\textbf{共享能力与受控资源}\\模型与工具、计算作业、资料数据、代码和运行环境；按任务授权并隔离写入};
\draw[scflow,<->](platform.south)--node[sclabel,right]{调用 / 记录 / 复用}(resource.north);
\node[font=\scriptsize\sffamily,align=center,text width=15cm] at (7.70,.20) {横向箭头表示按实际依赖交接资料、方法、接口、候选成果与问题反馈，不表示四类场景必须依次执行。};
\end{tikzpicture}
\caption{研发场景总体关系：不同任务共享平台与资源，按实际依赖组织成果交接。}
\label{fig:rd-scenarios}
\end{figure}

\subsubsection{资料研究与知识组织}

这类任务从一个明确问题和资料范围开始。例如，整理某行业的参与主体及其关系时，输入应包括研究问题、允许读取的资料集合、资料时点与需要回答的关系类型。算法角色负责信息提取方法、分类与关联规则，开发角色负责采集、解析和检索能力；资料规模较小时，也可由一个实例完成多项工作，但其任务责任仍须明确。

输出需要保留结论到原文的对应关系，说明来源、定位信息和不确定性。两份资料对同一事实描述不一致时，应保存差异及资料时点，不能通过一段合并摘要消除冲突。知识图谱、索引或资料库是知识组织的候选形式，采用哪种形式应由查询需求决定。例如，只要求查找原文时，全文索引可能足够；需要沿明确关系追踪来源时，关系结构才有直接价值。

\subsubsection{数据处理与算法实验}

数据处理任务接收确定的数据范围、字段含义与质量要求，产出处理后的数据集、转换程序及质量报告。开发角色负责读取、转换和接口实现，算法角色负责影响研究结论的数据口径与统计假设，运维角色提供受控的存储和计算环境。这里的数据职责由现有岗位协作承担，不另设数据部门。

算法实验以待检验假设为起点，把数据快照、基线方法、评价方法、参数与资源限制组成实验计划。输出不只是一个分数，还包括实验记录、代码版本、作业状态、误差分析与适用条件。例如比较两种实体匹配方法，应在固定样本上观察误合并和漏合并，而不能以某次演示成功代表整体有效。涉及时间序列研究时，还应在方案中约定样本切分和可获得信息的时点，避免把未来信息带入历史判断。

\subsubsection{软件开发与能力集成}

软件开发把稳定的需求和接口约定转化为代码、配置及可运行组件。开发组负责人负责技术方案、模块边界和集成安排；具体开发实例在分配的分支或工作副本中实现并自测。算法提供的实验代码若要成为系统能力，还需要明确输入输出、错误处理、运行依赖和资源消耗，不能直接把研究脚本视为可交付服务。

例如，把关系提取能力接入资料工作空间，开发任务除了调用提取方法，还包括资料标识传递、失败反馈、结果查询与证据定位。交接给独立测试时，应同时提交候选版本、接口说明、自测结果和已知限制。依赖尚未准备好时，可以先用约定样例验证接口，但必须标明模拟输入，最终仍需完成真实依赖下的集成检查。

\subsubsection{测试验证与运行维护}

测试从版本规划阶段参与，提前识别验收对象、测试数据与可验证条件。独立测试角色在候选版本提交后执行独立验收，检查功能、数据质量、方法证据和必要的异常场景；开发、算法和运维各自的自测是交付前提，不能替代独立验收。

运维角色负责研发环境可用性、作业运行观察、日志、备份恢复和经批准的研发成果部署。其输入是环境需求、候选发布包与运行限制，输出是环境记录、运行证据和交接说明。例如计算作业中断时，运维先核实作业是否仍在运行、是否已经写出结果，再报告恢复条件；不能把重新启动实例等同于恢复原实验。本章所说运行维护始终限于研发系统及获准的试运行环境。

\subsection{组织角色与系统职责}

图~\ref{fig:rd-roles} 展示版本责任、专业分工和交接关系。岗位定义持续存在，执行实例按任务安排；协调负责组织推进，专业角色负责相应技术判断，重要取舍由人决定。

\begin{figure}[H]
\centering
\begin{tikzpicture}[x=1cm,y=1cm,font=\footnotesize\sffamily,text=darkgray,
  robox/.style={rounded corners=4pt,draw=primaryblue!50,fill=white,align=center,inner sep=5pt},
  roflow/.style={-{Stealth[length=1.8mm]},draw=primaryblue,line width=.85pt},
  rolabel/.style={font=\scriptsize\sffamily,fill=white,inner sep=2pt,align=center}]
\path[use as bounding box](0,0)rectangle(15.4,8.12);
\node[robox,fill=lightgray,text width=8.6cm,minimum height=.78cm] (human) at (7.70,7.61) {\textbf{人类负责人}：目标、优先级、资源边界与重要取舍};
\node[robox,fill=primaryblue!8,text width=8.6cm,minimum height=.90cm] (coord) at (7.70,6.01) {\textbf{研发协调角色：版本负责人}\\维护范围与里程碑，安排依赖、核对交接、汇总验收与待决事项};
\draw[roflow,<->](human.south)--node[rolabel,right]{目标与授权 / 待决事项}(coord.north);
\draw[roflow,<->](coord.south)--(7.70,5.05);
\draw[draw=primaryblue,line width=.85pt](1.85,5.05)--(13.55,5.05);
\node[rolabel] at (7.70,5.05) {任务与专业意见；技术判断由相应专业角色负责};
\path[rounded corners=4pt,draw=primaryblue!50,fill=primaryblue!3] (.13,2.35)rectangle(3.57,4.55);
\node[align=center,text width=3.18cm] (devlead) at (1.85,4.02) {\textbf{开发组负责人}\\技术方案、接口与集成};
\node[align=center,text width=3.18cm] (dev) at (1.85,2.91) {\textbf{开发执行角色}\\代码实现、集成与自测};
\draw[roflow](devlead.south)--(dev.north);
\node[robox,fill=primaryblue!3,text width=3.08cm,minimum height=2.20cm] (algo) at (5.75,3.45) {\textbf{算法角色}\\[6pt]数据口径与研究假设\\方法实现与实验\\结论、局限及自测证据};
\node[robox,draw=softgreen!70,fill=softgreen!5,text width=3.08cm,minimum height=2.20cm] (test) at (9.65,3.45) {\textbf{独立测试角色}\\[6pt]提前制定验收计划\\独立运行检查与复核\\缺陷、证据及验收意见};
\node[robox,fill=primaryblue!3,text width=3.08cm,minimum height=2.20cm] (ops) at (13.55,3.45) {\textbf{运维角色}\\[6pt]研发环境与资源准备\\作业观察、日志与恢复\\经批准的部署与交接};
\draw[roflow,<->](1.85,5.05)--(1.85,4.55);
\foreach \id in {algo,test,ops}{\draw[roflow,<->](\id.north |- 0,5.05)--(\id.north);}
\node[robox,draw=softgreen!55,fill=softgreen!5,text width=14.23cm,minimum height=.75cm] (handoff) at (7.70,1.23) {\textbf{跨角色交接}：明确输入与版本 $\rightarrow$ 提交成果及证据 $\rightarrow$ 接收方核对确认；缺陷回到责任任务};
\draw[roflow,<->](1.85,2.35)--(1.85,1.63);
\foreach \id in {algo,test,ops}{\draw[roflow,<->](\id.south)--(\id.south |- 0,1.63);}
\node[font=\scriptsize\sffamily,align=center,text width=15cm] at (7.70,.23) {图中是持续的岗位职责，不是固定 Agent 数量或机器节点；实例可按任务增减和接续，测试独立性由职责、过程与证据落实。};
\end{tikzpicture}
\caption{组织角色与责任交接：协调负责版本推进，专业角色负责各自判断，重要取舍提交人决定。}
\label{fig:rd-roles}
\end{figure}

\subsubsection{人类负责人与研发协调角色}

人类负责人确定建设方向、阶段目标、优先级、可接受的资源投入与重要取舍，并对需要新增授权或改变目标的事项作出决定。人可以审阅技术方案，但不必承担每项技术细节的日常裁决。系统应把决策问题说明白，例如继续提高识别质量需要增加多少工作、会推迟哪些既定成果，而不是把全部执行日志交给人自行判断。

研发协调角色对一个版本的完整推进负责。它维护版本范围与里程碑，组织方案讨论，安排跨组任务和依赖，跟踪阻塞，核对交接条件并汇总验收结论。协调 Agent 可以由不同运行实例接续承担，版本责任却不能随某个实例结束而消失，因此版本状态必须记录在共享的任务管理中。

开发组负责人负责实现层面的技术细节，包括结构设计、接口、代码组织和集成路径；算法、测试与运维负责人分别承担本专业判断。协调角色组织这些意见并处理进度与范围的冲突，不能仅凭协调岗位身份替代专业验证；涉及改变目标或重大取舍时，再提交人类负责人决定。

\subsubsection{开发、算法、测试与运维角色}

四类岗位按其交付责任分工，各组可根据工作量设负责人和执行实例。它们通过任务与成果交接协作，不以聊天对象的数量定义组织规模。

\begin{table}[htbp]
\centering\small\renewcommand{\arraystretch}{1.22}
\begin{tabularx}{\textwidth}{P{1.6cm}P{3.5cm}L}
\toprule
\textbf{角色} & \textbf{主要输入} & \textbf{职责与应交付的证据}\\
\midrule
开发 & 需求、技术方案、接口及数据样例 & 实现代码与集成，自测后交付候选版本、接口说明和构建结果。\\
\rowcolor{lightgray}
算法 & 研究问题、数据口径、实验约束 & 提出方法并开展实验，交付实现、实验记录、结论及局限。\\
独立测试 & 冻结需求、验收计划、候选成果 & 自行执行测试与复核，交付缺陷、复现材料及独立验收意见。\\
\rowcolor{lightgray}
运维 & 资源需求、发布包、运行约束 & 准备和维护研发环境，交付运行、备份恢复及部署交接记录。\\
\bottomrule
\end{tabularx}
\caption{研发岗位的主要输入与交付责任}
\end{table}

交接需要由接收方核对并确认。例如，算法组提交“方法可用”后，开发组应检查接口是否明确、依赖是否可获得、样例能否运行，再接受集成任务。若失败来自算法输出格式不符，回到算法任务修正；若输出符合约定而调用端处理错误，则由开发修复。归属暂不清楚时，由协调角色组织共同复现，避免缺陷只在不同岗位之间反复转交。

\subsubsection{岗位角色、运行实例与执行节点}

本项目使用持续的岗位定义和按任务启动的运行实例。一个算法角色可启动多个实例，分别研究不同方法；一个开发实例也可依次完成若干相关任务。实例拥有自己的上下文、当前状态与获准工具范围，执行节点则承载实例程序、工具进程或计算作业。角色、实例和节点的标识应分开记录。

例如，算法实例甲和乙可以运行在同一台服务器，并调用同一个模型服务；独立测试实例也可以使用这个模型。它们是多个独立运行的“脑”，不要求采购多套不同模型。共享模型并不意味着共享全部上下文或权限，也不保证判断独立；测试的独立性要由独立的职责、检查过程和证据来落实。

并行工作的隔离应体现在实际资源中：代码修改使用独立分支或工作副本，实验输出使用独立目录与作业标识，固定输入只读引用，写入共享成果需要明确提交步骤。同一模块确需共同修改时，由开发负责人安排合并顺序与冲突处理。增加实例以前，先确认数据依赖、资源限额和写入边界允许并发。

\subsubsection{决策权限与独立验证}

版本内的日常技术选择由相应专业角色在约定范围内作出，协调角色负责跨组安排和版本一致性。新增需求、扩大数据使用范围、增加超出已定限额的资源投入或改变重要目标，需要提交人类负责人。模型提示词说明这些规则，但权限必须由工具封装、运行账户、资源配额和审批检查在执行侧限制。

独立测试不得自行修改候选成果后宣布其通过；发现问题应保存失败证据并退回责任任务。修复后形成新的候选修订，测试针对该修订重新检查。验收证据由确定性验证程序与测试 Agent 配合产生：程序检查格式、约束、测试断言和输出一致性，测试 Agent 设计与解释检查、识别遗漏并形成意见。两者均有局限，重要结论还需要抽样复核或人类裁决。

测试在规划时可以提出可测性要求和验收样例，却不因此成为实现方。实现组的自测记录可供参考，验收仍须针对固定的需求、数据和候选成果自行运行。对于研究任务，测试重点是方法是否符合计划、证据能否复现、结论是否超过证据，而不是要求每项假设都被证实。

\Needspace{12cm}
\subsection{系统组成与总体交互}

图~\ref{fig:rd-overall-design} 先给出组件与交互全貌，再逐项说明入口、协作控制、运行层、资源和成果管理的职责。图中的组件是逻辑划分，不预设独立服务或固定部署位置。

\begin{figure}[H]
\centering
\begin{tikzpicture}[x=1cm,y=.70cm,font=\scriptsize\sffamily,text=darkgray,
  rdarea/.style={rounded corners=5pt,draw=primaryblue!50,fill=primaryblue!2},
  rdbox/.style={rounded corners=3pt,draw=primaryblue!55,fill=white,align=center,inner sep=4pt},
  rdagent/.style={rdbox,text width=2.30cm,minimum height=.80cm,font=\scriptsize\sffamily},
  rdexternal/.style={rdbox,draw=darkgray!50,fill=lightgray,text width=2.25cm,minimum height=.68cm},
  rdflow/.style={-{Stealth[length=1.8mm]},draw=primaryblue,line width=.85pt},
  rdref/.style={-{Stealth[length=1.8mm]},draw=softgreen!70!black,dashed,line width=.85pt},
  rdlabel/.style={font=\scriptsize\sffamily,fill=white,inner sep=2pt,align=center}]
\path[use as bounding box](-.05,-12.0)rectangle(15.5,.55);
\path[rdarea](0,-.85)rectangle(11.65,-10.40);
\node[font=\small\bfseries\sffamily,text=primaryblue,anchor=north west] at (.20,-1.03) {研发系统};
\node[rdexternal,text width=5.8cm,minimum height=.64cm] (human) at (5.82,.05) {人类负责人：目标、约束、关键取舍与验收};
\node[rdbox,text width=7.2cm,minimum height=.68cm] (entry) at (5.82,-2.02) {用户入口与项目工作空间\\项目约定、需求、进展与成果导航};
\node[rdbox,text width=7.2cm,minimum height=.70cm,fill=primaryblue!8] (coord) at (5.82,-3.48) {研发协调 Agent + 协作控制程序\\版本负责、任务依赖、交接与待决事项};
\foreach \x/\id/\name/\detail in {1.48/dev/开发 Agent/技术方案与实现,4.37/algo/算法 Agent/研究与实验,7.27/test/独立测试 Agent/规划与独立验收,10.16/ops/运维 Agent/环境与运行交接}{
  \node[rdagent] (\id) at (\x,-5.43) {\textbf{\name}\\\detail};
}
\node[rdbox,text width=10.1cm,minimum height=.72cm,fill=softblue!8] (runtime) at (5.82,-7.25) {Agent 运行与资源接入\\协调及专业实例的上下文、模型与工具接入、权限及作业控制};
\node[rdbox,text width=4.56cm,minimum height=.89cm,fill=softgreen!6] (state) at (2.94,-9.33) {\textbf{任务与协作记录}\\版本、依赖、状态、消息关联\\项目执行状态的权威来源};
\node[rdbox,text width=4.56cm,minimum height=.89cm,fill=softgreen!6] (art) at (8.69,-9.33) {\textbf{数据、代码与成果管理}\\修订、来源、快照、验收证据\\交付内容的权威引用};
\node[font=\small\bfseries\sffamily,text=darkgray] at (13.72,-3.61) {外部资源与服务};
\node[rdexternal] (model) at (13.72,-5.22) {模型服务\\可供实例共享};
\node[rdexternal] (tool) at (13.72,-7.25) {工具与计算资源\\操作与计算作业};
\node[rdexternal] (data) at (13.72,-9.33) {资料及已有系统\\按授权交换};
\draw[rdflow,<->](human.south)--node[rdlabel,right]{意图 / 决策 / 结果}(entry.north);
\draw[rdflow,<->](entry.south)--(coord.north);
\draw[rdflow,<->](coord.south)--(5.82,-4.29);
\draw[draw=primaryblue,line width=.85pt](1.48,-4.29)--(10.16,-4.29);
\foreach \id in {dev,algo,test,ops}{\draw[rdflow,<->](\id.north)--(\id.north |- 0,-4.29);}
\node[rdlabel] at (5.82,-4.29) {任务、结果与专业协作};
\foreach \id in {dev,algo,test,ops}{\draw[rdflow,<->](\id.south)--(\id.south |- 0,-6.45);}
\draw[draw=primaryblue,line width=.85pt](1.48,-6.45)--(10.16,-6.45);
\draw[rdflow,<->](5.82,-6.45)--(runtime.north);
\draw[rdref,<->](runtime.south)--(5.82,-8.17)-|(state.north);
\draw[rdref,<->](5.82,-8.17)-|(art.north);
\node[rdlabel] at (5.82,-8.17) {读写状态与版本引用};
\draw[rdflow,<->](runtime.east)--(tool.west);
\draw[rdflow,<->](runtime.east)++(0,.23)--(11.95,-7.02)--(11.95,-5.22)--(model.west);
\draw[rdref,<->](art.east)--(data.west);
\node[rdlabel,text width=2.45cm] at (13.72,-10.59) {仅提供资源能力\\不承担项目决策};
\node[callnote,font=\scriptsize\sffamily,text width=14.9cm,anchor=north west] at (.22,-11.08)
{实线表示任务交互或能力调用，虚线表示记录与成果访问；双向箭头含请求和返回。组件是逻辑职责，可在一个程序中组合，也可按实际需要分开部署；人和外部服务不在内部自主研发流程之内。};
\end{tikzpicture}
\caption{面向资产管理的研发系统总体架构：协调角色组织版本，专业角色执行，运行程序落实访问与控制。}
\label{fig:rd-overall-design}
\end{figure}
\subsubsection{用户入口与项目工作空间}

统一入口负责接收意图、展示进展和成果、收集澄清与决策。入口可以有不同界面，但应指向同一个项目状态，避免网页、消息窗口与人工笔记分别维护互相矛盾的“最新进度”。用户看到的进度应能进一步展开为任务、阻塞原因和证据引用，不能只展示 Agent 的活动数量。

项目首次建立时，应完成一份简明的项目约定，记录目标、明确排除项、已有资料、术语、角色责任、代码与数据位置、可用资源和授权限制。协调角色据此组织后续版本。新实例进入项目时，从项目约定和当前任务组装必要上下文，不必读取全部历史聊天；无法确认的口径应先澄清，不能把自己的推测写成既定规则。

项目工作空间是逻辑上的统一入口，不要求所有文件存入同一目录或同一服务。本项目以 GitHub 私有仓库维护组织入口、岗位角色、skills 与项目索引；各项目仓库就近维护 requirements、design、tests、程序及相关文档。入口保存明确的仓库、路径和修订引用，便于不同执行环境读取同一套已批准规则。任务状态、对话、真实数据和运行日志不由 Git 维护，工作空间以受控引用连接这些主要位于本地的过程记录。

\subsubsection{项目任务管理与协作控制}

本设计将 Mac 上的平台任务数据库作为项目执行状态的权威来源，将 GitHub 私有仓库作为程序、组织规则与项目文档的权威版本来源，两者通过任务与提交标识关联。它维护版本归属、任务目标、责任角色、执行实例、依赖、状态、接收确认和结果引用。协调 Agent 提出拆分与调整，协作控制程序负责合法状态转换、依赖检查、并发限制和持久化。对话中的“已完成”只有在必要产物与检查条件满足并写入记录后，才构成可追踪的任务完成状态。

任务、消息、计算作业和成果需要各自可定位，不能混用一个会话编号代表全部对象。下表是本项目拟采用的逻辑关系，具体字段和存储产品留待实现设计。

\begin{table}[htbp]
\centering\small\renewcommand{\arraystretch}{1.2}
\begin{tabularx}{\textwidth}{P{1.75cm}P{3.4cm}L}
\toprule
\textbf{对象} & \textbf{主要表达什么} & \textbf{与其他对象的关系}\\
\midrule
版本 & 一次有边界的研发承诺 & 汇集需求、方案、任务及最终交付快照。\\
\rowcolor{lightgray}
任务 & 要完成的目标和责任 & 属于某版本，可拆子任务，可由多个实例先后推进。\\
消息 & 一次请求、通知或反馈 & 引用相关任务、版本或成果；消息送达不等于交接完成。\\
\rowcolor{lightgray}
运行实例 & 某次 Agent 的实际运行 & 记录采用的角色与上下文，可关联任务、消息和工具调用。\\
计算作业 & 提交给执行资源的工作 & 关联发起任务与实例；重试尝试单独记录，保留原作业关系。\\
\rowcolor{lightgray}
成果 & 代码、数据、报告等产出 & 以成果标识与修订定位，并关联来源、任务及相关运行或作业（如有）。\\
\bottomrule
\end{tabularx}
\caption{研发协作中的对象标识及关联关系}
\end{table}

例如，“关系提取实验”是任务；协调者发送的派发请求是消息；算法实例提交的一次批处理是作业；输出数据和实验报告是成果。一个任务可以产生多个作业和成果，失败作业也应保留记录。结果返回消息只携带必要摘要和成果引用，大文件不在各 Agent 的上下文中反复复制。

协作控制还应避免过期状态覆盖新状态。实例报告进展时携带其依据的任务修订，写入时检查该修订是否仍适用；若任务已被取消或已由其他实例接手，应重新读取状态再处理。这里要求的是可检查的一致性，不规定必须建设独立消息中间件或某种数据库集群。

\subsubsection{Agent 运行与模型接入}

运行层先按任务绑定的组织仓库及项目仓库 commit 读取角色、skills、项目约定、需求和方案，再将任务输入、相关成果与权限范围组装为一次执行上下文，启动实例，接收模型提出的行动并调用获准工具。运行记录保留任务关联、关键输入引用、模型与配置标识、工具观察和结束原因，使后续能够理解某项结论如何形成。长期事实应进入项目记录或成果，不能仅依赖一个持续增长的会话。

模型接入按能力需要选择服务和配置，为各角色提供统一的调用边界。协调可能需要任务规划与冲突解释，开发需要代码处理，测试需要结合规格分析失败证据；这些差异可以通过上下文、工具和检查要求体现，不必为每个岗位绑定不同模型。关键实验仍须记录实际使用的模型版本或可获得的服务标识，便于解释结果变化。

运行层记录单次 Agent 执行的预算、超时和暂停条件；持久调度程序负责跨运行的队列、租约、重试与恢复，不能由模型自行修改任务状态。模型服务不可用时，系统应保存已经完成的结果、当前等待原因和可恢复位置。切换模型若可能改变实验可比性或验收条件，须作为过程变化记录，而不是无痕替换。权限检查始终发生在实际工具和资源访问处，不能交给模型自行承诺遵守。

\subsubsection{工具、计算与研发资源}

工具层把资料读取、代码操作、文件处理、测试执行和作业提交等能力提供给 Agent。工具输入应清楚说明作用对象和范围，返回应区分完成、失败、等待及状态未知，并提供可核验的实际结果。读取资料与修改代码、提交作业与删除数据具有不同风险，不能因它们都以“工具调用”出现就授予同样权限。

执行能力既可由普通程序提供，也可由具有自身推理循环的开发 Agent 提供，例如通过适配器接入 Codex；工具入口并不意味着内部只是一条固定命令。上层按任务契约交接，不依赖具体 SDK 会话对象。

长计算由作业接口提交到合适的执行资源，返回作业标识，后续再查询状态与结果。Agent 等待计算期间可以释放推理运行资源，但任务仍保持未完成。计算所需的环境、依赖、输入快照、参数与输出位置应随作业留存；出现故障时先查原作业，确认能否安全重试，再决定是否另起一次尝试。

资源安排由运维角色维护，运行程序执行并发、配额和隔离限制。例如两个算法实例可以并行比较方法，但必须写入不同输出位置，并使用同一份冻结输入；开发和测试可以共享计算设施，但测试不得无记录地使用开发工作目录中随时变化的文件。本项目由 Mac 的可信 Worker 执行控制与轻量任务，4090 承担重计算；dot 云电脑作为独立轻量研发环境，阿里云作为一般 Web 服务与轻量接口验证环境。后两者不因参与研发就自动成为自建队列的 Worker，接入方式与数据许可须分别落实。

\subsubsection{数据、代码与成果管理}

本设计设置逻辑统一的成果登记与引用规则，作为“交付了什么”的权威依据，与任务记录共同构成项目主线。成果正文可以分散存放，但登记信息应包含类型、位置、修订、责任任务、来源与检查状态；必要时记录内容校验值，以发现文件被替换的情况。资料或数据本身的可信度还要根据来源判断，登记成功并不代表内容真实。

版本交付采用固定快照：列出所依据的需求修订、代码提交、数据快照、实验参数、模型标识、测试报告及运行说明。引用不能仅写“最新代码”或“共享目录里的结果”，因为这些内容会变化。候选成果可以产生新修订，但已验收快照不得被悄悄替换；修订和重新验收应留下对应关系。

程序、组织规范、角色、skills 及项目文档统一以 GitHub 私有仓库中的审阅后版本为准。代码和文档修改通过分支、差异审阅、必要测试与合并进入基线，任务绑定实际采用的 commit；不能直接读取不断变化的分支名作为已冻结输入。原始资料、处理结果和人工修订分别记录来源。知识图谱中的关系也应保留到原文位置的追踪路径。跨版本复用成果时，复用其明确修订并检查前提是否仍成立：前一版本的否定实验可以减少重复试错，却不能自动证明新数据和新方法仍会失败。

数据存放按“版本内容、过程数据、节点运行状态”分开。代码与获准文档进入 GitHub，密钥、敏感账户资料、真实业务数据、任务数据库和过程日志不进入 Git。Mac 保留主要状态记录和成果归档，4090 保留当前数据与实验记录；dot 和阿里云等节点按任务需要保留工作副本、临时文件及必要运行日志和状态，交接后将需保留的证据与归档回收至本地受控位置。备份、访问权限、保留期限和清理责任按类别配置，不能把“过程数据主要本地保存”误写为“云端没有数据”。

\subsubsection{总体架构与外部接口边界}

回看本节开头的图~\ref{fig:rd-overall-design}，上述组成通过任务、运行与成果引用连接起来。人通过入口与协调角色交互，协调角色把版本工作分派给各专业角色；协调及各专业 Agent 均由运行层承载，调用模型、工具与计算资源，再把结果写入任务和成果记录。角色间可以在授权范围内直接讨论接口、缺陷和专业问题，但重要结论及状态变化应回到共享记录，不能只存在于一条临时消息中。

外部接口需要区分“提出建议”和“执行操作”。调用模型返回的是候选判断或行动请求，调用工具才可能改变代码、文件和环境；运行层应在这个边界核对授权与任务范围。向外部模型或服务发送上下文前，还需按资料使用条件裁剪内容。外部系统返回的状态也应由对应接口核实，不能由模型生成文字代替。

图中的外部资源位于 Agent 协作应用的逻辑边界之外，可以由自有本机或远程节点提供，也可以来自第三方；“外部”不等于公有云或外部供应商。总体架构不要求为每个框建立一个微服务。起步时可用一个应用承载入口、协调控制和运行管理，配合已有代码与文件存储；需要独立扩容、故障隔离或更严格访问控制时，再拆分部署。逻辑边界在部署变化前后保持清楚，才能知道迁移影响的是执行位置、数据访问还是业务责任。

\Needspace{12cm}
\subsection{端到端研发运行流程}

图~\ref{fig:rd-end-to-end} 先串起从人的意图到交付反馈的完整流程，再说明各阶段的输入、输出与交接条件。开发、算法、测试和运维按依赖并行准备，候选成果须经过独立验证。

\begin{figure}[H]
\centering
\begin{tikzpicture}[x=1cm,y=.62cm,font=\footnotesize\sffamily,text=darkgray,
  lane/.style={draw=primaryblue!16,fill=primaryblue!2},
  phasebox/.style={callstep,font=\scriptsize\sffamily,text width=2.35cm,minimum height=.58cm,inner sep=3pt},
  phase/.style={font=\scriptsize\sffamily,text=darkgray,align=left},
  reflow/.style={-{Stealth[length=1.7mm]},draw=softred,dashed,line width=.8pt}]
\path[use as bounding box](0,-14.45)rectangle(15.4,.38);
\foreach \x/\name in {1.54/人类负责人,4.62/研发协调,7.70/开发与算法,10.78/独立测试,13.86/运维}{
 \path[lane](\x-1.47,.25)rectangle(\x+1.47,-12.90);
 \node[callhead,font=\footnotesize\bfseries\sffamily,text width=2.5cm,minimum height=.48cm,inner sep=3pt] at (\x,-.10) {\name};
}
\node[phasebox] (intent) at (1.54,-1.38) {提出意图\\确认目标约束};
\node[phasebox] (clarify) at (4.62,-1.38) {读取项目约定\\澄清需求};
\draw[callrequest,<->](intent.east)--(clarify.west);
\node[phasebox,fill=primaryblue!7] (vplan) at (4.62,-2.94) {版本需求包\\方案与任务依赖};
\draw[callrequest](clarify.south)--(vplan.north);
\node[phasebox] (design) at (7.70,-2.94) {技术方案\\研究与实验计划};
\node[phasebox,fill=softgreen!7] (testplan) at (10.78,-2.94) {验收计划\\独立检查条件};
\node[phasebox] (envplan) at (13.86,-2.94) {资源与访问条件\\运行交接要求};
\draw[callrequest,<->](vplan.east)--(design.west);
\draw[callrequest,<->](design.east)--(testplan.west);
\draw[callrequest,<->](testplan.east)--(envplan.west);
\node[phase,anchor=east] at (3.02,-2.94) {重要取舍\\提交人决定};
\draw[callreturn](vplan.west)--(3.10,-2.94)--(3.10,-2.02)--(1.54,-2.02)--(intent.south);
\node[phasebox,fill=primaryblue!7] (freeze) at (4.62,-4.63) {冻结版本需求\\派发任务};
\draw[callrequest](vplan.south)--(freeze.north);
\node[phasebox] (execute) at (7.70,-4.63) {实现与实验\\分支隔离、自测};
\node[phasebox] (testprep) at (10.78,-4.63) {准备测试样本\\验证程序与步骤};
\node[phasebox] (envprep) at (13.86,-4.63) {准备研发环境\\观察计算作业};
\draw[callrequest](freeze.east)--(execute.west);
\draw[callrequest](freeze.south)--(4.62,-5.48)--(13.86,-5.48)--(envprep.south);
\draw[callrequest](10.78,-5.48)--(testprep.south);
\node[calllabel,font=\scriptsize\sffamily] at (10.10,-5.48) {按依赖并行准备};
\node[phasebox] (candidate) at (7.70,-6.53) {集成候选修订\\自测与成果引用};
\draw[callrequest](execute.south)--(candidate.north);
\node[phasebox,fill=softgreen!7] (verify) at (10.78,-6.53) {固定输入验证\\程序检查与复核};
\draw[callrequest](candidate.east)--(verify.west);
\draw[callrequest](testprep.south)++(0,-.40)--(verify.north);
\node[phasebox] (evidence) at (10.78,-8.24) {验收意见\\结果与失败证据};
\draw[callrequest](verify.south)--(evidence.north);
\node[phasebox,fill=primaryblue!7] (check) at (4.62,-8.24) {核对验收证据\\安排返工或交付};
\draw[callreturn](evidence.west)--node[calllabel,font=\scriptsize\sffamily,above]{验收证据}(check.east);
\draw[reflow](check.north)--(4.62,-7.26)--(6.16,-7.26)--(6.16,-4.63)--(execute.west);
\node[calllabel,font=\scriptsize\sffamily,text=softred,rotate=90] at (6.16,-6.08) {有界返工};
\node[phasebox] (accept) at (1.54,-9.94) {审阅成果限制\\确认关键事项};
\node[phasebox,fill=primaryblue!7] (deliver) at (4.62,-9.94) {固定交付快照\\发起交接确认};
\draw[callrequest](check.south)--(deliver.north);
\draw[callrequest,<->](deliver.west)--(accept.east);
\node[phasebox] (handover) at (13.86,-9.94) {接收成果包\\确认运行交接};
\draw[callrequest,<->](deliver.east)--node[calllabel,font=\scriptsize\sffamily,above]{成果包 / 接收确认}(handover.west);
\node[calllabel,font=\scriptsize\sffamily] at (9.23,-10.66) {协调汇总接收确认后关闭本版};
\node[phasebox] (feedback) at (13.86,-11.76) {运行问题反馈\\日志与复现材料};
\draw[callrequest](handover.south)--(feedback.north);
\node[phasebox] (next) at (4.62,-11.76) {下一版本候选\\复用成果证据};
\draw[callreturn](feedback.west)--node[calllabel,font=\scriptsize\sffamily,above]{缺陷、改进与新问题}(next.east);
\node[phasebox] (priority) at (1.54,-11.76) {决定后续优先级\\进入下一版本};
\draw[callrequest](next.west)--(priority.east);
\node[callnote,font=\scriptsize\sffamily,text width=14.8cm,anchor=north west] at (.28,-13.24)
{时间总体向下。测试提前参与规划，候选成果提交后再独立验收；各实施角色持续自测。红色虚线为保持既定目标的返工，新需求进入下一版本；涉及重要取舍的变化回到人类负责人。};
\end{tikzpicture}
\caption{端到端研发流程泳道图：从意图、定版和执行到独立验证、交付与反馈。}
\label{fig:rd-end-to-end}
\end{figure}
\subsubsection{意图接收与需求澄清}

流程从人提交意图开始。协调角色先读取项目约定与已有成果，确认这是新项目、新版本还是既有任务的补充，再澄清目标对象、使用场景、范围、时间与资源约束。能从项目记录确认的问题不重复询问；会改变目标或验收口径的问题不能自行假设。

本阶段输出可评审的需求草案，说明要解决什么问题、预期交付什么、哪些事项暂不涉及，以及仍需确认的问题。例如“做一个知识图谱”需要进一步明确资料集合、实体和关系类型、查询方式与证据定位要求，才能成为可执行需求。尚未确定的目标可以先安排有边界的探索任务，其结果用于后续定版，不应伪装成已经确认的交付承诺。

\subsubsection{版本需求包与方案确定}

协调角色把经确认的需求组织为版本需求包，包括目标、范围、交付清单、验收条件、依赖、资源边界和关键决策。开发负责人组织技术方案，算法角色补充研究假设与实验计划，测试角色同步提出验收计划和数据准备要求，运维角色检查资源、访问条件和运行交接要求。存在分歧时，先明确分歧会影响哪项需求及何种证据，再由相应责任方决定。

方案确定后形成版本基线，冻结本版本需求，并在项目仓库中固定 requirements、design、tests 的审阅版本；派发任务同时记录组织仓库与项目仓库 commit、输入、责任和验收依据。冻结后不再向本版本加入新需求；新想法进入后续版本候选。为实现既定需求而调整规格时，应记录原规格、调整原因、影响任务、决策与新的修订，通知相关角色；若实际改变目标、交付边界或关键取舍，应重新提交人决定，不能用“规格调整”掩盖需求扩张。

\subsubsection{任务执行、实验与过程调整}

各角色接收任务后核对输入与前置条件，确认接收，再开展实现、实验或准备工作。开发和算法并行的前提是接口、样例和共享输入已经明确；测试可以同时准备独立检查，运维可以准备计算环境。协调角色跟踪必要依赖，遇到阻塞时先判断是输入缺失、技术问题还是范围问题，再安排补充、返工或决策。

执行过程中，实例把阶段性成果和证据登记到任务，而不是只报告百分比。实验失败应记录失败条件、已排除的原因和后续建议；重复失败达到预定停止条件时，应形成结论或提交取舍，不无限重试。跨版本研究可以继续积累证据，本版本只对约定的实验范围与交付负责，不承诺必须发现有效方法。

过程调整产生新的规格或方案修订时，受影响任务需确认采用的修订。已启动的作业可能仍依据旧输入运行，系统应决定保留为对照、停止或重跑，并记录理由。候选成果提交前，开发组负责人组织代码与能力集成，各责任角色完成相应自测，再将集成后的候选修订、运行条件、已知限制和结果引用一起交给接收方。

\subsubsection{独立验证与成果交付}

独立测试接收固定的需求修订与候选成果，核对交付是否齐全，再按计划执行验证。软件任务检查功能与异常路径，数据任务检查质量约束与来源，算法任务检查实验方法、复现性与结论边界。自动测试程序产生具体观察，测试 Agent 整理证据并给出通过、不通过或证据不足的意见；测试不能用实现 Agent 的自述代替验证。

未通过时，缺陷应指向需求、候选版本、复现步骤与失败证据，由协调角色确定责任与返工安排。新候选修订提交后，对受影响部分重新测试，并根据影响范围决定是否回归检查。研究得出否定结论时，只要完成了约定实验且结论有依据，可以通过研究任务验收；若目标是交付某项可用能力，则还须单独说明该能力是否已经具备。

协调角色汇总通过情况、未解决问题和使用限制，组织最终交付确认。交付包固定代码、文档、数据、报告和验证记录的版本，并包含操作说明及运行条件；GitHub 中的审阅提交与本地成果清单相互引用，部署目标只使用验收确认的提交或发布包。运维确认能够按说明准备或接收环境后，才算完成相应交接。研发验收通过不等于取得投资、交易或其他生产业务的使用授权。

本节开头的图~\ref{fig:rd-end-to-end} 已展示各角色在上述过程中的交接关系；交付后的运行反馈继续形成下一版本的输入。

\subsubsection{运行反馈与下一版本迭代}

交付后，运维根据实际运行观察提交问题，开发和算法根据使用反馈提出改进，测试补充遗漏的场景。每项反馈应尽量携带受影响版本、发生条件、日志或数据样例，协调角色区分既有缺陷、使用问题与新增需求，再安排进入相应的处理队列。急需修复的问题也要建立有范围的修订任务，不能直接覆盖已交付快照。

下一版本规划复用已有资料、代码和实验，但重新确认其适用前提。研究可能跨多个版本逐步推进：本版证实某一方法受限，下版可以更换数据或方法继续研究；已有否定结论连同限制一并保留。版本关闭表示本版承诺及遗留事项已经交代清楚，不表示该研究方向永久结束。

\subsubsection{知识图谱试点的流程示例}

具体项目另立需求文档与方案设计文档。以下仅以假设性的知识图谱试点说明平台交接流程，不代替项目文档，也不代表已确认的指标或完成结果。试点问题设为：能否从一组经确认的项目资料中，抽取有限类型的实体与关系，并从每条关系回到原文证据。示例只包含“项目、公链、赛道”三类实体，以及“项目部署于公链”“项目属于赛道”两类关系；不扩展到实时数据、投资建议或交易执行。

\begin{enumerate}[leftmargin=2em,itemsep=4pt]
  \item \textbf{澄清与定版。}人确认资料集合及预期查询方式。协调角色整理资料是否允许处理、哪些定义仍有歧义、交付原型还是方法报告。测试提出验收需有独立核对的样本及错误分类，最终评估样本与方法调试样本区分，并保留独立核对依据；抽取质量、证据完整性和运行成本的具体标准在需求包中商定，不在此预设数值。
  \item \textbf{方案与任务。}算法角色设计实体规范化和关系判定方法，开发角色设计资料解析、结果存取与证据查询接口，运维核对环境及资料访问条件，测试准备含同名项目、多链部署、赛道口径不一致和缺少证据等情形的检查样本。协调角色据依赖安排任务并冻结本版范围。
  \item \textbf{执行与成果登记。}原始资料固定快照，解析程序为片段保留文档标识和位置。算法实例在固定输入上运行实验，作业输出关系候选与原文引用，开发实例完成查询原型。无法确认的关系标记为待核，不用推断补齐；不同方法的结果分别登记，避免覆盖。
  \item \textbf{独立验收。}验证程序检查实体与关系类型、引用是否可定位、必填信息和重复记录；测试 Agent 对照独立样本复核关系是否由原文支持，并检查模糊案例是否被正确处理。发现错误后提交对应候选修订、样例和证据，由责任角色修复或说明方法限制，再重新验证。
  \item \textbf{交付与后续判断。}交付包包含原型代码、数据与模型配置引用、图谱样本、来源链路、测试报告和已知限制。若实验表明当前方法无法满足约定条件，则交付可复现的否定结论，明确尚未形成可用能力；由人决定停止、缩小范围或把替代方法列入下一版本。
\end{enumerate}

这个试点把系统各部分连成一个可检查的过程：人决定问题与取舍，协调角色维护版本，专业角色产出成果，程序执行真实检查，独立测试形成验收证据。后续可以据此验证系统本身是否能够可靠交接、定位来源和恢复状态，再决定扩展到更复杂的研发任务。


% Chapter 5, sections 5.1--5.2. Authoring fragment; no installation performed.
% Chapter 5, sections 5.1--5.2. Authoring fragment; no installation performed.
\section{研发平台实施与配置：OpenAI 主方案}

本章把第四章的逻辑设计落实为一套可分步建设、由用户长期掌握的研发平台。\textbf{长期保留的是平台核心：项目与版本、任务与消息、权限与审批、运行与成果、调度与恢复；OpenAI 是首期模型与执行能力的主方案。}Agents SDK、Codex 执行器及模型服务通过适配层接入，可以随能力、成本与部署条件替换，不能承担平台唯一的业务状态。GitHub 私有仓库保存代码与文档的权威版本，PostgreSQL 保存协作与运行状态，受控本地存储保存过程数据及大产物。

Mac 承担主要开发与首期平台控制，4090 承担训练和重计算；T640 在网络及适配验收后作为可选推理资源。现有阿里云服务器用于一般 Web 与轻量接口，dot 云电脑作为独立可用的文档、简单 Web 和轻量研发环境。下列软件、配置和命令是\textbf{待实施方案}，不表示已经安装、连通或验收；自研模块须先完成开发，才能执行相应启动命令。本文也不构成安装软件、配置凭据、修改网络或向 GitHub 上传内容的操作授权。

\Needspace{6\baselineskip}
\subsection{实施基线与总体部署}

\subsubsection{平台建设范围与交付清单}

首期建设的是可承接不同研发项目的公共平台，以小型数据任务和一般 Web 服务验证研发闭环，不预先建设某个知识图谱业务系统，也不接入实盘交易。阿里云已有 Web 服务与交易所网格程序属于现有业务，新增验证不改写其配置、账户、订单或运行流程。固定参数网格的后续整理及 RSI 简单策略模拟分别另立项目需求与方案，不能据此扩张本章范围。验收应以一条完整链路为对象：人提交需求，协调角色形成任务，专业角色调用受控执行能力，结果登记到任务及成果记录，独立测试核对证据，最后由入口呈现交付包和待决事项。

平台交付物分为五组：（1）可安装的 Python 项目、锁定依赖及执行器版本记录；（2）项目、任务、消息、运行、作业、审批和成果的数据结构及迁移脚本；（3）Web 入口、控制 API、协调循环、Agents SDK／Codex 适配和 Worker 程序；（4）组织入口、角色、skills、项目索引及模型、节点、权限和运行参数配置；（5）安装说明、自动化测试、备份恢复说明及一次端到端验收记录。交付代码不等于平台验收完成，必须同时证明状态可追踪、失败可识别、操作可约束；替换某一执行器时，应能保留原任务、成果与审计记录。

\subsubsection{OpenAI主方案与软件选型}

首期采用\textbf{自建平台核心 + Agents SDK 协调与专业 Agent + Codex 工程执行器 + 确定性计算 Worker}。平台核心掌握任务及消息路由、权限、调度、审批与交付；Agents SDK 提供一次 Agent 运行中的模型与工具循环，适合协调、研究、算法、测试等角色；Codex 用于受控工作区内的代码阅读、修改、命令执行和工程验证；数据处理、指标计算及训练则交给明确输入输出的程序。专业角色按任务选择执行能力，不要求每个角色占一台机器，也不要求所有步骤经过两个 SDK。

OpenAI Agents SDK 的 Python 包为 \texttt{openai-agents}，导入名为 \texttt{agents}；本方案以 Responses API 为通用模型接口，并由平台适配其结果与事件。SDK 的对话、handoff 与追踪有助于一次协作，但不替代持久任务库、审批账本或独立验收。控制服务无需另建厂商托管中枢，平台协议也不应暴露供应商特有对象作为唯一业务主键。\cite{rdRuntimeSdk}

官方 Codex SDK 文档已提供稳定版 Python 包 \texttt{openai-codex}，导入名为 \texttt{openai\_codex}；它通过 JSON-RPC 控制本地 Codex app-server，发布构建带固定的 Codex CLI runtime 依赖。实施时锁定 SDK 与 runtime 的兼容组合，不凭桌面应用版本推断后台执行版本，不默认使用预发布包或任意替换系统 \texttt{codex}。直接操作 app-server 协议的适配另行验证；首期不把实验性的远程 WebSocket 入口作为平台控制接口。\cite{codex-sdk,codex-app}

\begingroup\small\renewcommand{\arraystretch}{1.16}
\begin{longtable}{P{2.0cm}P{3.75cm}P{9.0cm}}
\toprule
\textbf{层次} & \textbf{首期选择} & \textbf{采用理由与自研边界}\\
\midrule
\endfirsthead
\toprule
\textbf{层次} & \textbf{首期选择} & \textbf{采用理由与自研边界（续）}\\
\midrule
\endhead
\bottomrule
\endfoot
\endlastfoot
平台核心 & 自研 \texttt{rd\_platform} & 长期掌握项目、任务、权限、审批、消息、成果和恢复；用内部契约接入可替换的模型与执行器。\\
\rowcolor{lightgray}
运行与依赖 & Python 3.12、uv & 控制程序统一语言；uv 管理虚拟环境与锁文件，不使用系统 Python 安装项目依赖。\\
模型与 Agent & \texttt{openai-agents}、\texttt{openai} & 复用通用 Agent 运行与 API 客户端；自研角色配置、上下文装配、工具封装和业务记录适配。\\
\rowcolor{lightgray}
工程执行 & \texttt{openai-codex} 及其固定 runtime & 在受限工作区执行工程任务；自研 Codex 适配、审批桥接、事件与成果登记，退出或报告完成不直接等于平台验收。\\
交互入口 & FastAPI、Uvicorn、Jinja2、HTMX & 服务端页面配合局部刷新；任务页、审批页、权限与会话保护由平台实现。Mac 桌面应用和 VS Code 供人交互。\\
\rowcolor{lightgray}
任务与队列 & PostgreSQL 17 & 同一事务库保存任务、消息、租约及事件；首期不增设 Redis、Celery 或独立消息中间件。领取、幂等与恢复逻辑自行实现。\\
数据访问 & SQLAlchemy、psycopg、Alembic & 分别承担数据访问、PostgreSQL 驱动及模式迁移；业务表与状态转换规则自行编写。\\
\rowcolor{lightgray}
计算与节点 & 自研 Python Worker & 领取受控任务，调用适配器或确定性程序并上报证据；不把 FastAPI 请求进程当作长计算执行器。\\*
代码与文档 & Git、GitHub 私有仓库、本地目录 & GitHub 是程序与文档的权威版本核心；组织入口、角色、skills、项目索引及项目需求、方案、测试随仓库审阅。过程数据与秘密不入 Git。\\*
\rowcolor{lightgray}
测试与部署 & pytest、Ruff、Docker Desktop / Compose & 自动测试与静态检查；Mac 数据库使用容器，控制程序先以原生 Python 进程运行，任务按边界隔离。\\
\bottomrule
\caption{首期软件清单：保留平台核心，复用并隔离可替换的执行能力。}
\end{longtable}
\endgroup

\textbf{人工入口与程序接入分别验收。}Mac 上集成的 ChatGPT／Codex 桌面应用及 VS Code 是人提出任务、查看修改与审阅结果的入口；程序化 Codex 由平台 Worker 调用 SDK／runtime。桌面可用、订阅可用或 IDE 已登录，均不证明后台进程的认证、模型权限、计费方式与工作区权限已经可用。Agents SDK 的 API 凭据和 Codex runtime 的认证按实际官方支持方式分别配置、分别测试，不能复制桌面会话秘密来替代接入设计。\cite{v121-codex-auth}

模型名称放入 \texttt{models.yaml}，经账号可用性、效果、费用和工具兼容性验证后确定，不随 SDK 默认值漂移。训练框架和 CUDA 属于计算节点环境，不安装进控制平台虚拟环境；具体研发项目也应有各自的依赖锁文件。dot 托管环境独立提供轻量研发能力，不是本平台 SDK 调用链中的默认 Worker。

\Needspace{22cm}
\subsubsection{协作环境、计算节点与交互概念图}

图~\ref{fig:rd-implementation-deploy} 同时表达\textbf{人的入口、长期保留的平台核心、按任务运行的多个 Agent、可替换执行器及设备落点}。Mac 上可运行多个协调／专业实例，它们的职责、进程和机器不是一一对应关系；4090 是计算资源而非“算法角色本身”。图中连线是拟建交互，不要求节点全互连，也不表示接口已经实现。

\begin{figure}[H]
\centering
\begin{tikzpicture}[
  x=1cm,y=1cm,
  every node/.style={font=\footnotesize\sffamily,text=darkgray,align=center,inner sep=2pt},
  concepttitle/.style={font=\small\bfseries\sffamily},
  conceptbody/.style={font=\fontsize{8}{10.6}\selectfont\sffamily},
  conceptlabel/.style={font=\fontsize{8}{10.5}\selectfont\sffamily,fill=white,inner sep=2.5pt},
  conceptline/.style={line width=1pt,line cap=round,line join=round},
  conceptflow/.style={{Stealth[length=1.7mm]}-{Stealth[length=1.7mm]},draw=primaryblue,line width=.9pt},
  conceptpublish/.style={-{Stealth[length=1.7mm]},draw=primaryblue,line width=.9pt},
  conceptoption/.style={{Stealth[length=1.7mm]}-{Stealth[length=1.7mm]},draw=softpurple,dashed,line width=.9pt},
  conceptpending/.style={draw=darkgray!60,densely dotted,line width=1pt}
]
\path[use as bounding box] (0,0) rectangle (15.5,17.6);

% Mac is a deployment boundary. Its inner roles are logical roles, not machines.
\fill[primaryblue!3,rounded corners=10pt] (3.2,5.3) rectangle (12.3,13.6);
\draw[primaryblue!30,rounded corners=10pt,line width=.7pt] (3.2,5.3) rectangle (12.3,13.6);
\begin{scope}[shift={(7.75,13.1)},conceptline,draw=primaryblue]
  \draw[rounded corners=1pt,fill=white] (-.4,-.12) rectangle (.4,.26);
  \draw[fill=primaryblue!9] (-.4,-.12)--(-.55,-.23)--(.55,-.23)--(.4,-.12)--cycle;
\end{scope}
\node[concepttitle] at (7.75,12.61) {MacBook Pro M3 Max：首期主控与开发};
\node[conceptbody] at (7.75,12.28) {64 GB / 2 TB；进程、工作区与凭据按任务隔离};

% Human and human-facing UI.
\begin{scope}[shift={(7.75,17.1)},conceptline,draw=accentorange!85!black]
  \fill[accentorange!12] (0,0) circle (.18);
  \draw (0,0) circle (.18);
  \draw[fill=accentorange!9] (-.34,-.53) .. controls (-.34,-.24) and (.34,-.24) .. (.34,-.53) -- cycle;
\end{scope}
\node[concepttitle] at (7.75,16.3) {人：目标、约束、授权与取舍};
\draw[conceptflow] (7.75,16.03) -- (7.75,13.64);
\node[conceptlabel] at (7.75,15.12)
 {人工入口：Mac ChatGPT／Codex、VS Code\\平台工作台：任务、审批与交付\\桌面可用不等于后台程序接入已验收};

% GitHub folder and version branch.
\begin{scope}[shift={(1.65,15.78)},conceptline,draw=primaryblue]
  \path[fill=primaryblue!7] (-.56,-.38)--(-.56,.25)--(-.12,.25)--(.02,.38)--(.56,.38)--(.56,-.38)--cycle;
  \draw (-.56,-.38)--(-.56,.25)--(-.12,.25)--(.02,.38)--(.56,.38)--(.56,-.38)--cycle;
  \draw (-.2,-.14)--(-.2,.08) (-.2,-.01)..controls(-.2,.08)and(.2,-.06)..(.2,.1);
  \foreach \xx/\yy in {-.2/-.2,-.2/.15,.2/.17}{\fill[white](\xx,\yy)circle(.055);\draw(\xx,\yy)circle(.055);}
\end{scope}
\node[concepttitle] at (1.65,15.14) {GitHub 私有仓库};
\node[conceptbody,anchor=north] at (1.65,14.8)
 {代码／规则／项目文档\\固定 commit；审阅归档};
\draw[conceptflow] (1.65,14.04)--(1.65,13.05)--(3.16,13.05);
\node[conceptlabel] at (1.65,12.55) {代码／文档版本交接};

% External model service, separately from dot and local execution.
\begin{scope}[shift={(13.86,15.78)},conceptline,draw=softpurple]
  \fill[softpurple!5] (0,0) circle (.48);
  \draw (-.33,.14)--(0,.31)--(.33,.14)--(.24,-.25)--(-.24,-.25)--cycle;
  \draw (-.33,.14)--(.24,-.25) (.33,.14)--(-.24,-.25) (0,.31)--(0,-.07);
  \foreach \xx/\yy in {-.33/.14,0/.31,.33/.14,.24/-.25,-.24/-.25,0/-.07}{\fill[white](\xx,\yy)circle(.06);\draw(\xx,\yy)circle(.06);}
\end{scope}
\node[concepttitle] at (13.86,15.14) {OpenAI 模型服务};
\node[conceptbody,anchor=north] at (13.86,14.8)
 {API／Codex 分开认证\\仅发获准的必要上下文};
\draw[conceptflow] (13.86,14.02)--(13.86,11.5)--(12.34,11.5);
\node[conceptlabel] at (13.86,12.57) {受控模型调用\\请求与响应};

% Durable, user-owned core.
\draw[primaryblue!55,fill=white,rounded corners=4pt] (3.67,10.89) rectangle (11.83,12.07);
\node[font=\small\bfseries\sffamily,text=primaryblue] at (7.75,11.84) {长期自有的平台核心};
\node[conceptbody] at (7.75,11.27)
 {项目／任务／消息／权限／审批／成果／恢复\\PostgreSQL 记状态；本地受控目录存过程数据};
\draw[conceptflow] (7.75,10.89)--(7.75,10.51);

% Multiple agent icons and the coordinator/professional hierarchy.
\draw[softblue!40,fill=white,rounded corners=4pt] (3.67,7.92) rectangle (11.83,10.51);
\node[anchor=west,font=\fontsize{8}{10}\selectfont\bfseries\sffamily,text=primaryblue] at (3.85,10.27)
 {Agents SDK 适配（可替换）};
\begin{scope}[shift={(7.75,10.01)},conceptline,draw=primaryblue]
  \draw[fill=primaryblue!9] (0,.04) circle (.11);
  \draw (-.2,-.28)..controls(-.2,-.09)and(.2,-.09)..(.2,-.28);
\end{scope}
\node[conceptbody,anchor=west] at (8.02,9.77) {协调 Agent};
\draw[conceptpublish] (7.75,9.56)--(7.75,9.33)--(4.64,9.33)--(4.64,9.1);
\draw[conceptpublish] (7.75,9.56)--(7.75,9.33)--(6.71,9.33)--(6.71,9.1);
\draw[conceptpublish] (7.75,9.56)--(7.75,9.33)--(8.78,9.33)--(8.78,9.1);
\draw[conceptpublish] (7.75,9.56)--(7.75,9.33)--(10.85,9.33)--(10.85,9.1);
\foreach \xx/\role in {4.64/开发,6.71/算法,8.78/测试,10.85/运维}{
  \begin{scope}[shift={(\xx,8.85)},conceptline,draw=softblue!85!black]
    \draw[fill=softblue!9] (0,.02) circle (.1);
    \draw (-.18,-.28)..controls(-.18,-.08)and(.18,-.08)..(.18,-.28);
  \end{scope}
  \node[conceptbody] at (\xx,8.29) {\role{} Agent};
}

% Codex and deterministic execution are alternative tools under platform control.
\draw[conceptflow] (5.55,7.92)--(5.55,7.35);
\draw[conceptflow] (9.95,7.92)--(9.95,7.35);
\draw[conceptflow] (11.83,11.37)--(12.07,11.37)--(12.07,7.54)--(9.95,7.54);
\draw[primaryblue!55,line width=.65pt] (9.95,7.54)--(5.55,7.54);
\node[conceptlabel,fill=primaryblue!3] at (7.75,7.64) {选择执行／回传证据};
\draw[softpurple!45,fill=white,rounded corners=4pt] (3.67,5.49) rectangle (7.47,7.35);
\begin{scope}[shift={(5.57,7.02)},conceptline,draw=softpurple]
  \draw[rounded corners=1pt] (-.35,-.19) rectangle (.35,.2);
  \draw (-.23,.06)--(-.11,-.01)--(-.23,-.09) (.02,-.09)--(.2,-.09);
\end{scope}
\node[font=\small\bfseries\sffamily] at (5.57,6.52) {Codex 工程执行器};
\node[conceptbody] at (5.57,5.94) {SDK／本地 runtime（可替换）\\受限读写、命令、diff 与测试};
\draw[softgreen!45,fill=white,rounded corners=4pt] (8.03,5.49) rectangle (11.83,7.35);
\begin{scope}[shift={(9.93,7.02)},conceptline,draw=softgreen!75!black]
  \draw[fill=softgreen!6,rounded corners=1pt] (-.33,-.2) rectangle (.33,.2);
  \foreach \yy in {-.11,0,.11}{\draw(-.21,\yy)--(.06,\yy);\fill(.21,\yy)circle(.022);}
\end{scope}
\node[font=\small\bfseries\sffamily] at (9.93,6.52) {确定性工具与计算};
\node[conceptbody] at (9.93,5.94) {本机测试／数据检查\\Worker 派发训练与重计算};

% dot remains independently useful, with no automatic SDK-worker promise.
\begin{scope}[shift={(1.65,10.13)},conceptline,draw=softblue!85!black]
  \path[fill=softblue!8] (-.64,-.2)..controls(-.86,-.2)and(-.87,.12)..(-.64,.17)
    ..controls(-.64,.5)and(-.19,.6)..(0,.35)..controls(.29,.47)and(.57,.3)..(.56,.08)
    ..controls(.85,.06)and(.83,-.22)..(.6,-.22)--cycle;
  \draw (-.64,-.2)..controls(-.86,-.2)and(-.87,.12)..(-.64,.17)
    ..controls(-.64,.5)and(-.19,.6)..(0,.35)..controls(.29,.47)and(.57,.3)..(.56,.08)
    ..controls(.85,.06)and(.83,-.22)..(.6,-.22);
\end{scope}
\node[concepttitle] at (1.65,9.53) {dot 云电脑};
\node[conceptbody,anchor=north] at (1.65,9.17)
 {独立可用的文档\\简单 Web 与轻量研发\\经审阅交接固定版本\\非自动 SDK Worker};
\draw[conceptpending] (2.48,10.13)--(2.65,10.13) (2.94,10.13)--(3.16,10.13);
\draw[darkgray!65,line width=.9pt] (2.73,10.02)--(2.8,10.24) (2.82,10.02)--(2.89,10.24);
\node[conceptlabel] at (1.72,11.19) {自动任务接口待另行设计};

% Device endpoints. Arrows cross the Mac boundary only for their named duties.
\draw[conceptpublish] (3.2,6.48)--(2.05,6.48)--(2.05,4.02);
\node[conceptlabel] at (1.7,5.3) {经授权发布\\Web／测试接口};
\draw[conceptflow] (9.93,5.49)--(9.93,4.73)--(7.75,4.73)--(7.75,4.02);
\node[conceptlabel] at (7.35,4.73) {任务／状态／结果引用};
\draw[conceptoption] (12.3,8.37)--(13.86,8.37)--(13.86,4.02);
\node[conceptlabel] at (13.86,6.48) {网络修复及适配后\\按需接入推理};

\begin{scope}[shift={(2.05,3.46)},conceptline,draw=accentorange!85!black]
  \draw[fill=accentorange!7,rounded corners=2pt] (-.5,-.37) rectangle (.5,.37);
  \draw (-.5,.17)--(.5,.17);
  \fill (-.35,.26) circle (.022) (-.24,.26) circle (.022) (-.13,.26) circle (.022);
  \draw (0,-.1) circle (.19); \draw(-.19,-.1)--(.19,-.1);
  \draw(0,.09)..controls(-.14,-.03)and(-.14,-.2)..(0,-.29) (0,.09)..controls(.14,-.03)and(.14,-.2)..(0,-.29);
\end{scope}
\node[concepttitle] at (2.05,2.81) {阿里云服务器};
\node[conceptbody,anchor=north] at (2.05,2.46)
 {2 核 / 2 GB / 40 GB\\现有 Web 与轻量测试接口\\先不扩容；既有网格不扩展};

\begin{scope}[shift={(7.75,3.46)},conceptline,draw=softgreen!75!black]
  \draw[fill=softgreen!8,rounded corners=2pt] (-.63,-.3) rectangle (.63,.3);
  \draw (-.77,-.4)--(-.77,.35)--(-.63,.35);
  \foreach \xx in {-.3,.27}{\draw(\xx,0)circle(.19);\draw(\xx,0)circle(.06);\draw(\xx,0)--(\xx+.15,.08)(\xx,0)--(\xx-.13,.08)(\xx,0)--(\xx,-.16);}
  \draw (-.3,-.3)--(-.3,-.4)--(.34,-.4)--(.34,-.3);
\end{scope}
\node[concepttitle] at (7.75,2.81) {i9 + RTX 4090};
\node[conceptbody,anchor=north] at (7.75,2.46)
 {1 TB；RAM 64／128 GB 待核实\\ML／DL 训练、实验与重计算\\本地过程数据，回传证据};

\begin{scope}[shift={(13.45,3.46)},conceptline,draw=softpurple]
  \draw[fill=softpurple!6,rounded corners=2pt] (-.37,-.42) rectangle (.37,.43);
  \foreach \yy in {.25,.03,-.2}{\draw(-.27,\yy)--(.1,\yy);\fill(.22,\yy)circle(.028);}
  \draw (-.22,-.34)--(.22,-.34);
  \draw[softred,line width=1pt] (.5,.2)--(.68,.38) (.5,.38)--(.68,.2);
\end{scope}
\node[concepttitle] at (13.45,2.81) {T640 + V100S 32 GB};
\node[conceptbody,anchor=north] at (13.45,2.46)
 {128 GB RAM / 2 TB\\网络不稳；CPU 待核实\\可选推理，首期不依赖};

% Legend defines proposal status and logical roles explicitly.
\draw[darkgray!20] (.15,1.18)--(15.35,1.18);
\draw[conceptflow] (.3,.86)--(1.05,.86);
\node[anchor=west,font=\fontsize{8}{10}\selectfont\sffamily] at (1.14,.86) {实线：拟建交互，待实施／验收};
\draw[conceptoption] (6.08,.86)--(6.83,.86);
\node[anchor=west,font=\fontsize{8}{10}\selectfont\sffamily] at (6.95,.86) {虚线：后续可选};
\draw[conceptpending] (10.3,.86)--(10.58,.86) (10.88,.86)--(11.16,.86);
\draw[darkgray!65,line width=.9pt] (10.66,.77)--(10.72,.95) (10.77,.77)--(10.83,.95);
\node[anchor=west,font=\fontsize{8}{10}\selectfont\sffamily] at (11.27,.86) {点线断口：未接通};
\node[conceptbody] at (7.75,.28)
 {双向箭头：请求与反馈；单向箭头：分派／授权发布。人物是逻辑角色，不是独占机器。\\核心也可直接调度执行器；所有返回的消息、产物均由平台登记和核验。};
\end{tikzpicture}
\caption{研发平台部署概念图：长期自有的平台核心组织多 Agent 协作，按任务调用 Codex 或确定性计算；设备按能力承载。全部自动链路均为待实施方案。}
\label{fig:rd-implementation-deploy}
\end{figure}

\textbf{协作与版本交接。}dot 云电脑适合文档、简单 Web 与轻量研发，并在其支持的环境内组织多 Agent 协作；这些能力可以独立于自建平台使用。它不是已注册的 SDK Worker，也不保证永久在线或提供平台可直接调用的通用外部 API。Mac 上的自建核心负责主要研发、任务状态与执行协调。两侧可使用获准的 GitHub 私有仓库，以固定 commit 或发布版本取得输入；候选代码、文档和必要验证记录经审阅后归档，不依赖同一对话或临时云环境永久保存成果。若未来需要自动分派或回传，须另行设计接口、身份、权限和验收，再确认可用性。

\textbf{角色、执行器与设备。}协调 Agent 把需求整理为有依赖和验收条件的任务；开发、算法、测试、运维等专业 Agent 按任务领取工作，通过平台选择 Codex 工程执行、确定性工具或计算作业。一个角色可有多个隔离运行实例，一台 Mac 也可顺序或有限并发承载多个角色。Codex 修改代码后应提交差异、命令结果与产物引用，由独立测试确认；算法 Agent 发起训练后应让 4090 的程序执行真实训练，而不是持续占用模型循环等待。

\textbf{计算、模型与对外服务。}Mac 按授权调用模型服务，并通过 Worker 协议向 4090 派发训练或重计算，接收状态、结果引用与证据。T640 当前网络不稳定，须先恢复网络，再核验驱动、模型适配和推理效果，之后才可作为可选模型服务；不预设增加第二张 GPU。阿里云现有资源用于普通 Web、轻量后台与经授权发布的测试接口；已有固定参数网格只是业务背景，本文不据此展开实盘系统或交易逻辑。

\textbf{数据与可用性边界。}GitHub 保存经批准上传的代码、文档和无秘密配置，PostgreSQL 保存平台持久状态；过程数据、原始资料、模型权重及大产物保留在受控本地存储。跨节点只传任务必需且获准的数据，运行服务及其日志保留在所属节点，Mac 登记位置、版本与校验值，不要求所有文件回流。Mac 休眠或关机时暂停其新任务派发与实时审批；已在 4090 独立启动的作业能否继续，由作业策略及恢复机制决定，不能依赖持续打开的远程终端。SDK 或 Codex 会话恢复不等于任务租约或审批自动恢复，必须由平台核对。

\subsubsection{首期配置与后续扩展边界}

首期按单用户、小并发配置：控制 API 一个实例，协调循环一个有效持有者，Agent 并发上限暂设为 2，Codex 工程任务先只允许 1 个，4090 初始只允许 1 个独占 GPU 作业。Codex 任务计入平台 Agent／执行额度，不能成为绕开并发或费用限制的嵌套执行入口；长计算等待时释放 Agent 槽。这些是保守实施起点，须按实测内存、显存、API 限额和耗时调整。角色数量不等于常驻进程数量，无任务时无需为每个岗位维持模型调用。

起步先建立 GitHub 版本基线、平台数据模型及本地模拟执行器，验收 Mac 的持久任务闭环；再分别接通 Agents SDK 与 Codex 的最小无副作用任务，并验证 dot 轻量开发成果的审阅交接。一般 Web 与阿里云轻量接口先在受控模拟环境形成证据，获准后再单独验证真实部署。随后启用 4090 Worker；T640 仅在网络、驱动和模型适配测试通过后接入推理。阿里云现有 2 vCPU、2 GB 内存、40 GB 存储先维持，不立即扩容；是否改为 2 核 4 GB 由资源监控、既有业务余量和验证结果决定。

只有出现明确的全天候调度需求，才讨论把控制服务及数据库迁到稳定在线的机器；迁移的是平台控制落点，不必重写角色、项目或执行器契约。新运行时通过适配器与契约测试接入，先验证权限、事件、取消、恢复和成果登记，再替换原实现。数据库集群、对象存储、专用向量库、复杂前端和外部队列均由实际容量或故障隔离需求触发，不能成为最小平台的前置条件。

\subsection{基础软件安装与环境配置}

\subsubsection{设备、系统与网络核验}

实施前先形成一份设备清单。Mac 核对 macOS 版本、Apple Silicon 架构、剩余磁盘及休眠策略；4090 台式机核对操作系统、CPU 型号、内存究竟为 64 GB 还是 128 GB、显卡和驱动；T640 核对 CPU、系统、GPU、可达性及机房访问方式。阿里云另核对当前系统、2 vCPU／2 GB／40 GB 规格、Web 与网格服务清单、端口、启动方式、磁盘和内存余量；dot 云电脑按每次任务实际可用环境记录工具与持久化条件，不套用 Mac 或云服务器的系统假设。未确认台式机和 T640 的系统前，不执行任何 Linux、Windows、驱动或 CUDA 安装脚本，也不为满足示例而重装系统。

\Needspace{8\baselineskip}
Mac 可先执行以下只读检查，输出供安装人员记录，不上传设备完整诊断：
\par\noindent\begin{minipage}{\linewidth}
\begin{lstlisting}
sw_vers; uname -m
sysctl -n hw.memsize
df -h "$HOME"
xcode-select -p; git --version
command -v uv; command -v docker
lsof -nP -iTCP:8000 -sTCP:LISTEN
lsof -nP -iTCP:5432 -sTCP:LISTEN
\end{lstlisting}
\end{minipage}\par


网络检查分别覆盖包下载、OpenAI API 出站 HTTPS、GitHub 访问和跨机控制链路；网络能打开网页不代表 API 账号、模型权限或 Worker 回报已经可用。先做无业务数据的连通测试，再做经授权的最小 API 调用。证书错误、访问拒绝或机房网络故障须查明原因，不关闭 TLS 校验，不临时把服务暴露到公网；网络尚未核验时仍可继续开发和本地模拟测试。

\subsubsection{Mac开发与控制环境安装}

若 Git / 编译工具缺失，由操作者运行 \texttt{xcode-select --install} 并完成系统安装界面。uv 使用官方安装脚本，先下载检查后执行；已有 uv 时直接核对版本。Docker Desktop 从官方页面选择 Apple Silicon 安装包，放入 Applications 后启动；由操作者确认许可与所需权限。Docker 虚拟机建议先预留 6 CPU、16 GiB 内存，供数据库及后续两个各 2 CPU、6 GiB 的本地沙箱作业使用；这是待实测的资源起点，不是性能保证。Mac 不安装 CUDA，数据库与作业使用独立容器。\cite{ch5-uv-project,ch5-docker-mac}

以下命令用于\textbf{首次创建尚不存在的项目目录}；已有代码仓库时应获取批准的版本并执行 \texttt{uv sync --locked}，不能再次初始化覆盖其配置：
\par\noindent\begin{minipage}{\linewidth}
\begin{lstlisting}
curl -LsSf https://astral.sh/uv/install.sh -o /tmp/uv-install.sh
less /tmp/uv-install.sh
sh /tmp/uv-install.sh
export PATH="$HOME/.local/bin:$PATH"
uv --version
uv python install 3.12
uv init --package --name rd-platform --python 3.12 "$HOME/agent-platform"
cd "$HOME/agent-platform"
uv add openai-agents openai openai-codex fastapi "uvicorn[standard]" \
  jinja2 python-multipart "psycopg[binary]" sqlalchemy alembic \
  pydantic-settings pyyaml httpx
uv add --dev pytest pytest-asyncio ruff
uv lock
uv sync --locked
uv run --locked python -c "import agents, openai_codex, fastapi, psycopg; print('imports OK')"
docker version
docker compose version
\end{lstlisting}
\end{minipage}\par


安装通过的判断至少包括：Python 显示为 3.12 系列、上述导入成功、Docker 客户端与服务端均能返回版本、Compose 可用。此时尚未证明模型调用、Codex runtime 启动、数据库迁移或业务流程正确。Codex SDK 应使用其发布构建固定的 runtime；记录实际 SDK 版本、runtime 版本与二进制来源，核对 ARM64 支持，并在隔离的临时工作区测试只读任务、事件回传及取消。认证由受保护的运行身份提供，不将主机秘密目录开放给任务，也不因桌面或 VS Code 已登录就跳过检查。SDK 可导入而 runtime、模型或授权不可用时，只禁用该执行能力，保留平台任务及模拟测试。\cite{codex-sdk,v121-codex-auth}

HTMX 采用官方发布的静态文件，固定版本及 SHA-256 后随应用提供，不要求另建 Node 前端工程，也不在页面运行时依赖第三方 CDN。

\subsubsection{平台目录、存储卷与依赖锁定}

目录按代码、运行数据和本机秘密分开。下表中的 \texttt{\textasciitilde} 表示当前操作者的主目录；自研配置加载器须展开为绝对路径，不能把它原样交给不支持展开的组件。

\begin{table}[H]
\centering\small\renewcommand{\arraystretch}{1.14}
\begin{tabularx}{\textwidth}{P{5.0cm}L}
\toprule
\textbf{路径或对象} & \textbf{内容与约束}\\
\midrule
\texttt{\textasciitilde/agent-platform/} & Git 工作区；含 \texttt{src/rd\_platform/}、\texttt{tests/}、\texttt{deploy/}、\texttt{migrations/} 与配置模板。\\
\rowcolor{lightgray}
\texttt{config/*.yaml} & \texttt{platform}、\texttt{agents}、\texttt{models}、\texttt{workers}、\texttt{policies} 五组配置；含角色、执行器选择、预算和节点能力。只提交无秘密模板及获准上传私有仓库的参数。\\
\texttt{\textasciitilde/agent-platform-data/} & \texttt{workspaces/}、\texttt{artifacts/}、\texttt{logs/}、\texttt{backups/}；原始数据和成果分别按项目授权访问。\\
\rowcolor{lightgray}
\texttt{\textasciitilde/.config/agent-platform/} & API 密钥、数据库口令、控制服务与 Worker 凭据引用；Codex 认证另按官方受保护存储管理。目录权限 700，秘密文件 600，不挂载给不可信任务。\\
Compose 命名卷 \texttt{pgdata} & Docker 管理 PostgreSQL 数据，挂载到容器的 \texttt{/var/lib/postgresql/data}；备份另存，不能靠复制运行中的目录代替。\\
\bottomrule
\end{tabularx}
\caption{平台文件、秘密与数据库存储边界。}
\end{table}

\texttt{pyproject.toml} 表达依赖约束，\texttt{uv.lock} 固定解析后的依赖，\texttt{.python-version} 在验收时固定实际 Python 补丁版本；部署使用 \texttt{uv sync --locked}，升级通过单独变更与回归测试完成。\cite{ch5-uv-project} 本文不虚构已经测试过的 SDK 版本号。还须记录 Agents SDK、Codex SDK 及其固定 runtime、uv、Docker、Compose、HTMX 版本与 PostgreSQL 镜像摘要；Mac ARM64 与计算节点 x86-64 分别测试，不复制 \texttt{.venv} 跨机使用。训练项目的 CUDA / PyTorch 环境另行锁定。

\subsubsection{服务地址、端口与启动配置}

Codex SDK 首期仅由受控本地 Worker 启动和管理本机 runtime，不开放 app-server 公网端口，也不把其 thread／turn 接口充当平台对外 API。跨机任务统一经平台 Worker 协议收发，远程 Codex 如未来需要，应在执行节点内就近启动并独立验收。\cite{codex-app}

初始入口是 \texttt{http://127.0.0.1:8000}，数据库仅监听宿主机的 \texttt{127.0.0.1:5432}；这两个地址都不能供另一台机器直接访问。4090 接入时再配置经核验的私有 HTTPS 控制地址与节点认证，不预设 VPN 已存在，不给 Worker 数据库口令，也不开数据库公网端口。

以下为 \texttt{deploy/compose.yaml} 的数据库基线。选择 \texttt{postgres:17-bookworm} 是固定大版本与发行基础，不是宣称它是最新版本；官方镜像提供 ARM64 架构。首次拉取和验收后，将标签替换为实际记录的 \texttt{postgres@sha256:...} 摘要。PostgreSQL 17 的数据卷路径与 18 及以后的默认路径不同，不能直接照搬新版本示例。\cite{ch5-postgres-image}
\par\noindent\begin{minipage}{\linewidth}
\begin{lstlisting}
name: agent-platform
services:
  postgres:
    image: postgres:17-bookworm
    restart: unless-stopped
    ports: ["127.0.0.1:5432:5432"]
    environment:
      POSTGRES_DB: rdp
      POSTGRES_USER: rdp_bootstrap
      POSTGRES_PASSWORD_FILE: /run/secrets/postgres_password
    secrets: [postgres_password]
    volumes: ["pgdata:/var/lib/postgresql/data"]
    healthcheck:
      test: ["CMD-SHELL", "pg_isready -U rdp_bootstrap -d rdp"]
      interval: 10s
      timeout: 5s
      retries: 5
secrets:
  postgres_password:
    file: ${HOME}/.config/agent-platform/postgres-password
volumes:
  pgdata:
\end{lstlisting}
\end{minipage}\par


由操作者在本机准备口令文件后才启动数据库；本文不提供可误用的默认密码。初始化管理用户 \texttt{rdp\_bootstrap} 禁止作为应用账号。安装脚本另建 \texttt{rdp\_migrator} 作为 \texttt{rdp} 模式及迁移对象所有者，\texttt{rdp\_app} 只取得运行所需表和序列权限；连接的 \texttt{search\_path} 明确为 \texttt{rdp,pg\_catalog}。初始化环境变量只对空数据目录生效，更改文件不会自动修改既有数据库口令。

\texttt{platform.yaml} 至少配置绑定地址、数据库主机 / 端口 / 库名 / 模式、数据根目录与并发限额；秘密值从本机受保护配置读取。\texttt{agents.yaml} 配置逻辑角色与默认执行器，\texttt{workers.yaml} 配置节点能力和执行槽；这些是平台自定义配置，须实现加载与校验，不能原样传入 SDK。\texttt{models.yaml} 只引用 \texttt{OPENAI\_API\_KEY} 等变量名，不写密钥正文。启动顺序为数据库、账号及迁移、控制 API、协调器、Worker，停止时先停止领取新任务，再处理在途作业并停止服务。

\par\noindent\begin{minipage}{\linewidth}
\begin{lstlisting}
cd "$HOME/agent-platform"
docker compose -f deploy/compose.yaml up -d postgres
docker compose -f deploy/compose.yaml exec postgres \
  pg_isready -U rdp_bootstrap -d rdp
# After rd_platform, role setup and migration code are implemented:
uv run --locked alembic upgrade head
uv run --locked uvicorn rd_platform.api:app --host 127.0.0.1 --port 8000
# Run the next two commands in separate supervised processes:
uv run --locked python -m rd_platform.coordinator
uv run --locked python -m rd_platform.worker --node mac-local
\end{lstlisting}
\end{minipage}\par


Alembic、Uvicorn 及两个 Python 入口是\textbf{待开发程序的目标启动约定}，还需配置加载、账号初始化和 Alembic 环境实现后才能运行。API 进程只处理入口及控制接口，协调循环不可随每个 Uvicorn 进程重复启动。运行检查应同时覆盖数据库健康、应用 \texttt{/healthz}、数据库就绪检查 \texttt{/readyz}、Worker 心跳以及一项无外部副作用的测试任务；仅有网页可打开不算通过。稳定后再配置 macOS 服务管理及日志轮转，且不向 Agent 或作业容器暴露 Docker 宿主机套接字。





\subsection{平台程序设计与接口实现}
\label{sec:platform-implementation}

\subsubsection{程序模块与代码仓库结构}

平台采用私有 GitHub 仓库 \texttt{agent-platform} 管理 Python 包 \texttt{rd\_platform}、网页模板、数据库迁移、配置模板及平台 requirements、design、tests；组织规则在独立的 \texttt{organization} 仓库维护，各业务或验证项目另有项目仓库。名称为拟定约定，实施时登记真实仓库地址与权限；任务状态与运行文件不提交 Git。控制端在 Mac 的 Python 3.12 虚拟环境运行，FastAPI 提供服务，Jinja2 渲染页面，HTMX 刷新局部状态；数据库为本机容器中的 PostgreSQL 17。自建 Worker 领取任务后，经适配器调用编排框架、开发执行器或确定性计算工具。\textbf{自有平台内核是长期维护的控制系统}：负责项目与版本、任务及依赖、权限和审批、持久状态、租约、成果与交付判定。OpenAI Agents SDK 是当前拟选的编排框架，Codex SDK 是拟接入的开发执行器；二者均位于可替换的适配层，不拥有平台任务的最终状态。\cite{rdRuntimeSdk,codex-sdk}

实现首先固定领域模型与版本化契约，再接通一个编排器和一个执行器。这样更换框架、SDK 版本或模型服务时，主要改动局限在适配、上下文转换与兼容性测试，不要求重做任务数据库、审批、工作台和成果管理。当前方案属于待开发的功能设计，不宣称这些模块已经实现。

\noindent\begin{minipage}{\linewidth}
\begin{lstlisting}[caption={平台仓库拟定结构；属于实施目标，尚非已建程序}]
agent-platform/
  pyproject.toml  uv.lock  .env.example  deploy/compose.yaml
  config/{platform,agents,models,workers,policies}.yaml
  docs/{requirements,design,operations}/
  prompts/{organization,coordinator,developer,algorithm,test,ops}.md
  # prompts are materialized from the pinned organization revision
  src/rd_platform/
    api/         # exports app; user/worker API, authentication
    coordinator.py # persistent coordination entry point
    web/         # Jinja2 templates, HTMX partial views
    domain/      # projects, versions, tasks, approvals, reviews
    contracts/   # versioned task, progress, result and adapter ports
    storage/     # SQL repositories, artifact registry, event ledger
    scheduler/   # claim, lease, dependency, fencing, reconciliation
    context/     # portable context package and source revisions
    runtime/     # orchestration driver, usage and checkpoint policy
    adapters/    # agents_sdk, codex_sdk, codex_exec, compute
    tools/       # guarded file, git, test and job gateways
    worker/      # __main__.py; pull, execution, heartbeat, journal
  migrations/  tests/{unit,contract,integration,acceptance}/
\end{lstlisting}
\end{minipage}


图~\ref{fig:platform-core-modules}~把内核、持久状态与可替换运行组件分开。\texttt{contracts} 定义边界，\texttt{domain} 实施状态与验收规则，\texttt{scheduler} 实施确定性调度，\texttt{context} 生成可迁移任务上下文；\texttt{runtime} 只协调适配调用、用量登记和恢复策略。\texttt{adapters} 将统一请求映射到具体 SDK／进程，将原始事件归一化后回报内核，不直接改写领域表。

\begin{figure}[htbp]
\centering
\begin{tikzpicture}[every node/.style={font=\footnotesize\sffamily,align=center},
 module/.style={rounded corners=3pt,draw=primaryblue!45,fill=white,minimum height=1.0cm,text width=4.15cm,inner sep=5pt},
 link/.style={-{Stealth[length=2mm]},draw=primaryblue,line width=.85pt}]
\node[module,fill=primaryblue!7,text width=14.2cm,minimum height=.65cm] (entry) at (0,0) {工作台与控制 API：接收意图、显示进度、提交决定};
\node[module] (domain) at (-4.9,-1.55) {领域与成果服务\\项目／任务／审批／独立验收};
\node[module] (schedule) at (0,-1.55) {调度与恢复服务\\依赖／配额／租约／故障核对};
\node[module] (guard) at (4.9,-1.55) {上下文与工具网关\\固定版本／权限／预算／审计};
\begin{scope}[on background layer]
\node[draw=primaryblue!55,fill=primaryblue!4,rounded corners=4pt,
 fit=(domain)(schedule)(guard),inner xsep=7pt,inner ysep=12pt] (core) {};
\end{scope}
\node[font=\footnotesize\bfseries\sffamily,text=primaryblue,fill=primaryblue!4,inner sep=2pt] at (0,-.82) {自有平台内核};
\node[module,fill=softgreen!7,text width=6.6cm,minimum height=1.05cm] (state) at (-3.8,-3.75) {平台持久状态\\PostgreSQL 17＋不可变成果存储};
\node[module,fill=accentorange!7,text width=6.6cm,minimum height=1.05cm] (adapters) at (3.8,-3.75) {可替换适配器\\Agents SDK 编排／Codex 开发／计算执行};
\draw[link] (entry.south)--(core.north);
\draw[link] (domain)--(schedule);
\draw[link] (schedule)--(guard);
\draw[link,{Stealth[length=2mm]}-{Stealth[length=2mm]}] (domain.south)--(state.north -| domain.south);
\draw[link,{Stealth[length=2mm]}-{Stealth[length=2mm]}] (guard.south)--node[right,font=\scriptsize\sffamily,fill=white]{任务／进度／结果契约}(adapters.north -| guard.south);
\end{tikzpicture}
\caption{程序模块边界：平台保有任务与成果权威，框架和执行器通过稳定契约接入，回报先由内核校验，再写入持久状态。}
\label{fig:platform-core-modules}
\end{figure}

入口模块不得直接修改数据库状态；应调用 \texttt{domain} 服务，在事务内检查版本、权限与状态前置条件。\texttt{runtime} 不直接拼接任意 shell 命令，\texttt{tools} 不自行判定任务完成。控制端代码位于 \texttt{\textasciitilde/agent-platform}，运行数据位于 \texttt{\textasciitilde/agent-platform-data}，主机私有配置位于 \texttt{\textasciitilde/.config/agent-platform}。各任务另建工作副本及输出区，不能在平台自身安装目录执行模型生成的代码。

五个 YAML 文件是\textbf{平台自定义配置}，需要实现加载、类型校验、环境变量解析和配置快照，不能直接交给 Agents SDK 或 Codex SDK 执行。未知字段、未解析模型名或缺少必要凭证时应拒绝启动对应能力。配置模板可入库，密钥只由受保护的主机配置注入；\texttt{uv.lock} 的真实依赖版本须在实施与兼容性检查后生成，不以本文示例冒充已测试锁定版本。

\subsubsection{项目、版本、任务与运行数据模型}

先区分三个生命周期：\textbf{task} 是有目标、授权范围和验收条件的持久业务任务；\textbf{attempt} 是某个 Worker 在一代租约下进行的一次执行尝试；\textbf{runtime session} 是某一运行后端的会话或线程。任务重试须新建 attempt，会话可在兼容条件下继续，也可重建；一个任务可以先后跨多个框架会话，纯计算任务则不需要模型会话。\texttt{run} 记录一次适配器调用或模型／工具运行；平台 \texttt{run\_id}、Codex thread／turn ID、提供方 response ID 和 trace ID 各自登记，不能互相充当主键。

主键使用 UUID，时间统一存为带时区时间；下例中的 \texttt{p1/v1/t1/at1/r1/a1} 仅是便于阅读的标识缩写。数据库使用 \texttt{rdp} 库与同名 schema，连接显式设置 \texttt{search\_path}；\texttt{rdp\_app} 只获运行期必需权限，DDL 由独立迁移身份执行。数据库中的对象关系必须建立外键，跨项目引用由服务层再次校验。任务执行所采用的 GitHub 仓库与不可变 commit 作为独立版本引用保存，不能与数据库对象自身的修订号混用。关键表的最小字段如下。

\begingroup\small
\begin{longtable}{P{2.55cm}P{5.65cm}P{6.5cm}}
\toprule
\textbf{表} & \textbf{主要字段} & \textbf{约束与用途}\\
\midrule\endfirsthead
\toprule\textbf{表} & \textbf{主要字段} & \textbf{约束与用途}\\\midrule\endhead
\texttt{projects}\newline\texttt{versions} & 项目标识、约定成果引用；版本标识、项目、需求修订、状态、冻结时间 & 版本隶属项目；冻结后保留需求快照。新增需求写入后续版本 backlog，不直接改写冻结包。\\
\texttt{inputs}\newline\texttt{tasks} & 输入正文或成果引用、提交人；任务标识、版本、输入、角色、类型、状态、修订号、优先级、可执行时间、组织与项目仓库 commit & 输入保留用户原意；任务保存验收条件、输出类型及所需 skills 路径／摘要。任务依赖另存 \texttt{task\_dependencies}，禁止成环。\\
\texttt{attempts}\newline\texttt{runs} & 尝试标识、任务、序号、Worker、租约代次、起止与错误；运行标识、尝试、适配器／版本、会话引用、角色／模型／配置快照 & 重试建立新尝试；一次尝试可包含多个后端调用。外部 thread、turn、response 标识另存映射，不替代平台标识。\\
\texttt{jobs}\newline\texttt{workers} & 作业标识、计算子任务／尝试、父任务、节点、输入清单、命令模板、外部作业标识、状态；Worker 能力、心跳、禁用状态 & 长训练与 Agent 推理解耦；作业返回不会自动验收任务。节点内存等资源以实测注册值为准。\\
\texttt{artifacts}\newline\texttt{artifact\_\allowbreak revisions} & 成果标识、类型；不可变修订、来源任务／运行、存储键、大小、SHA-256、Git commit、清单状态 & 只引用确定修订；物理存储与逻辑成果分开。文件校验通过后才可供下游读取。\\
\texttt{tool\_calls}\newline\texttt{approvals} & 操作标识、参数摘要、幂等键、外部操作标识／回执引用、结果状态；审批人、动作／参数哈希、范围、有效期、决定 & 请求与观察结果分别保存。批准只能匹配原动作与修订，参数变更须重新核对授权。\\
\texttt{reviews}\newline\texttt{events} & 候选修订、测试运行、结论、证据引用；事件序号、对象、前后修订、操作者、时间、摘要 & 评审必须引用候选快照。事件与业务变更同事务写入，网页断连后按序号补读。\\
\texttt{messages}\newline\texttt{handoffs} & 消息／交接标识、任务及修订、成果引用、发送／接收主体、幂等键、发送时间；接受／拒绝回执及时间 & 保存跨角色交接依据与接收确认。回执绑定原任务修订和输入清单；接受交接不表示任务已完成，资料改变须重新确认。\\
\bottomrule
\end{longtable}
\endgroup

\texttt{tasks} 还保存 \texttt{lease\_owner}、\texttt{lease\_until}、\texttt{lease\_generation}、\texttt{cancel\_requested\_at} 及重试计数。框架无关的上下文包、规范化会话条目与原始后端检查点分别保存。\texttt{sessions} 保存平台会话标识、任务、后端种类、后端会话标识和兼容版本；\texttt{session\_items/checkpoints} 按序号关联。服务重启不能使这些记录丢失，但“已保存”不等于“可在任意新框架直接加载”。对任务尝试序号、业务幂等键及成果修订设置唯一约束；JSONB 只承载可变参数，状态、外键和租约条件不能全部藏在不受约束的 JSON 中。

一次完整数据流为：用户输入 \texttt{i1} 属于项目 \texttt{p1}，版本 \texttt{v1} 引用冻结需求 \texttt{a0@1} 及对应仓库／commit／路径；调度 \texttt{t1} 产生尝试 \texttt{at1} 与运行 \texttt{r1}，必要时派出作业 \texttt{j1}；结果登记为 \texttt{a1@1}，独立测试任务引用该修订，交付清单再引用需求、代码提交及测试证据。网页展示的“完成”来自这些记录，不从聊天措辞或进程退出码推断。

\textbf{稳定契约。}下列字段属于本平台的版本化协议，存为可校验的结构化数据，不嵌入 SDK 的 Agent、Thread 或 RunState 对象。它们可用于进程内调用或 Worker API；协议版本变化须有迁移和兼容测试。

\begingroup\small
\begin{longtable}{P{2.0cm}P{6.2cm}P{6.5cm}}
\toprule\textbf{契约} & \textbf{最小信息} & \textbf{内核检查与含义}\\\midrule\endfirsthead
\toprule\textbf{契约} & \textbf{最小信息} & \textbf{内核检查与含义}\\\midrule\endhead
任务请求 & \texttt{schema\_version}、任务／版本／修订、目标与排除项、输入修订、验收与停止条件、角色、能力需求、权限及预算快照、源代码版本 & 区分“要完成什么”与“由谁执行”；领取后补充尝试、租约代次、截止时间及适配器选择。不得由模型自行扩大授权。\\
进度事件 & 任务／尝试／运行标识、事件幂等标识、生产者序号、阶段、阻塞原因、已观察事实、证据与用量引用 & 数据库分配全局事件序号；重复事件去重，迟到事件保留来源。阶段和估计进度不等于任务终态。\\
结果清单 & 任务／尝试、租约代次、候选成果修订、代码 commit、测试证据、操作账本引用、用量、错误及已知副作用状态 & 结果代表执行报告。内核检查租约、输入一致性、成果校验、独立测试与交付条件后，才决定任务状态。\\
\bottomrule
\end{longtable}
\endgroup

任务请求引用平台保存的授权快照，不携带密钥正文。进度中的时间和描述属于观测数据，不能替代数据库租约时间。模型文本、退出码、SDK 的“completed”以及 Codex 的最终回答都只能作为结果输入，不能直接把 \texttt{tasks.status} 改为 \texttt{succeeded}。

\subsubsection{用户工作台与后台 API}

工作台设项目与版本页、意图提交页、任务详情、待决事项和成果页。用户填写目标、约束与资料引用后即可提交；系统需要澄清或授权时才形成待决事项。任务详情呈现当前阶段、阻塞、已验证成果、尝试历史和用量，技术日志按需展开。UI 与 Worker 使用同一领域服务，身份和可操作范围不同；下表是拟实现契约，路径均以 \texttt{/api/v1} 为前缀。

单用户工作台仍需认证；回环监听不等于身份校验。首次初始化建立本地主体，登录后使用服务端会话与随机会话令牌，Cookie 设 \texttt{HttpOnly/SameSite=Strict}，HTTPS 入口再强制 \texttt{Secure}；明文 HTTP 仅允许首期回环访问。所有 Cookie 认证写请求校验 CSRF token、Origin 和允许的 Host，拒绝跨站提交；退出使会话失效。Worker API 使用独立 Bearer 凭证，不共用浏览器 Cookie；接口均检查项目范围。验收至少覆盖无登录读写、错误 CSRF、伪造 Origin／Host、跨项目和节点凭证越权请求均被拒绝。

\begingroup\small
\begin{longtable}{P{5.2cm}P{3.3cm}P{6.2cm}}
\toprule\textbf{方法与路径} & \textbf{关键输入} & \textbf{返回及语义}\\\midrule\endfirsthead
\toprule\textbf{方法与路径} & \textbf{关键输入} & \textbf{返回及语义}\\\midrule\endhead
\texttt{POST /projects/\{p\}/inputs} & 意图、约束、资料引用；幂等键 & \texttt{202}：输入与根任务标识、查询地址；表示已持久接收。\\
\texttt{POST /versions/\{v\}/freeze} & 需求成果修订、预期版本修订 & \texttt{200}：冻结快照；未经必要决策或修订冲突则拒绝。\\
\texttt{GET /tasks/\{t\}} & 已认证项目访问范围 & \texttt{200}：状态、修订、活动尝试、结果引用；不返回凭证。\\
\texttt{GET /projects/\{p\}/events} & \texttt{after\_seq}、\texttt{limit} & \texttt{200}：按序事件及下一游标；起步由 HTMX 轮询。\\
\texttt{POST /tasks/\{t\}/cancel} & 原因、预期修订号 & \texttt{202}：取消请求已登记；尚未证明执行停止。\\
\texttt{POST /approvals/\{id\}/decision} & 批准／拒绝、动作哈希 & \texttt{200}：决定已记录；执行前仍检查时效及任务状态。\\
\texttt{POST /handoffs/\{h\}/receipt} & 接受／拒绝、原因、交接与任务修订、输入清单摘要 & \texttt{200}：接收回执已登记；仅指定接收方可提交，不代表成果验收。\\
\texttt{GET /artifacts/\{a\}/revisions/\{n\}} & 成果与修订 & \texttt{200}：元数据、校验值、授权下载入口。\\
\texttt{POST /workers/claim} & 能力、空闲槽、请求标识 & \texttt{200}：任务与租约；无任务 \texttt{204}。\\
\texttt{POST /attempts/\{a\}/heartbeat}\newline\texttt{POST /attempts/\{a\}/complete} & 租约代次；心跳或结果清单 & \texttt{200}：确认续租或提交；旧代次／失效租约返回 \texttt{409}。\\
\texttt{POST /attempts/\{a\}/observations} & 旧尝试、外部作业状态与证据清单 & \texttt{202}：仅登记隔离证据，待控制器核对；不续租、不改变任务终态。\\
\bottomrule
\end{longtable}
\endgroup

写接口要求 \texttt{Idempotency-Key}，服务保存主体、路径、请求体哈希与响应；同键同请求返回原结果，同键异体返回 \texttt{409}。已接受的完成提交或审批决定发生网络重试时只重放原响应，不再次执行动作。状态写入另带 \texttt{expected\_revision} 防止覆盖他人修改。交接记录由领域服务在委派或成果交接时建立，回执与相应事件同事务保存；不增加消息代理，仍通过 PostgreSQL 与现有事件轮询传递。统一错误体含 \texttt{code/message/request\_id/retryable}；\texttt{401/403} 区分未认证与权限拒绝，\texttt{422} 表示字段错误，\texttt{409} 表示状态冲突。\texttt{retryable} 只是服务建议，调用方仍须核对动作的副作用状态。

\par\noindent\begin{minipage}{\linewidth}
\begin{lstlisting}[caption={意图接收与响应示例；标识均为缩写}]
POST /api/v1/projects/p1/inputs
Idempotency-Key: input-20261012-001
{"intent":"Build a reproducible data-validation capability",
 "constraints":{"scope":"research_and_development"},
 "artifact_refs":[{"artifact_id":"a0","revision":1}]}

HTTP 202
{"input_id":"i1","task_id":"t1","status":"ready",
 "revision":1,"status_url":"/api/v1/tasks/t1"}
\end{lstlisting}
\end{minipage}\par

\subsubsection{持久任务队列与调度状态机}

任务状态的唯一权威是自有 PostgreSQL 数据库；后端会话、SDK trace 与聊天记录均是附属证据。起步直接以 PostgreSQL 的任务表作为持久队列，不另加 Redis、Celery 或常驻聊天中转。任务主状态为 \texttt{draft/blocked/ready/running/waiting/awaiting\_review/succeeded/failed/cancelled/reconcile}。依赖、授权或资源不满足时不能进入 \texttt{ready}；等待长作业时进入 \texttt{waiting} 并释放 Agent 执行槽。\texttt{reconcile} 表示必须先核对外部事实。实现／研究任务须独立测试通过且交付条件满足，领域服务才能写入 \texttt{succeeded}；协调与评审任务按各自契约核验完成，评审任务不递归要求另一个评审任务。

\begin{figure}[htbp]
\centering
\begin{tikzpicture}[node distance=0.55cm, every node/.style={font=\footnotesize\sffamily},
 s/.style={draw=primaryblue!50,rounded corners,fill=primaryblue!5,minimum height=.68cm,align=center,inner sep=4pt},
 a/.style={-{Stealth[length=2mm]},draw=primaryblue,thick}]
\node[s] (b) {blocked};
\node[s,right=of b] (q) {ready};
\node[s,right=of q] (r) {running};
\node[s,right=of r] (v) {awaiting\_review};
\node[s,right=of v] (d) {succeeded};
\node[s,below=0.8cm of r] (w) {waiting / reconcile};
\draw[a] (b)--(q); \draw[a] (q)--(r); \draw[a] (r)--(v); \draw[a] (v)--(d);
\draw[a] (r)--(w); \draw[a] (w.west)-|(q.south);
\draw[a] (v.north) to[bend right=30] node[above]{退回修正} (q.north);
\end{tikzpicture}
\caption{任务主路径；失败与取消是经条件检查后的终态，不由心跳超时直接推断}
\end{figure}

多个 Worker 的领取由控制端在\textbf{同一数据库事务}中完成：筛选已就绪且能力匹配的任务，锁行，增加租约代次，建立尝试记录，再返回任务。数据库锁不能跨越模型调用或训练全过程。PostgreSQL 明确指出 \texttt{SKIP LOCKED} 可用于多消费者队列表的锁竞争处理，但不适合一般一致性查询；它也不提供端到端恰好一次执行保证。\cite{rdRuntimePg}

\Needspace{15\baselineskip}
\par\noindent\begin{minipage}{\linewidth}
\begin{lstlisting}[caption={领取事务的关键 SQL；省略已校验的能力与配额条件}]
BEGIN;
SELECT id FROM tasks
 WHERE status = 'ready' AND available_at <= now()
   AND cancel_requested_at IS NULL
 ORDER BY priority DESC, created_at, id
 FOR UPDATE SKIP LOCKED LIMIT 1;
-- :task_id is the row selected above, on this connection.
UPDATE tasks SET status = 'running', lease_owner = :worker_id,
 lease_until = now() + interval '90 seconds',
 lease_generation = lease_generation + 1, revision = revision + 1
 WHERE id = :task_id
 RETURNING lease_generation;
-- Insert attempts and events using the returned generation.
COMMIT;
\end{lstlisting}
\end{minipage}\par


租约初值 90 秒、心跳 20 秒是起步配置，不是跨公网可靠性的保证；续租使用数据库时间，并同时检查任务、尝试、Worker 身份、租约代次及尚未过期；过期租约不能被迟到心跳复活。额度、GPU 槽等共享资源须在同一事务锁定资源计数，不能先数空位再分别领取。任务完成也必须检查当前租约有效、代次一致且未被取消；旧 Worker 的迟到输出只进入隔离区，不能覆盖现行任务或成果。离线节点恢复后经 \texttt{observations} 接口上传证据；当前控制器核验取消、租约代次、作业事实和校验值后，才可显式采纳为新登记成果或恢复依据，不能让旧 \texttt{complete} 绕过拒绝。

交付语义按\textbf{至少一次尝试、幂等提交}设计。租约过期先转核对状态：只读计算或明确可重放的步骤可建立新尝试；已发出写请求而结果未知时，查询外部作业／操作标识，禁止盲目重发。取消分两步：先记录 \texttt{cancel\_requested\_at}，停止领取新步骤；收到进程退出、远程作业终止或可信查询证据后才写 \texttt{cancelled}。离线 Worker、失联请求和不能确认的副作用保持可见的核对状态。

\subsubsection{Worker 协议与故障恢复}

长计算沿用同一任务队列：提交时创建 \texttt{compute} 子任务，带 \texttt{parent\_task\_id} 及独立的尝试、租约和 \texttt{job\_id}。Mac 父任务保存检查点后转 \texttt{waiting}，结束当前尝试并释放 Agent 槽和租约；4090 只领取计算子任务，不能继承父任务租约。子任务结果验证通过后，事务内唤醒父任务为 \texttt{ready}，以新尝试装配持久上下文继续；只有后端与版本兼容时才直接加载原检查点。纯计算尝试的模型运行／会话引用可空，作业与来源 Agent 运行的关系另行保留。

Worker 是自建 Python 程序，包含领取循环、心跳、受限执行器、输入缓存和提交日志。Mac Worker 经适配层执行编排调用、Codex 开发任务与轻量工具；RTX 4090 节点优先执行 ML/DL 训练作业。Worker 通过经授权配置的私网 HTTPS 控制端主动拉取，使用独立节点凭证和最小权限；远程节点不持有数据库账户，也不直接连接 5432 端口。认证成功只证明节点身份，任务与资源授权仍由控制端决定。

\par\noindent\begin{minipage}{\linewidth}
\begin{lstlisting}[caption={workers.yaml 的最小配置片段；私网入口需实施时配置}]
schema_version: 1
control_url_env: PLATFORM_CONTROL_URL
credential_file_env: PLATFORM_WORKER_CREDENTIAL_FILE
transport: {tls_verify: true}
lease: {seconds: 90, heartbeat_seconds: 20}
workers:
  mac-local:
    capabilities: [orchestration, code_dev, cpu_test]
    max_concurrency: 2
  rtx4090:
    capabilities: [gpu_train]
    max_concurrency: 1
    resource_discovery_required: true
  t640:
    enabled: false
\end{lstlisting}
\end{minipage}\par


领取响应包含任务／尝试标识、可空的运行／会话引用、租约代次、角色配置哈希、组织与项目代码 commit、必要 skills 清单、输入成果修订清单和输出约束。\texttt{PLATFORM\_\allowbreak CONTROL\_URL} 在本机阶段指向回环服务；远程阶段才设为有可信证书的私网 HTTPS 地址。不能把 \texttt{127.0.0.1} 当成远程 Mac 地址，也不能为连通性关闭 TLS 校验。RTX 4090 主机的 RAM 尚未核定，首次注册前读取实际内存与可用磁盘；调度不能预设为 128GB。

文件通过成果 API 交换，清单携带 \texttt{artifact\_id/revision/size/sha256}，下载后先验证再放入任务沙箱；提交采用“上传临时对象、校验清单、登记不可变修订”三步。超时可重传同一内容，最终登记按幂等键去重。Mac 路径不应直接传给 Windows／Linux Worker，也不假定三台机器共享文件系统。训练输入可按哈希缓存，检查点按独立修订上传；跨弱网大文件失败不阻塞已可完成的本地任务。

Worker 本地日志在启动子进程前记录 \texttt{attempt\_id/job\_id} 和操作标识，重启先查询原进程或外部作业，再恢复提交。仅有 PID 不足以确认同一进程，应结合启动时间与作业令牌。任务目录按尝试隔离，工具执行前再验证当前租约；失去租约立即停止发起新副作用。已开始且事先获准离线运行的限时计算，可仅在既定资源与时限内继续写入本尝试的隔离本地输出；不得启动新作业、外部业务写入或发布。恢复联网后只先回传 \texttt{observations}，由控制器核对后决定是否采纳，不能以离线完成绕过 fencing。

执行适配器统一暴露 \texttt{submit/query/cancel/recover} 四类操作，但\textbf{统一函数名不代表各后端具备同等能力}。每次部署须登记并测试能力清单；不支持的操作返回 \texttt{unsupported}，状态不可证实时返回 \texttt{unknown}，由内核阻塞或转入 \texttt{reconcile}，不能伪造成功。

\begingroup\small
\begin{longtable}{P{2.65cm}P{2.9cm}P{2.9cm}P{2.9cm}P{2.85cm}}
\toprule\textbf{适配器} & \textbf{提交 submit} & \textbf{查询 query} & \textbf{取消 cancel} & \textbf{恢复 recover}\\\midrule\endfirsthead
\toprule\textbf{适配器} & \textbf{提交 submit} & \textbf{查询 query} & \textbf{取消 cancel} & \textbf{恢复 recover}\\\midrule\endhead
Agents SDK 编排 & 平台创建运行，再调用锁定版本 Runner & 平台事件与运行记录；不假设通用远程作业查询 & 停止运行并阻止新工具动作；在途外部动作另核对 & 兼容状态可继续；否则按持久上下文创建新运行\\
Codex SDK / app-server 开发 & 以固定工作区启动线程／turn，记录外部标识 & 映射后端事件及可用读取接口，补录候选成果 & 仅在所选接口支持时发中断；等终止事件并核对子进程 & 可恢复线程须校验存储、版本和工作区；否则重建上下文\\
Codex CLI 批处理备选 & 受控进程运行 \texttt{codex exec}，收集 JSONL & 进程、事件日志与平台本地账本；不假设远程查询 & 进程监管器请求停止，不能据此承诺撤销已生效动作 & 会话存在且兼容时按明确 ID resume；否则新尝试\\
确定性计算执行器 & 固定命令模板与作业令牌；参数校验后启动 & 本地监管器或已接入的作业系统查询 & 监管器终止进程组／后端取消，等待可验证终态 & 仅支持已登记的业务检查点；无检查点则核对后重跑\\
外部／人工环境 & 仅登记获准任务及交接包 & 人工回执与审阅证据 & 自动取消未验证则不支持 & 自动恢复未验证则不支持\\
\bottomrule
\end{longtable}
\endgroup

Codex SDK 用于程序化开发执行；当前文档包含 TypeScript 与 Python 路径；本方案拟选 Python \texttt{openai-codex} 包，其 \texttt{openai\_codex} 接口封装本地 app-server，并锁定配套运行时。app-server 提供线程／turn 与事件接口，CLI 非交互路径提供事件流及会话继续能力；具体可用操作须绑定实际采用的包、CLI 与协议版本。上表是\textbf{平台拟实现的能力映射}，不是声称 SDK 已统一提供这四个方法。\cite{codex-sdk,codex-app,v121-codex-exec}

提交前先持久登记 \texttt{operation\_id} 和幂等键，再调用后端并保存返回标识；后端若不支持幂等提交，断连后的重试先查原线程、进程或回执。\texttt{recover} 返回恢复计划与证据，不绕过新尝试的租约校验。\texttt{cancel} 被接受只表示请求已发出，\texttt{query} 报告的“已结束”还须与在途工具及成果记录核对。所有权威状态写入和结果提交保留租约代次 fencing；失去租约后，除上述已获准计算的隔离本地输出外，不再发起新副作用。

未来 T640 位于弱网 IDC，可另接模型服务适配器，默认关闭且不处于控制、数据库、队列或交付的关键链路。若其服务只实现 Chat Completions，应采用对应适配器并单独验证工具调用与结构化输出，不能只替换 Responses 的 \texttt{base\_url} 就认定兼容。

\subsection{Agent 编排、模型接入与运行配置}
\label{sec:openai-runtime-configuration}

\subsubsection{编排边界、模型路由与调用预算}

当前拟采用 Python OpenAI Agents SDK 实现需要模型推理的协调、研究和评审流程，经模型适配器调用 Responses API；代码开发可委派给独立的 Codex 执行适配器。Agents SDK 组织模型与工具循环，并可表达委派；它不是任务数据库、权限系统或持久队列，也不是平台永久不可替换的内核。以后改用另一编排实现，应继续遵守相同任务、进度、结果及授权契约。\cite{rdRuntimeSdk}

协调中的两类动作必须分开：\textbf{LLM 规划}根据目标、约束与证据提出拆解、依赖和返工建议；\textbf{确定性调度}由平台程序校验 DAG、权限、预算、资源和租约，在事务中领取与推进。模型不得直接决定“有配额”、给自己发租约或宣告独立测试通过。已定义的重复流程可以不调用 LLM，直接执行规则与状态机。

\begin{table}[htbp]
\centering\small
\begin{tabularx}{\textwidth}{P{2.35cm}LP{5.4cm}}
\toprule\textbf{配置维度} & \textbf{回答的问题} & \textbf{例子与边界}\\\midrule
组织角色 & 承担什么职责、按何种证据验收？ & 协调、开发、算法、独立测试、运维；不等于五个常驻进程。\\
编排框架 & 如何组织模型与工具循环？ & 当前 Agents SDK；可由经契约测试的其他实现替换。\\
模型路由 & 某次推理使用哪个模型／服务？ & \texttt{deep/standard/economical} 是平台别名；仅在相应后端能力范围内解析。\\
执行器选择 & 谁接受并执行该类工作？ & Codex 开发、受限 CPU 测试、GPU 训练；不是更换模型 ID。\\
运行节点 & 代码、工具或作业在哪里运行？ & Mac、4090、获准云环境；节点不决定模型账户或付费方式。\\
认证与计费 & 使用哪种身份、额度及账单？ & Platform API key 与 Codex 的 ChatGPT 登录分别核验、分别记账。\\
\bottomrule
\end{tabularx}
\caption{六个独立但相互约束的选择维度；不能以“换模型”代替执行器或认证配置。}
\end{table}

平台显式记录每次运行的模型或可获得的服务标识，避免 SDK 默认值随升级变化；部署时核对账户可用性、协议和所需工具能力，再将真实模型 ID 写入受控配置。较强模型不带来额外文件、网络或发布权限。Codex 有自己的模型与运行配置，平台路由别名不可未经映射就直接传入。\cite{rdRuntimeModels}

\par\noindent\begin{minipage}{\linewidth}
\begin{lstlisting}[caption={models.yaml 与 platform.yaml 中的拟定预算项}]
# models.yaml: required environment values are real account model IDs.
provider: openai_responses
api_key_env: OPENAI_API_KEY
routes:
  standard: {model_id_env: OPENAI_MODEL_STANDARD}
  deep: {model_id_env: OPENAI_MODEL_DEEP}
  economical: {model_id_env: OPENAI_MODEL_ECONOMICAL}
fallback: {automatic: false}
self_hosted: {enabled: false, protocol: chat_completions}

# platform.yaml: initial engineering limits, not spending approval.
runtime:
  max_agent_runs: 2
  max_turns_per_run: 12
  run_timeout_seconds: 600  # orchestration runs only
budget:
  per_task_usd: 2.00
  per_day_usd: 10.00
  on_limit: pause_for_decision
# adapters is a platform-owned mapping, not an SDK config format.
adapters:
  orchestration: {driver: agents_sdk, auth_mode: platform_api_key}
  code_dev: {executor_id: codex-mac}
  compute: {driver: restricted_job_runner}
  paid_fallback: false
\end{lstlisting}
\end{minipage}\par


上述 YAML 保持原有五个配置文件，只在 \texttt{platform.yaml} 增加适配映射；其中 \texttt{code\_dev} 通过 \texttt{executor\_id: codex-mac} 引用 \texttt{workers.yaml} 的同名执行器，驱动、认证方式、节点、工作区与权限均以该登记为准。\texttt{driver/auth\_mode/paid\_fallback} 是待实现的平台字段，不是 SDK 官方参数。配置加载器将抽象选择映射为锁定版本的实际接口，启动探针验证包与运行时、认证路径及支持能力；映射缺失即暂停该执行能力，不能悄悄改用另一个付费入口。

\textbf{认证与预算分开核验。}Codex 通过 ChatGPT 登录使用符合条件的订阅权益，与通过 API key 使用按量计费服务是不同路径；具体额度、账户权限和能力以上线时的账户及官方规则为准。即使开发执行器使用 Codex 订阅权益，协调 Agent 只要仍用 Platform API key 调用模型，就仍按 API 路径计费，不能把整个平台视为已由订阅覆盖。\cite{v121-codex-auth}

若将来希望私有平台客户端直接使用 ChatGPT plan usage，须先确认官方允许的客户端类型、注册／批准条件、账户资格和用户 OAuth 授权，并完成协议与范围校验；它是有条件的候选接入，\textbf{不作为当前可用或必然获批的前提}。不得读取或挪用另一个产品客户端的凭据来替代本客户端授权。认证失效、额度耗尽或路线不获支持时暂停并请求决定，不静默转成 API key 付费。\cite{v121-plan-quickstart}

预算值仅用于说明拟定的初始限额，启用前由负责人确认；费用构成与三档月度测算见第~\ref{sec:impl-budget}~节，不能把限额视为实际开销预测。月度预算控制须在自定义配置与预算管理器中另行实现并通过测试，本示例仅展示单任务和每日字段。预算管理器在调用前按输入、输出上限及已核验价格表预留额度，调用后记录实际用量和估算费用，无法确定账单的调用保留未结算额。订阅额度、API 货币费用及自有算力分别记录；订阅执行也受任务时长、并发和可见额度约束，不能把未返回货币金额记成无限免费。并行运行共享预留账本，不能各自认为余额充足；预留不足时暂停。模型最大输出量、超时和请求重试由适配层映射到锁定 SDK 支持的参数；\texttt{max\_turns} 不能替代费用控制。示例 600 秒仅约束编排运行，Codex 与计算执行器另有后端时限；实际截止时间取任务、尝试与后端截止时间中最早者，不能用较长的后端超时覆盖上层停止条件。租约仍须独立续期与校验。

模型 429 或临时 5xx 可在总超时内退避重试；预算同时计入重试，且不叠加多层无界重试。涉及工具副作用的运行超时不得从头自动重放。备用模型需先通过本项目契约测试，切换记录在运行中；测试与实验若要求模型一致，则暂停决策。API 密钥不进入提示词、Git、成果或普通日志；运行进程只读取所需凭证，默认不向训练子进程继承。

\subsubsection{组织规范、角色与项目上下文加载}

GitHub 私有仓库是不同执行环境共享程序与文档的版本核心。组织仓库维护入口、组织规则、岗位角色、skills 与项目索引；项目仓库维护项目约定、requirements、design、tests、源码及经审阅的交付文档。组织索引只导航到项目，不复制一套会与项目仓库失同步的需求或方案。Git 用提交记录可追踪修订，本平台另规定审阅、授权和冻结流程；采用 Git 并不会自动实施这些规则。\cite{ch5-github-git}

\par\noindent\begin{minipage}{\linewidth}
\begin{lstlisting}[caption={组织与项目仓库的拟定内容；真实地址及权限由实施时登记}]
organization/
  README.md              # organization entry and reading order
  organization.md        # organization rules and boundaries
  roles/{coordinator,developer,algorithm,test,ops}.md
  skills/<skill-name>/SKILL.md
  projects/index.yaml    # project ID, repository and entry paths

project-repository/
  README.md
  docs/requirements/     # scoped requirements and acceptance criteria
  docs/design/           # design and interface decisions
  tests/                 # synthetic fixtures and deterministic checks
  src/                   # implementation, where applicable
  docs/releases/         # reviewed delivery notes and artifact refs
\end{lstlisting}
\end{minipage}\par


岗位定义随组织仓库版本维护，Agent 实例按任务创建，不要求五个岗位同时在线。平台在领取任务后取得指定 commit 的只读快照，并把组织规范和角色文件映射到任务工作区的 \texttt{prompts/}；下例路径保持与平台模板一致。这里的文件是可重建的任务副本，权威修改应回到组织仓库审阅，不能在多个节点各自维护一套“最新提示词”。配置加载器先形成框架无关的角色、任务与上下文包，再由 Agents SDK 适配器构造 \texttt{Agent}，或由 Codex 适配器准备获准的工作区和任务输入。Agents SDK 不会自动读取本方案规定的组织入口、\texttt{AGENTS.md}、skills 或 YAML；Codex 自身的项目指令加载也须与固定版本清单核对，不能假定它等同于平台加载规则。

\par\noindent\begin{minipage}{\linewidth}
\begin{lstlisting}[caption={agents.yaml 的角色与上下文规则；均为平台拟建加载器字段}]
organization_prompt: prompts/organization.md
roles:
  coordinator: {prompt: prompts/coordinator.md,
    model_route: deep, policy: coordinate}
  developer: {prompt: prompts/developer.md,
    executor: code_dev, policy: implementation}
  algorithm: {prompt: prompts/algorithm.md,
    model_route: deep, policy: research}
  test: {prompt: prompts/test.md,
    model_route: standard, policy: independent_review}
  ops: {prompt: prompts/ops.md,
    model_route: standard, policy: operations}
context:
  require: [project_charter, task_contract,
            input_manifest, permission_snapshot, source_revisions]
  execution_requires: [frozen_requirements]
  skills_mode: task_selected
  history_mode: task_scoped
  revision_policy: pinned_commit
\end{lstlisting}
\end{minipage}\par


\texttt{source\_revisions} 是任务数据中的版本清单，至少记录组织仓库与项目仓库地址、完整 commit、所读文件路径及摘要，skills 一并绑定所属提交。\texttt{context\_builder} 按“组织入口与规范、角色职责、相关 skills、项目约定、冻结需求与方案、任务契约、必要证据”的顺序组装模型上下文，只加载该任务需要的材料。意图澄清阶段可尚无冻结需求，执行阶段则必须具备。任务开始后即使远端分支更新，当前输入也不自动改变；采用新提交须登记变更、影响范围和新的任务依据。

\texttt{executor: code\_dev} 引用平台适配映射；开发任务采用的 Codex 模型须由该适配器另行解析并记录，不继承协调器的 API key 或模型参数。其余角色中的 \texttt{model\_route} 用于当前编排路径。任务契约至少给出目标、排除项、依赖、输入成果修订、验收条件、输出类型和停止条件。岗位和 skills 提供做事方法，不能扩大权限；外部资料始终作为待分析数据，不能变成授权指令。研发协调者维护版本范围；发现新需求时提交变更建议或后续 backlog，不在实现过程中无声扩大范围。任务及运行记录另存 Mac 数据库，聊天正文、密钥和实验过程不随版本清单进入 Git。

自建平台由上述加载器自动装配；dot 云电脑则在获准的任务环境中读取相同仓库及明确 commit，按任务说明完成开发与多 Agent 协作，输出分支、差异及测试证据供审阅。二者共享版本依据，但各自维护执行上下文与状态，不能把读取同一仓库等同于共享会话或共享控制程序。

代码侧运行上下文另保存认证主体、项目范围、租约和已验证的工具客户端；这些值由程序注入，不能由模型提供的 \texttt{project\_id} 或声称的角色覆盖。SDK 的本地 run context 与模型可见对话属于不同通道，工具应使用前者核验身份与依赖。\cite{rdRuntimeAgent}

\subsubsection{工具接口、权限与执行校验}

起步只暴露少量业务工具：读取确定成果修订、在指定工作副本写代码、执行允许的测试入口、提交／查询训练作业、登记候选成果、提出子任务与提交测试证据。所有工具均返回 \texttt{operation\_id/status/result\_refs/error}；\texttt{status} 区分 \texttt{completed/failed/pending/unknown}，不把“请求已发出”返回为完成。

\Needspace{12\baselineskip}
\par\noindent\begin{minipage}{\linewidth}
\begin{lstlisting}[caption={policies.yaml 核心片段；由工具网关强制执行}]
default: deny
policies:
  coordinate: {allow: [task.read, task.propose, review.read,
                      artifact.read, artifact.propose]}
  implementation: {allow: [artifact.read, workspace.write,
                           test.self_run, artifact.propose]}
  research: {allow: [artifact.read, experiment.submit,
                     experiment.read, artifact.propose]}
  independent_review: {allow: [artifact.read, test.independent_run,
                               review.submit]}
  operations: {allow: [health.read, backup.verify, log.read]}
constraints:
  enforce_project_scope: true
  require_current_lease_for_write: true
  candidate_write_for_reviewer: false
  approval_required: [release.publish, dependency.install,
                      network.expand, data.delete]
\end{lstlisting}
\end{minipage}\par


\texttt{allow} 是能力上限，仍须叠加项目授权、输入数据范围、任务修订、当前租约及具体审批。协调角色的 \texttt{artifact.propose} 只向受信领域服务提交需求或方案候选，不能发布、冻结或覆盖已批准修订。路径规范化后检查是否位于获准工作区，拒绝符号链接逃逸；命令使用已登记的参数化入口，设置超时、输出限额与进程隔离。不为各 Agent 配置通用管理员 shell。Mac 控制端的虚拟环境只隔离依赖，不能作为恶意代码安全沙箱；执行不可信代码必须进入另行批准的受限环境。

开发执行器内部还可能运行自己的命令与工具，不能只在外层函数包装器检查一次便视为受控。Codex 适配器须落实工作区、只读输入、沙箱、网络、审批与日志映射；仅设置 \texttt{workspace-write} 不能证明新生成代码已隔离执行。须验证命令和测试实际进入获准的受限环境，默认不允许新生成代码在宿主机执行；未经过隔离验证的接口（包括 \texttt{thread/shellCommand}）不加入适配器白名单。不能落实的约束必须令任务停下。SDK 内部的批准也不能自行扩大平台授权。本文不使用已经移除的 \texttt{codex mcp-server} 作为接入方案，正式接入应选择当前支持的 SDK／app-server 或非交互接口。\cite{codex-sdk,codex-app}

下例仅说明当前 Agents SDK 适配器中的接口接入点，\texttt{RuntimeContext} 及 \texttt{read\_checked} 是\textbf{待实现的平台类型与校验服务}，不是 SDK 内置功能。执行端校验必须真实存在，提示词中的“遵守权限”不能替代它。\cite{rdRuntimeAgent}

\Needspace{14\baselineskip}
\par\noindent\begin{minipage}{\linewidth}
\begin{lstlisting}[language=Python,caption={SDK 接线片段；省略平台服务实现，不表示已运行测试}]
from agents import Agent, Runner, RunContextWrapper, function_tool

@function_tool
async def read_artifact(ctx: RunContextWrapper[RuntimeContext],
                        artifact_id: str, revision: int) -> str:
    """Read one authorized, immutable artifact revision."""
    return await ctx.context.artifacts.read_checked(
        principal=ctx.context.principal,
        project_id=ctx.context.project_id,
        artifact_id=artifact_id, revision=revision)

agent = Agent(name=role_name, model=resolved_model_id,
              instructions=assembled_prompt, tools=[read_artifact])
result = await Runner.run(agent, task_input, context=runtime_context,
                          max_turns=12)
\end{lstlisting}
\end{minipage}\par


\subsubsection{委派、并行、独立审查与返工}

跨步骤协作采用\textbf{协调 Agent 提案、平台校验后建任务、Worker 执行、成果驱动交接}。子任务带 \texttt{parent\_task\_id}、独立输入清单、验收条件和输出约束；调度器检查依赖、角色和预算后才允许领取。长训练与代码实现均走持久任务／作业接口，不能把一个内存中的并行协程视为跨进程可靠任务队列。

SDK 的 \texttt{agent.as\_tool(...)} 适合短小专业咨询，调用后协调者继续负责汇总；\texttt{handoffs} 表示同一 SDK 运行中转移后续处理权。二者均不自动创建本平台的持久任务、确认交接或实现独立测试；若使用，适配器仍需登记子运行、限额及证据。\cite{rdRuntimeOrchestration}

初始最多并行两个 Agent 运行、一个 4090 训练作业；待资源实测后调整。并行实现使用独立 Git 分支／工作副本和独立输出目录，同一输入快照只读复用。开发负责人安排合并并形成固定 commit，测试角色在新工作区获取该候选版本，不能复用开发者可变的工作目录。代码仓库属于平台成果管理；受控 Codex SDK 适配器可以参与平台内部自动开发，人工使用的 Codex／VS Code 则仍是外部入口。两条路径都提交候选版本与测试证据，均不能绕过任务数据库和独立验收。

测试任务使用独立运行与权限配置，依据冻结需求自行产生证据；实现者的自测记录仅作参考。\texttt{review.submit} 检查测试角色、候选修订和证据完整性，不能仅接受“通过”文本。发现缺陷而目标与责任范围不变时，原任务记录退回原因并以新尝试返工，修复形成新候选修订后复测；只有需要拆分责任或范围时才建立关联子任务，明确原任务的依赖。旧尝试、候选与审查记录全部保留，不得改写。研究结论为否定或未显著改善，只要方法、数据和论证满足验收标准，也可验收；不能让 Agent 为追求成功状态挑选有利结果。

\subsubsection{会话、检查点与上下文恢复}

恢复分成三个层次。第一层为 PostgreSQL 中的任务、尝试、作业、审批与工具操作账本，决定“哪些事确实发生”；第二层为\textbf{平台拥有的可迁移上下文包}，决定“下一次执行需要知道什么”；第三层为具体后端的会话、线程或中断检查点，用于同一兼容运行时继续。保存聊天、一个 OpenAI response ID 或 Codex thread ID，不能代替任务持久化，也不能恢复远程进程、工作目录和已经产生的副作用。

上下文包采用平台自己的版本化格式，至少保存目标与排除项、冻结需求、组织与项目 commit、输入及候选成果修订、已验证事实、未决假设、任务依赖、最新审批与停止条件、工具操作及测试证据引用、必要对话摘要。每条关键事实保留来源；摘要不丢弃原始证据，也不把模型推测升级为既成事实。此包不保存密钥、活连接或模型私有推理过程；原始 SDK 载荷另存，不混入通用结构。

\textbf{框架升级与切换按“事实核对、上下文重建、新运行接续”处理。}先停止旧代次的新副作用，核对已发出的外部操作，再基于当前任务修订和有效授权生成上下文包；新适配器装配输入并新建运行，保留与旧尝试和成果的来源关系。不能把一个框架的 RunState、另一个后端的 thread ID 或旧版本工具注册表直接导入新框架。更换编排组件不自动改变验收条件、任务状态或组织角色。

当前 Agents SDK 路径可在一次运行结束后用 \texttt{result.\allowbreak to\_input\_list()} 提取输入历史，按任务、后端与序号保存，作为重建材料之一；平台另保存可迁移摘要和证据。后续若采用 SDK \texttt{session}，同一会话不同时混用完整历史重放、\texttt{session}、\texttt{previous\_response\_id} 和 \texttt{conversation\_id}，以免重复上下文。这些 API 仅属于当前适配器实现，不成为平台公共契约。\cite{rdRuntimeResults,rdRuntimeRun}

当前 Python Agents SDK 的审批中断可通过 \texttt{result.\allowbreak interruptions} 和 \texttt{result.\allowbreak to\_state()} 取得中断状态。收到匹配原动作的有效决定后，适配器调用 \texttt{state.approve(...)} 或 \texttt{state.reject(...)} 处理对应项，再以 \texttt{Runner.run(agent, state)} 恢复原运行；不能批准所有中断，也不能另发“用户已同意”的文本替代审批账本。SDK guardrail 的位置存在边界，最终权限检查仍放在实际产生副作用的入口。\cite{rdRuntimeApproval}

\texttt{checkpoints} 拟保存运行／后端会话标识、序号、格式及 SDK／运行时版本、适配器版本、角色和工具注册表哈希、上下文包引用、状态载荷、SHA-256 与创建时间。数据库客户端、HTTP 客户端、密钥和进程对象不序列化；恢复时从可信服务重建身份、客户端、权限与租约。状态只保存在服务端，客户端提交决定而非修改后的运行状态；审批按动作及修订绑定，消费决定与推进运行采用原子转换。

\textbf{原始 SDK 检查点仅在同一受支持的兼容版本组合内恢复。}即使 SDK 提供保存与读取接口，跨 SDK、适配器、工具注册表或协议版本也默认视为不兼容，除非专门迁移与往返测试通过。当前 Agents SDK 提供 \texttt{state.to\_string()}、\texttt{RunState.from\_json(...)} 等状态保存／读取路径；适配器用 \texttt{context\_serializer} 限定数据字段，具体签名随锁定版本核验；本平台不依赖未测试的序列化字段保持永久不变。\cite{rdRuntimePersistent}Codex 线程继续同样依赖可用的后端会话存储、工作区和运行时；不能用“存在 thread ID”证明进程还在执行或工具没有重复执行。\cite{codex-sdk,v121-codex-exec}

上线门槛包括：审批中断后保存、销毁进程、用相同兼容版本恢复、拒绝或批准原动作并继续，验证工具不重复执行；另模拟框架升级不能读取原始状态，验证能够从平台上下文包建立新运行，保留任务、审批与成果引用。检查点只在明确的模型／工具边界或后端可恢复中断处保存，不承诺任意指令级断电续跑。

工具响应与检查点之间崩溃时，重启先按操作账本核对；发生过但未登记的外部结果进入 \texttt{reconcile}，核实后补记或按已批准的恢复方案继续。恢复还须检查取消、租约、配置兼容、输入校验值与凭证有效性；旧状态中的授权不自动延续至已变化的目标。敏感输入与 tracing 数据最小化，任务正文、密钥和私有文件不因调试而无条件外传。本章的执行能力仅用于获准研发、验证与运维，不设置实盘交易执行接口。




% v1.21 replacement for 5.5--5.7; proposed configuration, not deployed.
\subsection{开发执行器、云端服务与计算节点}
\label{sec:impl-nodes}

本节分别规定自有节点、托管开发环境和云端服务的职责。Mac 承担自建控制入口、主要状态数据库与受控 Codex 开发执行器；dot 云电脑用于文档与轻量开发及多 Agent 协作；阿里云承载既有服务并验证一般 Web 服务和轻量接口；RTX 4090 台式机承担数据与训练实验，IDC 的 T640 待网络恢复可靠后提供可选自托管模型服务。GitHub 连接这些环境的版本与审阅交接，不承担任务数据库或运行时文件系统。下列资源值是首轮验证的保守配置，不是设备性能测量；软件包、镜像、模块和命令中标明的拟建部分须在实施时完成，本文没有执行安装、联网配置或模型下载。

\subsubsection{Mac 与 Codex 开发执行器}

MacBook Pro M3 Max（64 GB 内存、2 TB 存储）作为首期个人研发工作台，运行原生 Python 3.12、uv、FastAPI、自建调度程序、Agents SDK 协调运行与 Codex 适配器；PostgreSQL 17 放在 Docker Desktop 的本地容器中，数据库与初始网页入口只绑定回环地址。平台源代码位于 \texttt{\textasciitilde/agent-platform}，状态与成果位于 \texttt{\textasciitilde/agent-platform-data}。平台按工作时段启动，不将合盖休眠的笔记本视为全天在线服务器。

\textbf{两个 SDK 分工明确，持久平台仍须自建。}Agents SDK 支撑研发协调、角色委派与工具调用；Codex SDK 封装开发执行，接收获准任务并在指定代码工作区内形成候选修改。官方当前 Python 包为 \texttt{openai-codex}，通过 JSON-RPC 控制本地 Codex app-server，发布包带有锁定的 Codex CLI 运行时依赖，适合与本平台 Python 主栈衔接。\cite{codex-sdk,rdRuntimeSdk} 项目、版本、队列、预算、权限、审批、成果和独立验收仍由自建平台维护；一次 Codex 会话成功结束不能直接关闭研发任务。

首期优先用 Python SDK 实现 \texttt{CodexExecutorAdapter}。若所需事件、审批或会话能力未由选定 SDK 版本暴露，可改为直接接入 app-server，按该版本生成的协议 schema 实现同一个适配器；这是两种接入方式，不是再建设一套平行任务平台。SDK 方式由 SDK 管理其本地运行时连接；直连方式使用私有 stdio 通道并完成初始化，再调用线程和轮次接口及读取事件。桌面内置 Codex、VS Code／Codex IDE 扩展是人工操作入口，安装或登录这些界面本身不构成自建平台的程序化接入。\cite{codex-sdk,codex-app}

\textbf{一次开发任务按以下顺序执行。}
\begin{enumerate}[leftmargin=2em,itemsep=4pt]
  \item \textbf{冻结要求并准备代码。}协调 Agent 形成与已批准版本需求对应的任务包；确定性调度程序校验需求修订、依赖、预算与权限，领取租约并固定组织／项目 commit。可信 Worker 从允许的仓库创建本次尝试独占的 worktree 或隔离克隆，登记基线与真实路径。只提供必要的 requirements、design、tests、相关 skills 和输入引用，不复制用户主目录；并发任务不得共享可写工作树；采用 Git worktree 时须验证公共 Git 元数据与 hooks 的访问边界，隔离不足则用独立克隆，可信 Worker 负责固定候选和版本操作。
  \item \textbf{启动受控 Codex 开发。}适配器以固定工作目录、已验证的运行时和权限配置启动或恢复 Codex，传入范围、禁止事项、验收条件与成果格式。会话／轮次标识映射到平台尝试；持续接收进度、文件变更、审批与终态事件，并登记事件序号和来源。普通实现细化可在范围内推进，新增需求、扩大目录或网络访问、改动验收标准必须作为变更请求暂停处理，不能藏在“修复”中。
  \item \textbf{交回可核验的候选成果。}Codex 返回修改摘要、差异、已执行检查、失败项和未完成项。可信 Worker 从实际文件重新提取 diff 与文件清单，固定候选提交或不可变补丁摘要，关联测试命令、环境、退出码和报告；没有运行的测试须明确列为未运行。自然语言报告与真实证据不一致时标记待核实。
  \item \textbf{独立测试与返工。}平台创建独立测试任务，从固定候选重新建立只读代码副本或测试工作区，用受控容器执行既定测试。测试角色可以使用 Agents SDK 或另一独立执行实例分析结果，但不能复用开发者的自评作验收。范围不变的失败返工保留原逻辑任务，追加新尝试和新候选，旧评审与证据不可改写；只有职责拆分才另建关联子任务。若需改变需求或验收标准，先走版本变更决定，不用放宽测试掩盖缺陷。
  \item \textbf{审阅、批准与交付。}协调 Agent 汇总证据供人审阅；调度程序检查独立测试及对应审批，通过后才由可信发布入口执行获准的合并／发布。批准绑定候选 commit／补丁摘要与发布目标；候选变化须重测并重新判断批准是否仍有效。拒绝或暂缓时保留候选与原因，不修改权威版本。
\end{enumerate}

\begin{figure}[H]
\centering
\begin{tikzpicture}[x=1cm,y=1cm,
  lanehead/.style={callhead,text width=2.57cm,minimum height=.82cm},
  lanebox/.style={callstep,text width=2.43cm,minimum height=.72cm,font=\scriptsize\sffamily,inner sep=4pt},
  lanetext/.style={font=\scriptsize\sffamily,align=center,fill=white,inner sep=2pt}]
\foreach \x/\name/\title in {0/c/{协调 Agent},3.05/p/{平台确定性调度},6.10/d/{Codex 执行器},9.15/t/{独立测试执行},12.20/h/{人类批准}} {
  \node[lanehead] (\name head) at (\x,0) {\title};
  \draw[calllife] (\x,-.48) -- (\x,-11.45);
}
\node[lanebox,draw=accentorange] (h0) at (12.20,-1.18) {批准版本需求与边界};
\node[lanebox] (c0) at (0,-1.18) {形成任务包\\要求／约束\\既定验收条件};
\draw[callreturn] (h0.west) -- (1.37,-1.18);
\node[lanetext] at (6.10,-1.18) {已批准需求修订};
\node[lanebox] (p0) at (3.05,-2.45) {核验权限与额度\\登记租约与基线\\独占 worktree};
\draw[callrequest] (c0.south) -- (0,-2.45) -- (p0.west);
\node[lanebox] (d0) at (6.10,-3.65) {按范围修改代码\\受控检查\\提出阻塞};
\draw[callrequest] (p0.south) -- (3.05,-3.65) -- (d0.west);
\node[lanebox] (p1) at (3.05,-4.91) {记录事件与审批\\核验差异\\冻结候选成果};
\draw[callevent] (d0.south) -- (6.10,-4.91) -- (p1.east);
\node[lanetext] at (7.88,-4.90) {diff、检查证据\\失败／未完成项};
\node[lanebox] (t0) at (9.15,-6.14) {独立任务与副本\\新容器／固定候选\\执行独立测试};
\draw[callrequest] (p1.south) -- (3.05,-6.14) -- (t0.west);
\node[lanetext] at (6.07,-6.14) {候选摘要＋既定验收集};
\node[lanebox] (c1) at (0,-7.45) {汇总差异与测试\\失败返工\\范围变更上报};
\draw[callreturn] (t0.south) -- (9.15,-7.45) -- (c1.east);
\node[lanetext] at (5.90,-7.45) {独立报告先登记，再提交审阅};
\node[lanebox,draw=accentorange] (h1) at (12.20,-8.70) {审阅候选与证据\\批准／拒绝／暂缓};
\draw[callrequest] (c1.south) -- (0,-8.70) -- (h1.west);
\node[lanetext] at (6.10,-8.70) {批准必须绑定版本及合并／发布目标};
\node[lanebox] (p2) at (3.05,-10.02) {复核批准再交付\\获准合并／发布\\登记不可变证据};
\draw[callreturn] (h1.south) -- (12.20,-10.02) -- (p2.east);
\node[lanetext] at (8.18,-10.02) {返回审批决定；仅有效批准后交付};
\node[lanetext,text width=14.7cm] at (6.10,-11.24) {平台全程追加任务事件与成果引用：\texttt{task / attempt / thread / turn / artifact}\\实线：执行请求；虚线：结果／决定；点线：运行事件。任一范围变化回到需求批准节点。};
\end{tikzpicture}
\caption{从已批准需求到 Codex 开发、独立测试与批准交付的泳道图。图中动作均为拟建流程；运行事件和模型判断不能自行成为审批。}
\label{fig:codex-development-swimlane}
\end{figure}

\textbf{适配器配置落在平台，而非提示词中。}下面是拟建 \texttt{workers.yaml} 中的执行器配置，所有字段均由平台自定义，不是可以直接传给 SDK 的官方参数；其中 \texttt{workers} 登记节点，\texttt{executors} 登记引用节点的执行器。配置加载器将它们映射到锁定 SDK／app-server 版本，并拒绝缺项、不支持的选项和未通过安全验收的配置。SDK 与 CLI 的确切版本、包哈希、协议 schema 和验证记录写入依赖锁定文件及执行器清单，不在方案中猜测版本号。

\par\noindent\begin{minipage}{\linewidth}
\begin{lstlisting}
# workers.yaml: proposed fields implemented by platform adapter.
executors:
  codex-mac:
    node_id: mac-local
    enabled: false                 # enable only after contract tests
    driver: codex_sdk             # alternative: codex_app_server
    runtime_manifest: config/locks/codex-runtime.json
    workspace_root: ~/agent-platform-data/workspaces/codex
    codex_home_ref: secrets/codex-runtime-home
    auth_profile_ref: secrets/codex-auth-profile
    auth_mode: codex_local_login
    paid_fallback: false
    policy_profile: codex-edit-only # defined in policies.yaml
    max_concurrency: 1
    turn_timeout_seconds: 1200
    task_attempt_limit: 3
    max_patch_bytes: 2097152
    host_generated_code: deny
    test_worker_id: mac-local
    test_capability: cpu_test
    event_sink: platform_attempt_events
    unknown_event: quarantine
    reconnect: reconcile_before_resume
\end{lstlisting}
\end{minipage}\par

\texttt{codex-runtime.json} 至少登记包版本、实际二进制版本及摘要、适配器提交、协议 schema 摘要与已验证功能；若显式覆盖 SDK 自带运行时，必须重新验收。\texttt{policies.yaml} 中的 \texttt{codex-edit-only} 配置要求限定可读／可写根目录、模型服务网络、审批桥接与允许命令；模型返回内容不能修改这些规则。平台调用接口只开放创建／恢复线程、运行／中断、读取状态与必要事件／审批处理，不对 Agent 暴露任意 app-server 方法。尤其排除会绕开线程沙箱的 \texttt{thread/shellCommand}。\cite{codex-app} \texttt{platform.adapters.code\_dev.executor\_id} 指向 \texttt{codex-mac}，其 driver 为 \texttt{codex\_sdk}，不得从前端绕过映射自行选执行器。首期并发一个 Codex 开发尝试，独立测试另占容器槽；1200 秒仅为 Codex 后端初始软期限，与协调运行的期限分别配置，实际期限取任务、尝试与后端截止时间中的最早值。超时先请求中断并核实进程终态，不能把发出取消视为已停止。

事件接收必须与执行并行，先持久保存标识与脱敏摘要，再推动状态；审批请求绑定任务、尝试、动作及目标，由平台显示或按已批准策略判断，未知请求默认阻塞。协议不能保证重放所有流式增量，断线后须读取线程与轮次状态、核对实际工作树及操作账本，再决定继续订阅／恢复或创建新尝试；不能仅重新提交同一提示词。适配器保存平台自己的接收序号和去重键，不把它伪装成服务端可恢复游标。连接关闭、调用返回、轮次终态与平台任务完成分别记录。SDK 无法暴露所需能力时保持该能力禁用，或通过同一适配器经契约测试后改用直连 app-server。

\textbf{认证、计费和宿主边界分别验证。}本人在本地 Codex 已完成的 ChatGPT 登录，可在确认账户、工作区、执行身份和该运行环境确实可复用后作为本地认证来源；界面已登录不保证后台进程具备相同凭据。订阅配额、允许的自动化方式与托管场景须按实际账户及官方条件核实，不能推断覆盖所有后台任务。API key 模式按 OpenAI Platform 另行计费，启用前须获准预算并显式选择；订阅不可用、限额耗尽或认证失效时进入等待，禁止静默切换付费 API。\cite{v121-codex-auth} 身份秘密只供执行器认证层使用，不进入协调 Agent、任务提示词、代码工作区或测试容器。

Codex 原生运行须使用隔离工作区与经过验证的权限配置；\texttt{workspace-write} 本身不等于禁止宿主执行任意代码，也不等于隔离所有可读文件。只允许经审查的静态文件操作，生成脚本、构建、依赖安装和测试一律由可信 Worker 放入一次性容器执行。实施必须证明所选运行时的命令／文件策略确实阻止宿主执行此类代码、读取秘密及越权访问；若做不到，则保持原生执行器禁用，将整个 Codex 执行进程置于另行验收的隔离运行环境，不以提示词代替安全边界。不得向 Codex 工作区或执行容器提供 Docker 套接字、用户主目录、SSH 密钥及平台配置；Docker 控制能力仅留在可信 Worker 内部，Agent 不能提交任意挂载、特权模式或容器参数。容器隔离仍需安全更新与最小权限。\cite{ops-docker-security}

Docker Desktop 首轮预留 6 个逻辑 CPU、16 GiB 内存；容器执行最多并发 2 个作业，每个上限 2 CPU、6 GiB、256 个进程，默认 10 分钟超时。数据库、缓存和 SDK 进程纳入宿主总预算，内存不足时停止领取任务。以下是\textbf{拟建可信 Worker 生成的测试约束}：变量由 Worker 解析至该作业目录，镜像预先审查并固定摘要；代码副本只读，需要编译或缓存时写入独立的临时／输出位置。

\par\noindent\begin{minipage}{\linewidth}
\begin{lstlisting}
docker run --rm --network none --read-only \
  --user 10001:10001 --cap-drop ALL \
  --security-opt no-new-privileges \
  --cpus 2 --memory 6g --pids-limit 256 \
  --tmpfs /tmp:rw,noexec,nosuid,size=512m \
  --mount type=bind,src="$JOB_CODE",dst=/workspace,readonly \
  --mount type=bind,src="$JOB_INPUT",dst=/input,readonly \
  --mount type=bind,src="$JOB_OUTPUT",dst=/output \
  --workdir /workspace "$VERIFIED_ARM64_IMAGE" \
  python -B -m pytest -q -p no:cacheprovider \
    --junitxml=/output/tests.xml
\end{lstlisting}
\end{minipage}\par

Worker 负责计时、终止与回收容器，输出目录预先赋予对应非特权用户写权限。每秒统计写入量，超过 2 GiB 即取消，这是可能短时越过的软阈值；需要硬限额时另验收文件系统或卷配额。依赖由独立受控步骤准备，测试阶段默认无网络。Codex 的自检与独立验收分别记录，后者从固定候选重新取得输入，不读取仍在修改的开发工作区。

\Needspace{8\baselineskip}
\subsubsection{dot 云电脑轻量开发与协作}

将 dot 云电脑作为轻量开发环境，适合整理和修改文档、LaTeX 编译排版、简单 Web 原型、轻量代码和测试，以及在任务内组织多个 Agent 分工研究、实现与检查。开始前确认该任务实际可用的工具、文件、依赖和访问权限；这是一项工作分工，不代表已安装本章全部依赖，也不把托管环境当成长期固定的自有主机。

其交接以 GitHub 版本为中心：获准读取组织入口与项目仓库，固定组织和项目 commit，按角色与相关 skills 完成有边界的工作；修改保留在任务分支或补丁中，附可复核的测试及生成文件清单，经审阅后进入权威版本。若该环境当次不能直接读写 GitHub，可导出完整源文件、补丁和证据，由获准的环境核对基线并提交，不能假装仓库已经同步。Mac 随后在固定提交上复核需要本地数据或权限的步骤。

不把 dot 云电脑自动列入 \texttt{workers.yaml}，不要求它领取自建数据库租约，也不假定它永久在线或提供可被外部调用的任务 API。若以后确有自动任务接入需求，应另行验证官方支持的接口、认证、生命周期与恢复行为，再决定是否集成。本章的 Agents SDK 协调程序、Codex 开发适配器及模型认证仍独立配置；dot 不是保证可供外部调度的 SDK Worker，dot 内部如何调度不能替代本平台的实现与验收证据。

云电脑会产生工作副本、临时文件和运行记录；仅放入本次任务获准处理的数据，不上传真实账户资料、密钥或不必要的过程数据库。交付时将经审阅的程序和文档纳入 GitHub，将需要保留但不宜进 Git 的产物与日志证据回收至本地，并登记来源和校验值。任务结束后按该环境实际支持的保留与清理机制处理，不能承诺云端零数据或未验证的自动清除。

\subsubsection{阿里云一般 Web 服务与轻量接口}

现有阿里云服务器配置为 2 vCPU、2 GB 内存、40 GB 存储，已运行 Web 服务和交易所网格程序。首期维持该规格，新增范围限定为一般 Web 服务及轻量任务接口验证；自建平台的主数据库、完整协调循环、大规模 Agent 并发、训练及模型推理不迁入该小实例。已有网格作为独立既有业务保留，不因新增研发项目而改动账户、订单、策略参数或启动配置。

实施前先盘点现有进程、端口、反向代理、开机启动、内存和磁盘余量，保留可恢复配置及发布版本。新增 Web 验证项目使用独立仓库、应用身份、目录、日志和进程入口；只发布 GitHub 已审阅且通过测试的 commit 或固定发布包，认证配置和秘密在节点单独注入。发布、停止、回退与访问控制均按具体部署任务授权执行，不能把本章的部署建议视为已完成部署。

最小接口只接收经验证的合成输入或任务引用，返回请求标识、接受状态和可查询结果；可以先用本地或容器模拟服务端，不要求立刻对公网开放。若后续需要与 Mac 控制 API 交接，须单独核实双方网络、认证、幂等与超时处理；Mac 休眠时接口只能诚实返回等待或不可用，不能假定它能继续推进本地队列。未实现、未认证的接口不计为平台 Worker，也不向其提供数据库账户。

容量采用“先监控、再决定”的方式。记录 CPU、可用内存、交换分区、磁盘和 I/O、请求耗时、错误率及既有服务的资源高峰；可优先复用阿里云主机监控和应用现有日志，不另建监控集群。\cite{ch5-aliyun-monitoring} 首轮不新增重型中间件。若出现反复 OOM、持续内存压力或换页，并且减少非必要进程、控制并发和日志轮转后仍不能满足既定 Web 验收目标，再比较 2 核 4 GB 的成本与效果并提交决定。CPU、磁盘或网络瓶颈不能用仅增加内存来解释或掩盖，规格变更也不自动保证性能。

阿里云服务必需的当前状态、访问与错误日志在该节点保存，并设置轮转和空间限制；按获准清单把必要证据、配置备份和归档拉回本地。40 GB 磁盘应为系统、既有业务、回退包和日志留出余量，不上传训练数据或模型权重。固定参数网格的后续调整独立于新增 Web 验证，RSI 简单策略仅作为另立项目的模拟流程候选；两者均需自己的 requirements、design、tests，本章不扩展实盘交易范围。

\subsubsection{RTX 4090 训练节点}

台式机已知为 i9、RTX 4090、1 TB 存储；内存是 64 GB 还是 128 GB、当前操作系统及驱动尚须盘点。优先准备 Linux x86\_64 训练环境；若保留 Windows，则先验证 WSL2、GPU 访问、文件路径与休眠行为，不据本文重装系统。首期只运行一个 GPU 训练作业，系统按 64 GB 内存下限安排：单作业内存上限 32 GiB、CPU 上限 8 个逻辑核，至少为宿主保留 16 GiB 内存和 2 个逻辑核；实测资源不足时降低配额或拒绝领取任务，不假定 i9 的具体型号与核数。

训练栈采用固定版本的 PyTorch 与匹配的 CUDA 运行环境，宿主另需兼容 NVIDIA 驱动；采用容器时再配置 NVIDIA Container Toolkit。先按官方安装矩阵确认操作系统、驱动、框架和 GPU 的组合，再固定镜像摘要及依赖锁定文件，不将宿主 \texttt{nvcc} 版本、驱动支持的 CUDA 版本与 Python 包携带的运行库混为一谈。\cite{ops-pytorch-install,ops-nvidia-container}

Mac 的 Apple Silicon/ARM64 开发环境与 4090 的 x86\_64/CUDA 环境分别构建、分别验收；共享的是源码提交、数据清单和参数，不是复制虚拟环境目录或直接复用二进制镜像。平台仓库保留基础依赖与各平台锁定结果，训练项目另保存自己的环境清单。Worker 登记 Python、PyTorch、驱动、CUDA 运行库、GPU 名称及镜像摘要，并在领取作业前完成 GPU 枚举和小张量计算检查。训练记录保存随机种子、划分方法、超参数、输入摘要与检查点；相同种子不保证跨平台逐位一致。\cite{pytorch-repro}

\noindent\begin{minipage}{\linewidth}
\begin{lstlisting}
# workers.yaml: proposed limits; hardware discovery is mandatory
workers:
  rtx4090:
    enabled: false          # enable after inventory and acceptance
    execution_kind: training
    platform: linux/amd64
    gpu_count: 1
    max_concurrency: 1
    cpu_limit: 8
    memory_limit_gib: 32
    host_memory_reserve_gib: 16
    workspace_soft_limit_gib: 100
    required_checks: [driver, torch_cuda, disk, input_hash]
\end{lstlisting}
\end{minipage}


显存限额由模型、批大小、序列长度及框架实际占用共同决定，不能仅靠容器内存上限约束。先以小批量试跑并记录峰值，再决定正式训练参数；显存不足时保留失败与检查点，由任务计划决定降批量或重跑。1 TB 本地盘保存获准数据、当前输入缓存、实验记录、检查点与未完成交接的成果；长期保留内容按清单归档至本地受控存储，模型与数据累计达到磁盘预警线后不再派发新作业。

\subsubsection{跨节点任务与数据传递}

本小节所述远程节点是已接入自建协议的 4090 等 Worker；dot 通过仓库交接、阿里云通过独立 Web 接口验证，不自动套用该协议。远程 Worker 主动领取允许的作业，经已授权配置的私网 HTTPS 接口报告状态；初期 Mac 的回环接口不能被远端直接访问。启用远程执行前须完成私网可达性、服务证书、节点身份和访问规则验证，数据库始终不向 Worker 开放。长训练对应独立计算子任务，拥有自己的尝试、租约和作业标识，不复用 Mac 上父 Agent 的推理租约；父任务保存检查点后进入 \texttt{waiting}，结束当前尝试并释放推理槽和租约；计算成果核验后唤醒父任务为 \texttt{ready}，由新尝试引用检查点继续。任务包携带代码提交、环境摘要、输入成果标识及校验值、参数、资源要求和输出规范；文件通过受控传输接口获取，不信任任意远程路径。

\begin{enumerate}[leftmargin=2em,itemsep=3pt]
  \item \textbf{领取与本地落盘。}Worker 以作业与尝试标识领取租约，先将输入传入临时文件，校验长度、摘要及路径后再原子改名。完整输入与执行计划落盘后才启动计算；重启时先查本地作业日志，避免重复启动同一尝试。
  \item \textbf{断线与续传。}网络中断时，仅预先批准可离线继续、输入齐全且已启动的有界计算进程可在原限额内运行，写入本地隔离输出；不得新启进程或作业、对外写入或发布。状态与结果进入持久化待发送队列。Mac 休眠期间不领取新任务；控制端把失联作业标为待核实，不立即派发副本。恢复连接后先提交观察记录，对账同一作业的进程、检查点、取消与终态，再决定采纳输出或经重新校验建立新尝试。
  \item \textbf{成果接收。}输出先写入临时区域，再上传文件、长度与 SHA-256 清单。接收端拒绝越界路径、符号链接逃逸和不符的摘要；文件校验完成后才发布成果目录并提交数据库登记。若发布成功而登记失败，按唯一标识补登记；未完成的上传不计为成果。
\end{enumerate}

回执重发必须幂等；租约过期后的终态报告不得直接走有效租约的完成接口，而应提交到 \texttt{/api/v1/attempts/\{a\}/observations} 对账入口。迟到结果先归档到原尝试，由当前控制程序核对取消状态、当前代次、待处理子任务、实际进程及输入与成果后决定是否采纳，不能覆盖当前任务或跳过独立验收。远端在收到成果接收确认前保留本地输出；拟设 7 天未确认即告警，达到缓存空间阈值则停止新作业，不自动删除唯一副本。检查点恢复由训练程序明确支持，平台不能把任意进程从中断指令处恢复。

\subsubsection{T640 自托管模型服务}

IDC 的 T640 已知有 128 GB 内存、2 TB 存储及 Tesla V100S 32 GB；CPU、操作系统、驱动与磁盘可用空间须现场核验。当前网络较差且常离线，故首期 \texttt{enabled: false}，不承担控制数据库、必须实时响应的模型调用或训练调度。其进入条件是先修复网络并完成私网访问验证，再建设可被 Mac 调用的自托管模型服务；4090 仍是训练节点。

\textbf{V100 的软件兼容性须单独处理。}NVIDIA 官方表将 V100 列为计算能力 7.0，本次核验的 vLLM 最新 NVIDIA 安装要求为 7.5 及以上，因此不能直接以最新版 vLLM 作为该节点基线。CUDA 13.0 已移除 Volta 的离线编译与库支持；“更新到最新 CUDA”在此不是有效方案。\cite{ops-nvidia-legacy,ops-vllm-gpu,ops-cuda13}

首轮候选为固定提交的 \texttt{llama.cpp} CUDA 后端：其官方构建文件仍在 CUDA 13 以前的工具链路径保留 \texttt{sm\_70}。实施时选择经验证支持 Volta 的 CUDA 12.x 小版本、匹配驱动和引擎提交，以显式架构编译后实机验证；支持编译不等于所有模型、量化格式和内核都能运行。\cite{ops-llama-cmake,ops-llama-build} 下列仅为\textbf{依赖和模型已准备后的候选构建与启动形式}，路径、提交、工具链小版本与模型文件均须先写入环境清单，不包含自动下载。

\par\noindent\begin{minipage}{\linewidth}
\begin{lstlisting}
cmake -S . -B build-v100 -DGGML_CUDA=ON \
  -DCMAKE_CUDA_ARCHITECTURES=70 \
  -DCMAKE_CUDA_COMPILER="$VERIFIED_CUDA12_NVCC"
cmake --build build-v100 --config Release -j 4
./build-v100/bin/llama-server \
  --model /srv/models/verified-model.gguf \
  --host 127.0.0.1 --port 8081 --ctx-size 4096 \
  --parallel 1 --n-gpu-layers 1 \
  --api-key-file /run/secrets/model-service-key
\end{lstlisting}
\end{minipage}\par


先用单并发、短上下文和少量 GPU 层确认加载与推理，再逐步增加 GPU 层数，测量显存、首字延迟、完整响应时间和连续运行稳定性；示例参数不代表推荐性能。原始服务只绑定节点回环地址，由另行批准配置的私网 HTTPS 入口转发，并限制来源和请求大小。健康接口仅证明进程或模型就绪，验收还必须发出真实短请求并检查结果。\cite{ops-llama-server}

模型名称、规模和量化方式须以实际任务质量、许可、显存与上下文试验决定，本版不预选一个“32 GB 一定能跑”的模型。第二块 GPU 仅保留扩展可能，不形成采购或部署承诺；须先核对 T640 的插槽、电源、散热、连接方式及引擎并行支持，两张 32 GB 卡不会自动成为无缝的 64 GB 显存。

\subsubsection{OpenAI 与自托管模型的接入关系}

协调 Agent 的当前主链路使用 OpenAI API，Codex 开发执行器按自己的已验证认证配置运行；T640 通过验收后作为新增模型后端，先用于适配过的摘要、分类或辅助编码任务。模型注册表记录 \texttt{provider}、\texttt{endpoint}、\texttt{model\_id}、\texttt{api\_shape}、认证引用、超时、上下文上限与能力测试结果。OpenAI 使用既定的 Responses 适配；自托管后端先用 Agents SDK 的独立 Chat Completions 适配，不能只替换默认 Responses 客户端的地址，但须分别验证文本、结构化输出、工具调用、多轮上下文和流式返回。接口形式相似并不保证行为或功能等价。该模型适配不自动把 T640 变成 Codex 后端，Codex 的其他供应商支持须单独验证。\cite{agents-models,ops-llama-server}

本地模型的工具调用只是向平台提出请求，仍经过同一授权和执行检查；推理服务器不直接取得研发工具或训练节点权限。自托管端离线时仅对允许出域且预先允许切换的任务回退到 OpenAI，记录切换理由及实际模型；有本地处理限制、比较试验或固定模型要求的任务进入等待，不通过回退绕过限制。T640 未启用期间，这些任务应明确显示当前缺少可用后端。

\subsection{数据、安全与运行维护配置}
\label{sec:impl-operations}

\subsubsection{代码、数据、配置与成果的存放规则}

平台数据库中的项目、版本、任务、尝试、审批与成果构成持久业务记录。Agents SDK 的运行检查点和 Codex 的 thread／turn／item 属于执行器状态，前者记录协调运行的可恢复位置，后者保留代码执行会话；两者通过平台映射记录关联，不能代替任务表，也不能凭会话被删除或压缩而删除任务历史。一个任务可经历多个执行尝试、多个 SDK 运行或 Codex 会话；取消、重试、返工和更换执行器均追加关系及原因，不覆盖旧尝试。

GitHub 私有仓库是程序与文档的权威版本核心，不只是代码备份。组织仓库维护组织入口、角色、skills 和项目索引；\texttt{agent-platform} 维护平台程序、配置模板、迁移与测试；每个项目仓库维护自己的 requirements、design、tests、实现和交付说明。项目目录可有本地或云端工作副本，但审阅后提交才构成共同基线，不能各自维护没有关联的“最终版”。修改走分支、差异审阅、检查和合并，任务绑定完整 commit；新版本不静默替换已运行任务的输入。

密钥、令牌、敏感账户资料、真实业务数据、过程数据库、对话正文、原始运行日志、模型权重与大型实验产物不入 Git。可提交合成测试数据和经检查不含敏感正文的验收摘要，完整过程证据按受控引用保存。\texttt{.gitignore} 只是第一道排除规则，提交前仍须检查差异、秘密标记和数据分类；仓库为私有并不解除数据上传授权和许可限制。

\begingroup\small\renewcommand{\arraystretch}{1.15}
\begin{longtable}{P{3.0cm}P{8.7cm}P{3.0cm}}
\toprule
\textbf{位置} & \textbf{内容与写入规则} & \textbf{恢复依据}\\
\midrule\endfirsthead
\toprule\textbf{位置} & \textbf{内容与写入规则（续）} & \textbf{恢复依据}\\\midrule\endhead
GitHub 私有仓库 & 组织与项目程序、文档、skills、配置模板、合成测试；经审阅提交。密钥及过程数据排除。 & 仓库地址、commit、锁定文件与本地克隆。\\
\rowcolor{lightgray}
Mac 工作区与文件 & \texttt{agent-platform/} 及项目工作副本；\texttt{agent-platform-data/} 中的 \texttt{inputs/}、\allowbreak\texttt{workspaces/}、\allowbreak\texttt{artifacts/}、\allowbreak\texttt{logs/}、\allowbreak\texttt{backups/} 分开。 & 固定提交、文件清单、摘要及成果登记。\\
Mac PostgreSQL 卷 & 项目、版本、任务、作业、消息、审批及成果元数据；保存平台尝试与 SDK／Codex 标识的映射及去重事件。 & 数据库备份、迁移版本及关联文件快照。\\
\rowcolor{lightgray}
执行器受控状态目录 & Agents SDK 检查点、Codex 本地会话与适配器事件缓存；身份秘密分开保护，任务正文按分级保留。 & 运行时版本、兼容性校验与平台映射；不独立决定任务完成。\\
4090 本地磁盘 & 数据、实验记录、输入缓存、检查点和待交接产物；按作业隔离，长期保留内容归档至本地受控存储。 & 作业清单、输入摘要、环境和检查点。\\
\rowcolor{lightgray}
dot 云电脑 & 获准代码与文档的工作副本、临时文件及必要记录；交付后将版本成果入库，需留存的非 Git 证据回收本地。 & 仓库 commit、导出清单与已核实的本地副本。\\
阿里云服务目录 & 固定发布包、当前必要状态与轮转日志；与既有 Web／网格业务隔离。按清单本地归档，不作为主数据库。 & 发布版本、节点配置备份及本地归档。\\
\rowcolor{lightgray}
各节点秘密位置 & 本机真实配置与凭据；与工作区隔离，只引用秘密名称。Mac 使用 \texttt{.config/agent-platform/}。 & 单独保护的备份或重新注入。\\
\bottomrule
\caption{程序文档、主要本地过程数据与各节点运行状态分别保存。}\\
\end{longtable}
\endgroup

\textbf{本地优先不等于云端零数据。}Mac 是主要状态与归档节点，4090 是主要数据和实验节点；dot、阿里云及可选模型节点为完成任务仍可能保存必要工作文件、状态、缓存和日志。每类数据须明确产生节点、保留期限、获准传输范围、归档目的地和删除条件。本地归档在校验并确认接收前不能代替原节点的唯一副本；云端保存行为还须遵守服务实际条件。

每个成果清单至少包含任务与尝试标识；\texttt{job\_id} 仅在关联计算作业时填写，代码成果可为空；纯计算尝试没有 SDK 运行时，\texttt{run\_id} 可为空。清单还记录组织与项目提交、输入与输出摘要、模型及环境标识、生成节点、存放位置、归档状态和验收状态。Codex 候选另记录基线与候选提交、补丁摘要、SDK／CLI／适配器版本、会话与轮次标识、权限策略修订，以及自检与独立测试证据的各自来源。文件名不能替代不可变成果标识。临时文件、失败输出、候选成果和已验收成果采用不同状态；清理程序只删除符合策略且已有所需副本的临时数据，不覆盖已验收修订。

\subsubsection{身份、凭据与访问权限}

本地入口先服务单一负责人，回环绑定仍需会话认证、跨站请求保护及明确的审批会话。应用、数据库、代码仓库和各远程 Worker 使用不同身份；dot 任务与阿里云部署只取得各自所需仓库和项目权限，不因共享 GitHub 就取得本地数据、交易账户或跨节点管理权限。Worker 凭据只能领取分配给该节点的作业、读取对应输入及上报结果，不能更改全局策略或审批自己的请求。数据库迁移账户与日常应用账户分开；模型密钥只交给相应模型适配进程；Codex 登录状态只由执行器认证层使用，不转交协调 LLM、代码工作区或测试容器。两个 SDK 不共享可被模型读取的秘密文件，实际认证资料及平台数据库凭据不进入任务载荷。

真实秘密由负责人通过安全方式注入受限文件或进程环境，权限按本机账户收紧；本文仅给出变量名与空值模板，不生成密钥、不要求在聊天中提供秘密。\texttt{.env.example} 可进仓库，真实 \texttt{runtime.env} 不可进仓库；禁止在异常堆栈、进程参数、报告或成果清单中打印其值。凭据泄露时暂停相应调用、由人撤销并重新注入，再核对受影响记录。

\Needspace{10\baselineskip}
\par\noindent\begin{minipage}{\linewidth}
\begin{lstlisting}
# .env.example: names only; real values are injected separately
OPENAI_API_KEY=
DATABASE_URL=
PLATFORM_SESSION_SECRET=
PLATFORM_WORKER_CREDENTIAL_FILE=
SELF_HOSTED_API_KEY=
CODEX_AUTH_PROFILE_REF=
CODEX_RUNTIME_HOME_REF=
OPENAI_AGENTS_DISABLE_TRACING=1
OPENAI_AGENTS_TRACE_INCLUDE_SENSITIVE_DATA=0
\end{lstlisting}
\end{minipage}\par


以上 Codex 两项是平台自定义的受控引用，不是官方 SDK 自动识别的环境变量，也不承载令牌值。执行器启动器以允许列表清理继承环境，隔离平台数据库、GitHub 写入与发布凭据；即使同在 Mac，也不把控制进程的全部环境传给 app-server。需交互登录时由负责人完成官方认证流程，不复制秘密到会话或仓库。

策略执行采用默认拒绝：文件访问按项目和作业目录限制，Shell 命令按已审查模板执行，网络按目的地与用途开放；安装软件、修改网络、扩大目录访问、上传受限资料及正式对外发送等操作依类别进入审批。策略由程序加载并校验，Agent 的自然语言同意不能修改策略。审批记录绑定动作、对象、数据范围和有效条件；对象改变后重新判断。审计事件采用追加写入和受限维护权限，日常运行身份不能删除审批及关键历史。终端数据与备份应使用经确认的静态加密保护，实施前核对已有状态，本方案不直接改变加密或安全设置。

\subsubsection{模型请求、日志与追踪的数据范围}

发送给模型的内容由任务需要和资料分级共同决定：可外发资料可裁剪后进入 OpenAI 请求；仅限自有节点处理的资料只能交给已验收的自托管后端；禁止外传的秘密不进入任一模型上下文。文件名、完整路径、错误信息和工具结果也可能含敏感数据，不能只检查用户输入。默认传递必要片段和成果引用，不自动上传整库、整盘或全部对话。

Agents SDK 默认开启追踪，并默认允许追踪记录包含模型及工具输入输出。首期在进程启动前关闭默认云端追踪，同时显式关闭敏感内容采集，只保留平台自身的本地脱敏事件记录；以后若启用追踪出口，须单独确认数据范围并在发送前脱敏。仅增加一个追踪处理器不代表默认导出器被移除，运行中关闭追踪也不能撤回已经缓冲或发送的数据。\cite{ops-agents-tracing}

OpenAI API 默认不使用 API 数据训练模型，除非客户主动选择共享；这不等于零保留。官方说明滥用监测日志默认可保留最长 30 天，特殊情形可能更久，部分功能还保存应用状态。首期在平台自建、支持该选项的 API 调用中显式设置 \texttt{store=false}，不据此承诺消除所有保存；是否具有零数据保留资格、具体模型和端点限制，须以账户实际条件及当期文档核对。\cite{ops-openai-data}

Agents SDK 的 tracing 开关只约束对应 SDK；不能据此宣称已关闭 Codex／app-server 的会话记录、诊断或上游留存。Codex 认证方式对应的服务政策与工作区控制另核对，执行器本地目录另定权限、保留与备份范围；订阅登录不能套用 API 账户的数据政策。\cite{v121-codex-auth}

自建平台及各服务节点的日志默认只记时间、对象标识、状态、耗时、用量、错误分类及摘要，不记原始提示词、工具正文或凭据；dot 云环境仅能对任务可控制的文件和输出采用该规则，不能据此承诺改变托管服务底层留存。必要云端服务日志在节点轮转，需保留的证据按清单归档回本地。初始保留建议为运行诊断日志 30 天、审批和状态事件 180 天，交付证据随版本归档；需要正文排障时另开有时限的受控记录，并在工单中写明范围和清理条件。这些周期是平台建议配置，不覆盖资料本身的许可或用户另定要求。

\subsubsection{监控、预算与异常告警}

首期用结构化日志、数据库查询和轻量状态页即可，不先建设监控集群。状态页区分控制服务可用、Worker 在线、作业执行状态和模型可调用；“在线”不代表正在执行的结果正确。按任务关联记录队列等待、运行时间、模型请求数、令牌用量、预计费用、重试次数、GPU/内存峰值、磁盘占用及最近成功备份。协调服务存活、调度器推进、Codex 连接状态、活跃轮次、待审批和独立测试结果分别显示；有流式输出并不证明任务在有效推进，轮次完成也不证明测试通过。

\begin{itemize}[leftmargin=2em,itemsep=3pt]
  \item \textbf{健康与容量。}活动时段内 Worker 连续 3 个心跳周期失联即提示待核实；磁盘剩余低于 20\% 预警，低于 10\% 停止领取新作业；任务输出软阈值或已验证硬配额另行生效。GPU 内存不足、驱动异常和输入校验失败立即记录，禁止无界重试。
  \item \textbf{预算。}按任务汇总协调 Agent 的 OpenAI API 费、Codex 的订阅配额或 API 费、计算资源和其他服务费用，并分别保留计量来源；订阅配额不能换算成免费的通用 API 额度，同一笔用量也不因 SDK 层级重复计费。日预算、单任务预算及价格表修订由人设置；未配置时不启动付费验收。可计量调用先预留再结算，重试消耗同一预算；Codex 限制同时包括尝试数、运行时长、并发及可核验用量，不能仅靠提示词控制费用。达到 80\% 提醒、达到上限停止新调用；延迟／未知用量保守占额，已有在途调用可能继续计费，另留安全余量。订阅受限时等待或请人明确选择已获准的替代方式，禁止静默转 API。
  \item \textbf{异常通知。}只对需行动的失败、待审批、预算阻塞、无备份及持续失联告警；同一根因合并通知。默认在本地状态页显示，外部通知渠道启用前另定接收人、信息范围和发送许可，不自动发送含数据正文的消息。
\end{itemize}

阿里云另关联主机资源与一般 Web 的延迟、错误率、磁盘增长和重启／OOM 记录，扩容判定见第 5.5 节；现有网格监控属于既有业务边界，不把策略收益或交易成功作为平台验收指标。dot 任务的实际可用性按任务状态检查，不以自建 Worker 心跳衡量。

休眠时段不把 Mac 计划离线当作故障。恢复启动时集中展示休眠期间的远端结果、失败和待决项；不能声称本机休眠期间仍具备实时监控或告警保证。

\subsubsection{备份、升级、休眠与恢复}

备份按“数据库加文件快照加清单”形成一套恢复点。首期在工作日结束和升级前暂停新任务与成果提交，等待正在提交的成果完成，再执行 PostgreSQL 逻辑备份并复制对应的输入与成果清单；不直接复制仍在写入的数据库卷。\texttt{pg\_dump} 提供一致性导出，但数据库引用与外部文件仍需应用层协调。\cite{ops-postgres-backup} 归档与备份以本地受控存储为主，须覆盖 Mac 状态库、4090 实验清单及获准回收的云端服务证据。备份目录不能仅在同一块内部磁盘；应选定经授权的独立本地介质并确认加密保护，目的地未落实前不宣称已具备灾难恢复能力。GitHub 的提交历史不能恢复不入 Git 的数据库和过程文件；如以后采用异地云备份，须另定数据范围和授权。

拟保留最近 7 份日备份和 4 份周备份，已验收交付快照另行归档；清理前核对可恢复副本。首期恢复目标为最多损失一个工作日已记录的数据、在环境可用后 2 小时内恢复基本入口，均须通过恢复演练验证，不作为已实现的服务承诺。每月及重大升级前，把备份恢复到独立测试数据库，抽查成果摘要、任务状态与权限，并留下实测耗时。恢复演练必须沿原任务打开不可变成果链接，验证引用、版本、文件 SHA-256 和读取权限仍对应原证据；只有数据库行恢复而附件缺失时不得标为恢复成功。

Codex 的本地会话目录需要在受控停写条件下纳入相匹配的恢复点；认证秘密另行注入，不能把包含令牌的整个运行目录打包为普通成果。SDK 会话恢复仅在版本、目录和权限兼容时尝试，恢复失败仍保留历史差异、测试报告及平台任务。升级或替换 Codex 适配器时保留旧标识、旧事件和旧成果，通过新的适配器运行重新关联的新尝试，不伪造旧线程已经续跑。

升级先固定代码、依赖、镜像与迁移版本，备份后在测试副本验证，再切换；数据库破坏性迁移不能只靠回退代码恢复。Mac 休眠前停止领取任务，落盘协调检查点、执行器标识和工作树清单，核实 Codex 是否已中断，并明确远端作业是否可继续；恢复后先检查时钟、数据库、原进程／轮次、租约与未回传结果，再开放调度。处于未知状态的尝试进入对账，不立即启动副本；原执行器仍活跃时不得同时恢复第二个写入者。恢复连接不自动批准继续宿主副作用，Codex 的下一轮或新操作必须重新核对当前权限、取消状态与批准条件；离线计算仅按第 5.5 节限定的已启动有界进程执行。T640 模型与 GPU 驱动分别记录版本，升级前保留上一套可启动的引擎及模型清单，各节点分别验证升级；dot 输出经仓库审阅交接，阿里云新增服务保留原发布包和回退路径，不因研发平台升级而重启或改变既有网格。依赖维护记录漏洞、版本支持期与升级责任；断电演练核对数据库恢复和未完成文件清理，未配置受验证的备用电源时不承诺断电继续运行。

\subsection{平台验收与分阶段交付}
\label{sec:impl-acceptance}

\subsubsection{最小测试任务与验收环境}

首个验收任务采用通用的“小型 CSV 汇总工具”，配套一个最小 Web 页面及接口验证研发平台闭环：读取合成 CSV，按组输出数量、合计与均值，提交代码、确定性结果、测试报告和简短说明。固定样本为 A 组数值 10、20、30，B 组 5、15；期望 A 为 3/60/20，B 为 2/20/10。异常样本另含缺列、非数值、空文件及重复提交；具体错误处理写入这次测试需求，不借此定义任何业务项目需求。

验收先分两层：模拟模型与临时数据库完成不产生外部费用的程序测试；通过后，用已获准预算的真实 OpenAI 请求跑一次端到端任务。协调 Agent 从已批准需求形成任务，Codex 适配器在固定基线上产生实现与证据，独立测试角色在固定候选与只读样本上另行执行测试并核对输出，协调程序只有在成果及验收记录齐备后才能关闭任务。模型措辞可变化，结构化数据、权限检查和状态转换必须可确定验证。

Web 验证另立项目需求：页面提交合成 CSV 或固定样例标识，接口返回任务／请求标识，状态页展示等待、成功、失败及确定结果；重复提交不得产生重复任务，越权和错误输入须明确拒绝。先在受控本地环境或模拟阿里云目标完成“需求、开发、审阅、发布包、接口调用、结果、回退”闭环，可由 dot 生成轻量代码并经 GitHub 审阅交接。该模拟通过不等于阿里云已部署，真实节点验证须在盘点既有业务、核准操作和回退条件后单独记录；全程不连接真实交易账户或提交订单。

环境证据包含实际执行节点及操作系统／CPU 架构、Python 与依赖锁定摘要、PostgreSQL 版本、组织和项目代码提交、配置修订和所用模型标识；另记录 Codex SDK／CLI／适配器版本、认证类型（不含秘密）、权限配置、基线与候选摘要，以及平台尝试到 SDK 运行／Codex 线程与轮次的映射。Mac 的核心验收不依赖 dot、阿里云、4090 或 T640 在线；训练与自托管推理各有独立扩展验收，不能用一次云端模型任务通过来代表 GPU 节点已通过。

\Needspace{12\baselineskip}
\subsubsection{功能与集成测试}

\begin{table}[H]
\centering\small\renewcommand{\arraystretch}{1.18}
\begin{tabularx}{\textwidth}{P{1.6cm}P{3.2cm}L}
\toprule
\textbf{类别} & \textbf{验证对象} & \textbf{通过条件与证据}\\
\midrule
功能 & 配置校验、状态转换、任务依赖、权限、预算、成果登记。 & 必填项缺失即拒绝启动相关能力；非法转换和越权请求均被拒绝；同一提交键重放不新增作业；CSV 数值全部符合预期。\\
\rowcolor{lightgray}
集成 & 页面/API、协调运行、模型适配、Worker、数据库、文件成果。 & 同一任务标识贯穿全链路；数据库记录与成果摘要一致；测试独立运行并留下退出码和报告；无实际成果时不能显示已交付。\\
开发执行契约 & 固定任务、Codex 适配、候选与证据回传。 & 相同输入定位同一基线；会话事件可关联尝试；实际 diff、退出码、报告、失败／未运行项齐备；只有“完成”文字不能交付。\\
\rowcolor{lightgray}
独立审阅 & 独立测试任务、候选冻结与人工批准。 & 开发自检不能满足验收；测试只读取固定候选；候选或发布目标改变后旧批准不可误用；未批准不得合并／发布。\\
版本交接 & GitHub 组织／项目基线、dot 产出与审阅。 & 任务可定位组织和项目 commit；角色及 skills 来源可追溯；候选差异与测试证据齐备；未审阅修改不成为交付基线。\\
\rowcolor{lightgray}
Web 闭环 & 合成输入、请求状态、发布包、目标服务与回退。 & 页面到接口到结果可复核；未授权请求被拒绝；记录模拟／真实环境；真实验证不影响既有服务，未部署时明确标注。\\
训练扩展 & 4090 的领取、GPU 执行、检查点与回传。 & 使用合成数据完成小型训练；日志确认使用指定 GPU、指标为有限数、检查点可加载；中断后按程序支持方式续跑并只登记一次成果。\\
\rowcolor{lightgray}
模型扩展 & T640 的加载、请求、认证与能力适配。 & 网络进入条件通过；短请求、超时、未认证拒绝及已声明能力逐项测试；记录延迟、显存与质量结果后再决定可承担任务。\\
\bottomrule
\end{tabularx}
\caption{功能和集成验收；通过条件是待执行标准，不是已有测试结果。}
\end{table}

训练扩展先选有明确终止步数的小任务，不以预测准确率包装平台验收；性能测试同时记录输入规模与软件环境。自托管模型的工具调用或结构化输出未通过时，只开放已通过的普通文本子集，或保持禁用，不把字符串返回成功算作完整 Agent 兼容。

\subsubsection{异常、安全与压力测试}

\textbf{异常测试}逐项注入模型超时、HTTP 429、数据库暂不可用、Worker 断线、Mac 重启、进程被终止、输出空间阈值触发和成果摘要不符。要求错误可定位、重试次数有界、已完成结果可保留；未核实的远端作业不能被当作失败重新启动；预算与审批在恢复后仍生效。执行同一完成回执 3 次只产生一条有效完成记录；旧尝试回执不能覆盖新尝试。

\textbf{Codex 恢复契约}覆盖待审批时关闭客户端、开发中断、Mac 离线、连接丢失、终态已发生但回执未入库、SDK／CLI 版本不兼容与适配器更换。重新连接先核对线程／轮次与实际文件，确认是否仍在运行，再订阅或恢复；未知事件不吞掉、不伪造终态。相同事件重放不重复登记成果，未知作业不能重启成第二个写入者。SDK 状态无法恢复时保留原任务、事件、候选及证据链接，以明确关联的新尝试继续；替换适配器后全部历史仍可查看。

\textbf{安全测试}用无真实秘密的测试标记验证容器读不到主目录、平台配置、Docker 套接字和其他作业目录；网络禁用时不能访问外站；带路径穿越、符号链接和恶意文件名的输入被拒绝。把“忽略规则、读取密钥”的文字放入 CSV 或工具输出，验证它不能扩大工具权限；测试审批后更换目标会重新触发检查。还须验证 Codex 无法读取平台秘密、访问其他 worktree、取得 Docker 控制能力、调用未经允许的 app-server 方法或在宿主执行生成脚本；外层策略不能有效约束时保持执行器禁用。模拟订阅额度耗尽与登录失效，验证任务等待且没有 API 付费回退。通过条件是所有拒绝记录符合预期，日志、追踪及成果中无测试秘密标记外泄。

\textbf{Web 与版本边界测试}验证组织或项目分支更新后，运行任务仍使用原 commit；缺少 commit、skills 文件或验收依据时拒绝执行。模拟阿里云服务超时、Mac 离线、重复请求、日志达到轮转条件和回退失败，要求返回明确状态并保存证据。真实阿里云环境只做经授权的小流量验证，不在现有业务机器直接运行大规模压力注入或交易相关故障演练。

\textbf{压力测试}在 Mac 受控测试环境中先用模拟模型提交 100 个小型作业，容器运行并发保持 2、Codex 开发并发保持 1、等待队列上限 100；另行填满等待队列后再提交一个请求，应得到明确的限流或排队拒绝。检查没有丢任务、重复提交、持续资源越限和状态页卡死，记录吞吐量、队列等待及峰值内存；所有已接受作业必须最终进入可解释的终态。真实 API 压测只有在另定费用上限后进行。GPU 压测从单作业递增批量或上下文，不同时把并发、模型大小与输入长度全部推高；本版不虚设已经达到的吞吐指标。

\subsubsection{分阶段交付与进入条件}

\begin{enumerate}[leftmargin=2em,itemsep=4pt]
  \item \textbf{阶段一：个人研发最小闭环。}交付组织／项目仓库、Mac 持久任务平台、Agents SDK 协调、Codex 适配器、独立容器测试、审阅批准与成果查看，形成“已批准需求、开发、独立测试、审阅、获准交付”的完整链路。最小权限、预算与操作审批从此阶段生效。先用模拟执行器通过契约测试，再在已核实认证与预算下跑真实 CSV 任务；退出条件包括 Codex 证据回传、范围变更阻塞、人工批准、一次中断恢复和备份恢复。dot 的仓库交接及合成 Web 流程可独立验证，不作为核心闭环必须在线的依赖。
  \item \textbf{阶段二：4090 与多项目协作。}在阶段一稳定后，增加多个项目的独立权限／配额、依赖交接和恢复测试，完成一般 Web 的受控模拟闭环；确认 4090 内存、系统、驱动、磁盘及私网通信后接入单 GPU Worker。退出条件是多项目隔离、任务幂等、训练检查点、失联回传及成果核验通过。真实阿里云验证另设发布检查点，须先盘点既有服务、核准操作、资源余量与回退；模拟通过不代表真实部署，不为验收立即扩容或改动网格业务。
  \item \textbf{阶段三：按实际需要扩展常在线与其他后端。}仅在工作时段不足、并发量或可用性要求确实变化时，再决定常在线控制节点与迁移方案；先验证数据库、文件、身份和审批恢复，不把小型阿里云实例直接当完整平台。T640 先修复网络，再验证 V100 兼容引擎与模型；建议连续 3 个工作日、每 5 分钟进行认证健康检查和短请求，成功率至少 99\% 且无未解释持续中断，并另行通过质量／能力验收。其他模型供应商或开发执行器通过相同适配契约后逐项接入，替换执行器不得丢失任务历史与证据；未通过则保持原后端或明确等待。
\end{enumerate}

阶段标准可在开始实施前依据实际使用时段调整，但须先记录再测试；变更阈值不能用于事后掩盖失败。平台验收只证明研发设施满足约定范围，具体业务项目仍分别提交需求、方案、试验和验收文档；固定参数网格与 RSI 简单策略模拟均在另立项目中规定验收，不把交易盈利作为平台标准；本章不授权实盘交易、资金操作、云服务器扩容或第二块 GPU 采购。

\subsubsection{配置参考与运行手册}

五组 YAML 文件作为\textbf{拟建平台的自定义配置接口}，由应用进行类型、范围、路径和交叉依赖校验，不是 Agents SDK、Codex SDK 或 Docker 自带的字段：\texttt{platform.yaml} 管目录、入口、数据库、运行时段和预算；\texttt{agents.yaml} 管角色、指令、工具与模型引用；\texttt{models.yaml} 管后端与能力；\texttt{workers.yaml} 管节点、环境、配额以及开发执行器适配、运行时、认证引用与会话恢复；\texttt{policies.yaml} 管权限、审批、网络、日志和保留规则。配置提交不得包含秘密，实际值在部署时从受保护位置注入。

\texttt{deploy/compose.yaml} 只在首期编排本地 PostgreSQL 等基础服务，使用明确版本及经验证镜像摘要、回环端口、健康检查和持久卷；原生控制进程由 Mac 上的明确启动入口管理。实现完成并生成锁定文件后，运行手册可采用下列\textbf{拟定命令约定}；其中 \texttt{rd\_platform} 模块、检查接口及命令行入口均需编写，不表示当前已有这些文件或已经运行成功。

\Needspace{17\baselineskip}
\par\noindent\begin{minipage}{\linewidth}
\begin{lstlisting}
cd ~/agent-platform
uv sync --locked --group dev
docker compose --env-file ~/.config/agent-platform/runtime.env \
  -f deploy/compose.yaml up -d postgres
uv run --locked python -m rd_platform.cli config-check --config-dir config
uv run --locked alembic upgrade head
uv run --locked uvicorn rd_platform.api:app --host 127.0.0.1 --port 8000
# Run coordinator and worker as separate managed processes:
uv run --locked python -m rd_platform.coordinator
uv run --locked python -m rd_platform.worker --node mac-local
# Verify process health and dependency readiness separately:
curl --fail http://127.0.0.1:8000/healthz
curl --fail http://127.0.0.1:8000/readyz
\end{lstlisting}
\end{minipage}\par


应用启动入口负责从受保护位置加载真实配置；\texttt{config-check} 只输出缺失字段与校验结论，不回显秘密。\texttt{/healthz} 表示进程响应，\texttt{/readyz} 必须实际核对数据库、配置与可写工作目录，不能假装验证 OpenAI 账户或远程 GPU。手册须另列 Codex 执行器启动前的运行时核验、认证类型检查、权限自测、事件通道与中断／恢复契约；只输出认证状态，不打印凭据。接口未实现或契约测试未通过时配置保持禁用。手册还须包含启动与停止顺序、私网接入前置条件、备份恢复、密钥失效、预算阻塞、作业待核实、磁盘清理、版本回退及人工接管步骤，并给出每一步的检查命令、预期结果和失败去向。运行手册另登记组织与项目仓库地址、上下文 commit、dot 成果交接方式、阿里云既有服务清单、新增 Web 发布／回退步骤、监控判断和本地归档位置；只有未来确认自动接入协议后才把新执行端加入 \texttt{workers.yaml}。最终交付应附安装记录、配置摘要、测试证据和已知限制；未实现、未授权和未通过的项目分别列明。


\newpage


\subsection{费用预算与成本控制}
\label{sec:impl-budget}

首期沿用现有计算设备，以 Mac 运行控制服务，GitHub 管理程序与文档，dot 承担获准的轻量开发工作，Codex 承接通过适配器提交的开发任务；其他 Agent 的模型调用按所配置的服务和认证路径分别计费。\textbf{当前不需要把新增计算设备作为启动前提。}本节先给出可复核的 API 用量预算，再用待填账单项形成全系统月度费用；负责人已确认现有 ChatGPT Pro 订阅为每月 100 美元；其余订阅、IDC、云服务器及电费仍须核实，不能把 API 金额直接称为系统总预算。所有估算用于规划，不构成采购、续费、扩容或提高运行限额的授权。

\subsubsection{预算口径与费用构成}

预算同时保留两种视图：\textbf{现有支出加新增支出的全量视图}，用于判断长期承担能力；\textbf{本项目带来的增量视图}，用于判断试点和扩展是否值得。设备已购置不等于运行免费；已付年费按月摊列用于比较，实际付款月份仍按合同登记。折旧与维护工时可在经济成本视图单列，不能与现金支出重复相加。

\begingroup\small\renewcommand{\arraystretch}{1.15}
\begin{longtable}{P{2.3cm}P{6.8cm}P{5.6cm}}
\toprule
\textbf{费用类别} & \textbf{当前基线} & \textbf{新增部分与登记要求}\\
\midrule\endfirsthead
\toprule\textbf{费用类别} & \textbf{当前基线（续）} & \textbf{新增部分与登记要求}\\\midrule\endhead
一次性计算设备 & 已有 Mac M3 Max 64 GB／2 TB；i9／RTX 4090 台式机；T640 128 GB／2 TB、V100S 32 GB。4090 主机内存 64／128 GB 待核。 & 本轮不新增计算设备；历史购置额不再算作启动新增现金。第二块 GPU、显存或整机扩容另行论证。\\
\rowcolor{lightgray}
一次性配套投入 & 备份介质与网络状态尚须盘点；T640 网络修复是接入前置条件。 & 备份盘、必要网络维修按实际缺口询价，金额待定；未确认前不记为零，也不默认购买。\\
固定云端费用 & 阿里云 2 核／2 GB／40 GB 已承担 Web 与网格业务；原套餐、续费日和金额待核。 & 先复用已批准的余量；因新业务产生的套餐升级、额外盘或带宽另列，不能用新增方案覆盖原业务成本。\\
\rowcolor{lightgray}
固定软件与服务 & 已有 ChatGPT Pro 100 美元／月（负责人确认）；GitHub 为版本核心。其他席位与附加费用待核。 & Pro 纳入已有费用，不重复记为本项目新增订阅。额度、超额、税费及其他服务以实际账户账单核对。\\
固定托管费用 & T640 所在 IDC 的机位、网络、电力收费及合同未知。 & 先查是否已含电费或带宽，避免重复计算；恢复服务带来的合同变更单列。\\
\rowcolor{lightgray}
按量运行费用 & 实际走独立付费路径的模型 API；本地计算电力；必要存储、备份与网络传输。 & 按供应商账单与本地测量登记，区分既有基线和本项目增量；工具调用等附加费不并入纯 token 单价。\\
\bottomrule
\caption{现有资产、固定支出与新增按量费用分开登记。}\\
\end{longtable}\endgroup

\textbf{复用订阅的边界。}Mac 上的 ChatGPT 桌面应用及 VS Code 插件属于人工入口；本地 Codex 的程序化开发任务应核验实际登录方式和额度。使用受支持的 ChatGPT 账号认证时，可消耗相应订阅包含量；使用 API Key 时按 API 计费。dot 内获准的文档和轻量 Web 开发按该产品实际额度结算，不再假设同一工作还发生一次独立 API 调用。开发结果部署后的云主机、存储与网络费用仍然存在。多个 Agent 或多个任务不会新增订阅额度。\cite{v121-codex-auth,v121-codex-exec,v121-plan-usage}

Pro 的 100 美元不能直接作为 API 充值抵扣，也不能保证覆盖外层 Agents SDK 运行所需的全部模型调用。订阅、API 和本地计算是三类成本来源，SDK 名称不是计费方式。通过 Sign in with ChatGPT 让自建应用直接使用订阅属于另一条需核验资格、授权范围与接口限制的路径；本方案不以该私有客户端必然获准为前提。\cite{v121-plan-quickstart}

\textbf{软件许可与服务费用分开。}GitHub Free 可用于普通私有仓库，个人使用 Docker Desktop 可符合免费条件，OpenAI Agents SDK 采用 MIT 许可；在这些适用条件下，首期可不增加对应的基础软件许可费。GitHub Actions、LFS 的超额用量、商业支持及其他付费服务仍须单列；现有账户可能已有付费套餐，不能据免费选项反推实际账单为零。Docker Desktop 的组织使用条件须另行核对。\cite{budget-github-plans,budget-github-actions,budget-github-lfs,budget-docker-license,budget-sdk-license}

\Needspace{9\baselineskip}
\textbf{月度汇总模板。}令 $U_0$ 为已有美元月度运行基线，$F_0$ 为已有人民币月度运行基线，二者均包含原有固定及按量费用（含已归集的原有电费、存储与网络费）；$\Delta U$ 为新增美元月费（包括 API、工具及新增订阅），$\Delta F$ 为新增人民币固定月费；$E$ 为本项目实测增量电费，$S$ 为新增且未在前项重复计入的人民币存储、备份与网络费，$r$ 为付款时实际人民币／美元结算汇率，则：
\[
C_{\mathrm{month}}=(U_0+\Delta U)\,r+F_0+\Delta F+E+S.
\]

按目前已知信息，$U_0$ 中已确认的 Pro 项为 100 美元／月，其余既有美元费用仍待补齐；不把 $U_0$ 整体直接定为 100。若订阅内用量足够且不购买额外额度，该部分新增订阅现金可记为零，但其使用量仍应记录。

本项目新增月度费用另记为 $\Delta C_{\mathrm{month}}=\Delta U\,r+\Delta F+E+S$，便于与既有费用基线分开比较。

此式是月度费用归集口径；\textbf{首次启动现金另加已批准的一次性投入}，年付合同的付款月现金流另列。$U_0$、$F_0$ 等未知项须标注“待填”，并附币种、周期、发票／账单来源与是否包含电力或流量；缺项未补齐时只报告已知部分，不能输出一个貌似完整的总额。其他币种另按实际结算换算，不采用未核验的实时汇率。

本地增量电费按各节点工作负载相对既有运行基线增加的墙插平均功率计算：
\[
E=\sum_j \frac{\Delta P_j\,(\mathrm{W})\times h_j\,(\mathrm{h})}{1000}
       \times q_j\,(\mathrm{CNY/kWh}).
\]

$\Delta P_j$、运行时长 $h_j$ 与电价 $q_j$ 均待测／待填；若原本关机而为本项目开机，应计入该时段的整机新增用电。GPU 的 TDP、铭牌功率和电源额定功率不能代替墙插实测。IDC 已打包收费的电力不重复相加；自托管推理还占用硬件、运维和备份资源，不能视为零成本替代。

\subsubsection{模型费率与计费方法}

下表采用 2026-10-12 核验的 OpenAI 官方\textbf{标准文本 API}价格，仅选两档模型作为预算样例，不声称其为最新型号，也不代表最终模型路由；不表示已在当前账户开通、已完成任务质量验证或保证未来价格。费用单位均为美元／百万 token；本表只适用于独立 API 付费路径，不适用于订阅包含量的折算；dot 与 Codex 的订阅用量及可能的额外收费须按实际账户核实。\cite{budget-openai-pricing,budget-mini,budget-gpt54,budget-chatgpt-api}

\begin{table}[H]
\centering\small\renewcommand{\arraystretch}{1.2}
\begin{tabularx}{\textwidth}{P{3.5cm} L L L}
\toprule
\textbf{预算用模型档位} & \textbf{普通输入} & \textbf{缓存输入} & \textbf{输出}\\
\midrule
经济档 GPT-5.4-mini & \$0.75 & \$0.075 & \$4.50\\
\rowcolor{lightgray}
较强档 GPT-5.4 & \$2.50 & \$0.25 & \$15.00\\
\bottomrule
\end{tabularx}
\caption{标准 API 样例费率；较强档按输入不超过 272K token 的定价条件估算。}
\end{table}

GPT-5.4 超过 272K 输入 token 时适用更高的整会话费率，不能继续套用上述短上下文单价；官方说明相应输入为两倍、输出为 1.5 倍。区域处理端点等特殊选项也可能加价。本节情景不采用长上下文加价、区域选项、优先级服务、批处理折扣或缓存优惠；若配置改变，先重新估算。\cite{budget-gpt54,budget-mini}

令模型 $m$ 的普通输入、缓存输入、计费输出量分别为 $I_m$、$K_m$、$O_m$，单位均为\textbf{百万 token}，对应费率为 $p_{i,m}$、$p_{k,m}$、$p_{o,m}$，则：
\[
C_{\mathrm{API}}=\sum_m
 (I_m p_{i,m}+K_m p_{k,m}+O_m p_{o,m})+C_{\mathrm{tools}}.
\]

缓存输入是总输入中的一部分，必须先拆分，不能与普通输入重复相加。输出量以供应商计费 usage 为准，包含计费的推理 token，不能只数用户最终看到的文字。\cite{budget-reasoning} $C_{\mathrm{tools}}$ 单列供应商工具、搜索、容器、文件存储等实际附加收费；下节纯文本测算取“不启用这些收费项目”的边界，并非断言所有工具免费。\cite{budget-openai-pricing}

\textbf{多 Agent 按所有实际调用累计。}一个任务的需求澄清、规划、实现、独立评审、测试修复、重试和恢复调用均会消耗额度；重新附带的历史上下文、工具返回和共享文档也计入相应输入。月度量是所有任务和所有角色的计费总量，不是单次提示词长度乘任务数的简单承诺。计量记录绑定任务、尝试、角色、执行器、认证／计费路径、模型、usage 与费率版本；供应商账单是最终对账依据，超时但已产生的计费同样要追踪。

\subsubsection{三档月度估算}

三档均为\textbf{实际进入独立 API 付费路径的假设用量}，尚无试运行数据支持其对应的任务数或交付能力。每档的输入 token 与输出 token 都分别按经济档 80\%、较强档 20\%分配；这是 token 比例，不是请求数量、Agent 数量或工作时间比例。采用普通输入、不计缓存优惠，得到混合输入费率 $0.8\times0.75+0.2\times2.50=1.10$，混合输出费率 $0.8\times4.50+0.2\times15.00=6.60$，单位均为美元／百万 token。

\begin{table}[H]
\centering\small\renewcommand{\arraystretch}{1.22}
\begin{tabularx}{\textwidth}{P{4.1cm} L L L}
\toprule
\textbf{月度情景／计算项} & \textbf{低量试点} & \textbf{日常使用} & \textbf{高量压力}\\
\midrule
计费输入／输出（百万 token） & 10／2 & 50／10 & 200／40\\
\rowcolor{lightgray}
经济档输入／输出（百万） & 8／1.6 & 40／8 & 160／32\\
较强档输入／输出（百万） & 2／0.4 & 10／2 & 40／8\\
\rowcolor{lightgray}
输入费用（美元） & 11.00 & 55.00 & 220.00\\
输出费用（美元） & 13.20 & 66.00 & 264.00\\
\rowcolor{lightgray}
API 小计（美元） & 24.20 & 121.00 & 484.00\\
20\% API 预留（美元） & 4.84 & 24.20 & 96.80\\
\rowcolor{lightgray}
含预留金额（美元） & 29.04 & 145.20 & 580.80\\
便于规划的取整额度（美元） & 30 & 150 & 600\\
\bottomrule
\end{tabularx}
\caption{三档纯文本 API 月度情景。取整额度仍须审批，且不含全系统其他费用。}
\end{table}

已经由订阅承担的开发任务不再重复计入本表；不能从这里按 API 单价算出其“订阅抵扣额”。只有订阅可用性、任务分配及试运行用量得到验证后，才能重新估计独立 API 部分。因此，30／150／600 美元不是使用本平台必然增加的月费。

例如日常使用档为 $50\times1.10+10\times6.60=121$ 美元，预留后为 $121\times1.20=145.20$ 美元。20\%用于用量偏差和额外重试的规划缓冲，不意味着价格、任务完成率或最终账单得到保证；明显增加较强模型比例、长上下文或收费工具时，应重新做情景估算，不能假定缓冲足够覆盖。

\textbf{与现有运行阈值的关系。}第 5.4 节的单任务 2 美元、每日 10 美元是示例保护阈值，不是每项任务的价格或月费预测。按 30 日比较，高量档 API 小计平均为 16.13 美元／日，含预留平均为 19.36 美元／日，已超过每日 10 美元；\textbf{当前示例日限额下不能实现该高量情景}。如确需该用量，须另行批准并调整配置，本节不改变旧阈值。低量和日常档平均值虽较低，集中在少数天运行仍可能触发日限额；单个长任务也可能先触发单任务限制。

建议先复用现有订阅完成可支持的文档与开发任务，再把\textbf{约 30 美元／月的独立 API 试点额度}作为确有调用需求时的待批准备选，观察实际质量、重试比例与任务分布后再决定是否进入约 150 美元档。约 600 美元档是压力情景，不是首期默认额度，更不是全系统月费。T640 未完成网络和质量验收前不预扣云端费用；验收后用同一任务集比较 API 节省额、增量电费和运维成本，再决定可迁移比例。

\subsubsection{预算控制与复核}

\begin{enumerate}[leftmargin=2em,itemsep=4pt]
  \item \textbf{启动前补齐账单基线。}登记已知 Pro 100 美元／月，核对阿里云套餐／续费、GitHub 席位与附加用量、订阅可用额度、IDC 合同、电价和备份缺口。分别填写原有月费、新增月费与一次性投入；确认可接受月度上限及付款币种后再启用付费主链路。
  \item \textbf{调用前预留，结束后结算。}对独立 API 付费调用，平台按预计输入、最大计费输出及适用费率预留任务、每日和月度额度，并把并发在途请求一起占用预算；实际 usage 返回后释放多余预留。无 usage、超时或连接中断的调用先标记待对账，不立即当作零成本释放额度。预算不足时暂停派发，不靠重试绕过限制。订阅路径单独读取可获得的额度与限流状态；未知剩余额度时保守限制并发，不将美元 token 公式套作订阅余额。
  \item \textbf{限制上下文与无效循环。}按角色只传必要材料；长记录先做有来源的摘要，保留引用与可恢复原文。为修复轮次、工具调用次数、重试和单轮输出设置上限；较强模型用于确有质量需要的阶段，不因失败而无限升级模型。缓存或批处理只有在任务语义和服务条件适合时才纳入节省额。
  \item \textbf{预警与停止分开。}沿用第 5.6.4 节的 80\%预算告警检查点，触发后核对剩余任务与在途预留，并提醒负责人；达到已批准可用额度上限时暂停对应路径的新任务并请求决定。订阅耗尽、认证失效或模型不可用，不得静默切换到 API Key、购买额度或提高套餐；转入待决状态并记录原因。供应商预算通知不替代应用端的并发预留与派发检查，已在途调用的费用也须计入。
  \item \textbf{每周核对，首月重估。}每周按角色、模型和任务汇总输入、输出、缓存、重试、成功交付数及待对账金额；与供应商用量和账单核对差异。首月结束后再估算下一月的 token 量及固定费用，保留失败任务消耗，避免只用成功样例推算成本。
\end{enumerate}

验收时应能回答：本月已花多少、在途预留多少、剩余额度多少；哪类任务增加了成本；供应商账单与内部台账为什么不同。价格、模型、工具、地区、订阅或运行时段变更时更新预算版本；自动化可以暂停、告警和准备申请，不能自行增加订阅、提升额度、扩容云资源或采购 GPU。


\newpage
\section*{附录}
% ════════════════════════════════════════════════
\phantomsection\addcontentsline{toc}{section}{附录}

\subsection*{A. 术语对照表}
\phantomsection\addcontentsline{toc}{subsection}{A. 术语对照表}

\begin{table}[H]
\centering\small
\renewcommand{\arraystretch}{1.25}
\begin{tabularx}{\textwidth}{P{3.5cm} L}
\toprule
\textbf{组织规范术语} & \textbf{系统实现对应} \\
\midrule
人类负责人 & 在系统外提出需求与约束，接收成果；必要时作出决定 \\
\rowcolor{lightgray}
组织层协调 & 协调角色，由编排机制支撑任务流转 \\
研发协调负责人 & 协调角色 + 项目与版本视图 \\
\rowcolor{lightgray}
各组负责人 & 对应专业 Agent 及其工具授权范围 \\
需求包 & 项目管理工具中的一个迭代所含需求集合 \\
\rowcolor{lightgray}
规格调整 & 变更请求 + 审批流 + 关联记录 \\
交接 & 结构化交接消息 + 接收方确认 \\
\rowcolor{lightgray}
接收确认 & 接收方校验完整性后返回确认或退回缺项清单 \\
版本验收 & 测试结论、接管条件与需求满足反馈分别呈现；重要取舍由人类参与 \\
\rowcolor{lightgray}
独立测试 & 独立依据与结论；不能自行更改被测实现并批准修复 \\
科研任务的否定结论 & 项目管理中科研类任务的独立验收标准 \\
\bottomrule
\end{tabularx}
\end{table}

\subsection*{B. 能力清单}
\phantomsection\addcontentsline{toc}{subsection}{B. 能力清单}

下表概括框架需要支撑的能力，不表示每项都要独立部署一套产品。实施时可复用已有工具，按实际需要组合。

\begin{table}[H]
\centering\small
\renewcommand{\arraystretch}{1.25}
\begin{tabularx}{\textwidth}{P{3.5cm} L}
\toprule
\textbf{能力} & \textbf{作用} \\
\midrule
编排与工具执行 & 流程推进、模型路由、工具校验、状态管理与 HiL \\
\rowcolor{lightgray}
角色与项目协作 & 职责、进度、依赖、交接、成果和待决事项可见 \\
意图与结果交互 & 需求输入、结果与证据交付、必要澄清和追加授权 \\
\rowcolor{lightgray}
代码与交付管理 & 版本、评审、构建、测试证据和变更追溯 \\
实验与计算 & 支撑数据处理、实验、训练和结果比较 \\
\rowcolor{lightgray}
数据管理 & 必要资料、认知与运营记录的取用、留存和恢复 \\
安全与可观测性 & 身份、权限、凭证、审计、状态和用量可观测 \\
\rowcolor{lightgray}
运行保障 & 部署、监控、维护、备份与故障恢复 \\
\bottomrule
\end{tabularx}
\end{table}

\subsection*{C. 模型 API 接入原则}
\phantomsection\addcontentsline{toc}{subsection}{C. 模型 API 接入原则}

模型与推理投入按任务配置，运行控制程序执行预算与并发限制，模型适配器负责接入、计量和错误处理，模型供应商提供实际推理服务。模型负责推理并返回请求，执行环境负责校验、执行和证据回传；更换供应商后仍需验证行为。凭证、传输、重试和费用应受控，调用失败时保留任务状态与恢复路径。第二章已说明单个 Agent 的模型与工具调用及上下文组织，第三章说明实例间交换与共享服务机制；实施时对选定 SDK 和模型版本进行契约测试，验证权限拒绝、超时、工具错误和用量记录。

\subsection*{D. 部署架构形态参考}
\phantomsection\addcontentsline{toc}{subsection}{D. 部署架构形态参考}
\label{app:deployment}
\begingroup\renewcommand{\thesection}{D}\renewcommand{\theHsection}{appendixD}\setcounter{subsection}{0}
% ════════════════════════════════════════════════

本附录保留此前的部署比较，作为后续实施设计的参考。以下本地优先基线属于起步候选建议，本项目当前采用的节点分工与配置基线见第五章；本附录的比较与早期起步建议作为背景参考。其中涉及具体产品、任务服务或产物登记方式的内容均为候选实现，不表示已经完成选型，也不构成第四章逻辑设计的前置限制。本附录比较执行系统的三种部署形态：本地单机、本地多机与全云。云端模型 API 在三种形态中都可以使用，不能据此把本地方案称为全云；执行系统迁到云端则属于另一个部署决定。

\subsection{形态划分的依据}

\subsubsection{两个需要分别决定的维度}

部署方案至少需要分别考虑以下两个维度，它们的成本、安全与可用性仍会相互影响：

\begin{table}[H]
\centering
\renewcommand{\arraystretch}{1.4}
\begin{tabularx}{\textwidth}{P{3cm} P{4.5cm} L}
\toprule
\textbf{维度} & \textbf{取值} & \textbf{影响} \\
\midrule
执行系统的部署位置 & 本地单机 / 本地多机 / 全云 & 可用性、成本结构、数据主权、运维负担 \\
\rowcolor{lightgray}
模型推理的提供方式 & 云端模型 API / 自行部署模型 & 能力上限、离线能力、边际成本 \\
\bottomrule
\end{tabularx}
\caption{部署方案的两个维度}
\end{table}

本附录主要按第一个维度展开；模型推理当前选择云端 API。自行部署模型也需另行决定其承载位置，执行系统与推理可以形成多种组合，但都要满足数据、性能和运行条件。云端执行环境还需区分自管云服务器与托管服务：前者仍需维护应用和环境，后者仅在约定范围内转移部分运维责任。

\subsubsection*{当前模型推理方式的取舍}

当前基线调用云端模型，无需在本地常驻大模型推理服务，本地算力可优先用于训练、数据处理、开发与测试；相应地，需要管理 API 的可用性、超时与版本变化。下表比较当前选择与本地模型推理的主要差异。

\begin{table}[H]
\centering
\renewcommand{\arraystretch}{1.35}
\begin{tabularx}{\textwidth}{P{2.4cm} P{3.8cm} L}
\toprule
\textbf{比较维度} & \textbf{云端模型 API（当前选择）} & \textbf{本地部署大模型} \\
\midrule
GPU 用途 & 优先用于训练与本地任务 & 推理与其他负载需分配资源 \\
\rowcolor{lightgray}
模型能力上限 & 取决于可用服务与调用条件 & 受硬件及模型规模约束 \\
运维复杂度 & 无需自建模型服务，仍需维护接入 & 需维护模型、硬件及服务 \\
\rowcolor{lightgray}
离线能力 & 推理依赖网络，其他任务视依赖而定 & 推理可离线，外部依赖另计 \\
成本模式 & 按服务条款与调用量付费 & 硬件、电力、维护及更新成本 \\
\bottomrule
\end{tabularx}
\caption{云端模型 API 与本地模型推理的取舍}
\end{table}

当前阶段以验证架构可行性为目标，因此优先考虑可用模型能力和维护负担，选择云端模型 API。若出现强离线需求，或调用量变化使成本关系改变，可再评估本地模型；这一选择与执行系统采用单机、多机或全云分别判断。

\subsubsection{组件职责与运行约定保持稳定}

三种形态下，应尽量保持以下逻辑约定稳定：

\begin{itemize}[leftmargin=2em,itemsep=2pt]
  \item Agent 的四要素构成，以及模型提出请求、运行环境校验执行的分工；
  \item 人与系统的边界，以及协调、执行、验证和运行控制程序的职责；
  \item 数据按用途管理、必要证据留存与受控传输的原则；
  \item HiL 机制与审计要求。
\end{itemize}

迁移会改变服务与数据的位置，也可能改变网络、身份、可用性与成本。保持逻辑约定稳定有助于减少重做，但仍需验证相关接口和运行行为，不能把迁移视为只移动容器或文件。

\subsubsection{形态之间不是优劣排序}

三种形态没有绝对的优劣，只有与当前条件的匹配程度。判断匹配度需要考虑：数据量与算力需求、是否需要长期在线服务、是否有协作需求、可投入的运维精力、成本承受度、对数据主权的要求。

选择形态时，正确的问题不是哪个最好，而是当前条件下哪个够用且代价可承受。

\subsection{单机架构}

\subsubsection{形态特征}

编排、工具、存储、消息接入与开发环境集中在同一台本地设备上，设备同时承担服务端与客户端角色。当前基线通过云端 API 获取模型推理；这一选择不属于单机形态本身的定义。

\subsubsection{适用条件}

\begin{itemize}[leftmargin=2em,itemsep=2pt]
  \item 处于架构验证阶段，需要快速跑通端到端流程；
  \item 实际数据、计算和并发需求在单设备承载范围内；
  \item 可以接受设备休眠、关机或故障时的服务中断；
  \item 主要由单人使用，访问与任务并发在设备可承载范围内；
  \item 训练任务尚未成为常态，或规模在设备能力范围内。
\end{itemize}

\subsubsection{能力边界}

\begin{table}[H]
\centering
\renewcommand{\arraystretch}{1.35}
\begin{tabularx}{\textwidth}{P{3cm} L}
\toprule
\textbf{约束} & \textbf{说明与应对} \\
\midrule
内存竞争 & 服务容器与开发环境争抢内存。应对：精简服务集，按实际需要精简或移出非核心负载 \\
\rowcolor{lightgray}
无冗余 & 设备故障即全面中断。应对：定期备份，接受中断作为阶段性代价 \\
在线能力随设备状态 & 休眠或关机时服务不可用。应对：明确活跃时段和待机策略 \\
\rowcolor{lightgray}
算力上限 & 训练任务受单设备能力限制。应对：小规模实验为主，大型任务另行评估本地扩展或临时云资源 \\
环境混杂 & 开发、测试、服务运行共用环境，相互干扰。应对：容器化隔离，明确资源配额 \\
\bottomrule
\end{tabularx}
\caption{单机架构的能力边界}
\end{table}

\subsubsection{成本结构}

若已有合适设备，单机形态可以降低起步时的新增投入。总成本仍包括设备折旧、存储备份、电力、维护时间与模型调用费用，不能只看新增账单。

需要留意的是，调用费用通常随模型配置、用量和服务计费规则变化，重试与长上下文也可能明显增加消耗。预算选择与超限处理机制（见第 \ref{sec:orchestration} 节）在此形态下尤为重要。

\subsubsection{向多机演进的触发信号}

出现以下情况时，应评估多机或其他调整，而不是自动扩容：

\begin{itemize}[leftmargin=2em,itemsep=2pt]
  \item 训练任务成为常态，且与开发环境产生明显的资源争抢；
  \item 需要长期在线的服务（如定时数据采集、持续监控、事件响应）；
  \item 内存或存储成为瓶颈，精简服务集已无法解决；
  \item 需要区分开发环境与准生产环境以支持独立测试。
\end{itemize}

\subsection{多机架构}

\subsubsection{形态特征}

多台设备通过局域网协作，按角色分工。典型的角色划分：

\begin{table}[H]
\centering
\renewcommand{\arraystretch}{1.35}
\begin{tabularx}{\textwidth}{P{2.6cm} P{4.2cm} L}
\toprule
\textbf{角色} & \textbf{承担内容} & \textbf{特征} \\
\midrule
算力节点 & 模型训练、数据处理、计算密集任务 & 按任务配置算力，必要时配备 GPU \\
\rowcolor{lightgray}
服务节点 & 常驻服务容器、数据库、监控体系 & 长期在线，稳定性优先 \\
接入终端 & 提交需求、查看结果、必要澄清与授权 & 便携性优先，可随时接入 \\
\bottomrule
\end{tabularx}
\caption{多机架构的角色分工}
\end{table}

三种角色可以由不同设备分别承担，也可以部分合并（如算力节点同时充当服务节点）。合并会牺牲一部分隔离性，换取设备数量的精简。

\subsubsection{适用条件}

\begin{itemize}[leftmargin=2em,itemsep=2pt]
  \item 训练任务常态化，需要专用算力；
  \item 存在需要长期在线的服务；
  \item 已有多台可用设备，无需额外采购；
  \item 具备基本的网络配置与运维能力，或有自动化手段降低运维负担。
\end{itemize}

\subsubsection{核心设计问题：数据存放与同步}

多机架构最关键的设计决策是数据放在哪里、如何在设备间保持一致。有三种基本模式：

\begin{table}[H]
\centering
\renewcommand{\arraystretch}{1.4}
\begin{tabularx}{\textwidth}{P{2.2cm} P{3.8cm} P{3.2cm} L}
\toprule
\textbf{模式} & \textbf{机制} & \textbf{优势} & \textbf{代价} \\
\midrule
集中式 & 数据只存于一个节点，其余设备通过网络访问 & 主数据位置清楚，减少副本分歧 & 依赖网络；该节点成为单点 \\
\rowcolor{lightgray}
副本式 & 各节点自持副本，定期同步 & 各节点可离线工作 & 可能不一致；同步策略复杂 \\
混合式 & 按数据类型分别处理：代码走版本控制，数据走共享存储，缓存各节点自持 & 兼顾一致性与灵活性 & 需要明确规则，纪律要求高 \\
\bottomrule
\end{tabularx}
\caption{数据存放的三种模式}
\end{table}

可按数据用途组合以上模式，当前参考设计采用以下明确规则：

\begin{itemize}[leftmargin=2em,itemsep=2pt]
  \item 代码与配置。通过版本控制系统同步。这本身就是软件开发的标准做法，天然支持多机协作与历史追溯；
  \item 认知数据与运营数据。由中心任务服务管理写入，远程 worker 通过接口提交；不复制多个可写数据库文件来同步状态；
  \item 外部资料与产物。通过 artifact ID 和 manifest 按需取得并验证哈希，节点缓存不可反向覆盖中心版本；
  \item 训练用数据集。可按固定 manifest 拉取本地只读副本，保留可复现实验的版本与划分；大文件存储以后可改为对象存储，逻辑引用保持稳定。
\end{itemize}

多个远程计算节点可先共用一台控制服务，SQLite 留在服务本机；出现多个应用写入者、持续写竞争或更高可用性需求时，再迁移 PostgreSQL。PostgreSQL 的 MVCC 支持事务并发，但不会自动解决业务 revision 冲突或重复外部副作用\cite{postgres-mvcc}。

\subsubsection{网络与远程访问}

多机需要明确服务如何发现、访问身份如何验证，以及数据如何安全传输。远程接入不应以无保护地开放内部服务为前提；建议将控制接口限制在受保护网络内，验证客户端身份并加密传输；所选网络与账户配置需要单独确认和验证。

应识别网络中断和节点不可达的影响，避免未确认的远程操作被误判为完成，并说明哪些任务可以独立继续、哪些任务需要等待恢复。

\subsubsection{成本结构}

多机成本包括设备投入与折旧、电力、备份、网络、维护时间和模型调用。已有设备的利用率较高时可能具有成本优势，但扩容、闲置和故障恢复也会增加成本，应与云端的实际用量方案比较。

一次性投入可能包括网络、存储与供电保护。投入顺序由故障影响和恢复要求决定，不预设某类设备始终优先。

\subsubsection{约束与风险}

\begin{itemize}[leftmargin=2em,itemsep=2pt]
  \item 无服务等级保障。可用性依赖自有设备、网络与维护安排，核心故障需自行处置；
  \item 单点故障影响面扩大。服务节点故障会影响所有依赖它的设备；
  \item 运维负担随节点数上升。每增加一台设备，就多一份环境维护、监控与故障排查工作；
  \item 扩容需要采购。算力不足时通常不能立即获得新增物理容量，需要提前规划。
\end{itemize}

运维负担是多机形态最需要正视的代价。缓解手段是提高自动化程度：配置即代码、部署脚本化、监控告警自动化、故障恢复预案化。运维 Agent 的价值在此形态下最为明显。

\subsection{全云架构}

\subsubsection{形态特征}

执行系统的编排、工具、存储与消息接入部署于云端，用户设备作为接入终端。模型推理由调用的模型服务或另行部署的模型提供，应与执行系统的位置分别决定。仅调用云端模型 API 的本地执行系统不属于全云形态。云端执行环境可以自管，也可以使用托管服务，两者的责任与运维负担不同。

\subsubsection{适用条件}

\begin{itemize}[leftmargin=2em,itemsep=2pt]
  \item 多人协作，需要统一的访问入口与权限管理；
  \item 对在线访问、恢复或可用性有明确要求，且选定服务与部署能满足；
  \item 算力需求波动大，需要弹性伸缩；
  \item 运维精力有限，愿意为合适的托管能力付费；
  \item 数据出域在合规与安全上可接受。
\end{itemize}

\subsubsection{优势与代价}

\begin{table}[H]
\centering
\renewcommand{\arraystretch}{1.4}
\begin{tabularx}{\textwidth}{P{3cm} P{5cm} L}
\toprule
\textbf{优势} & \textbf{说明} & \textbf{对应代价} \\
\midrule
责任可部分转移 & 硬件由平台维护，托管服务进一步承担约定责任 & 应用、安全配置、数据与费用仍需管理 \\
\rowcolor{lightgray}
弹性资源 & 可按需申请，受配额、容量与资源类型限制 & 需管理费用及资源可用性 \\
可配置高可用 & 可利用冗余与恢复服务，效果取决于具体设计 & 额外费用、配置与责任分工 \\
\rowcolor{lightgray}
远程协作 & 容易提供集中入口和共享服务 & 仍需控制访问、数据地域与迁移依赖 \\
\bottomrule
\end{tabularx}
\caption{全云架构的优势与代价}
\end{table}

\subsubsection{成本结构}

全云成本可能包括订阅、计算、存储、网络传输、托管服务与模型调用费用，具体随服务和用量结构变化。比较时应同时计入本地设备折旧、电力、备份和维护时间，而非只比较云账单与本地电费。

全云是否划算，取决于负载持续性、资源利用率、所需服务等级与维护能力。按需使用可能适合低频或波动任务，长期稳定负载则应另行核算，不能仅凭规模断言哪种形态更便宜。

\subsubsection{混合形态}

除纯粹的本地或全云部署外，还可以采用混合形态：

\begin{itemize}[leftmargin=2em,itemsep=2pt]
  \item 本地持有数据与编排，云端提供弹性算力。核心存储保留本地；需要大算力时可评估临时云资源，但训练所需的数据与代码是否可传输须单独判断；
  \item 本地承担常态负载，云端应对峰值。日常服务本地运行，突发的高负载任务分流到云端；
  \item 云端托管协作服务，本地保留计算。项目管理、文档协作等需要多人访问的服务托管于云端，计算密集型工作留在本地。
\end{itemize}

混合形态可以在控制与弹性之间取舍，代价是跨环境管理更复杂。需要明确数据和责任边界、传输范围及故障影响，不能因为核心副本留在本地就忽略云端收到的数据。

当前基线在整体资源使用上结合了本地执行与云端推理；按本附录的执行系统部署口径，它仍属于本地形态。只有执行系统的职责或数据分布在本地和云端时，才属于本节讨论的混合部署。

\subsection{形态选择与演进}

\subsubsection{三形态对比}

\begin{table}[H]
\centering
\renewcommand{\arraystretch}{1.4}
\begin{tabularx}{\textwidth}{P{2.6cm} P{2.8cm} P{2.8cm} L}
\toprule
\textbf{维度} & \textbf{单机} & \textbf{多机} & \textbf{全云} \\
\midrule
成本重点 & 已有设备与调用成本 & 设备利用率及维护 & 用量、托管费与传输费 \\
\rowcolor{lightgray}
可用性 & 取决于设备与恢复 & 取决于冗余与维护 & 取决于部署与服务约定 \\
数据控制 & 本地持有，受控外发 & 本地多节点，需控访问 & 受地域、配置与合同影响 \\
\rowcolor{lightgray}
运维负担 & 与服务数量有关 & 需维护节点及连接 & 自管与托管差异较大 \\
算力上限 & 单设备容量 & 自有设备总容量 & 配额、可用容量与预算 \\
\rowcolor{lightgray}
扩容条件 & 升级或新增设备 & 调整或新增节点 & 申请资源并验证可用性 \\
协作支持 & 可支持，受在线状态约束 & 可设共享服务 & 易远程访问，仍需授权 \\
\rowcolor{lightgray}
起步条件 & 有合适设备即可验证 & 需明确节点与网络 & 需账户、费用与数据安排\\
\bottomrule
\end{tabularx}
\caption{三种形态的对比}
\end{table}

\subsubsection{选择依据}

形态选择应围绕实际约束，而不是把单机、多机和全云视为固定的成长阶梯：

\begin{enumerate}[leftmargin=2em,itemsep=4pt]
  \item 需要验证哪些能力？先确认任务闭环、协作和质量要求，再判断现有设备是否足够；
  \item 需要怎样的在线与恢复能力？单机也可以长期运行，但仍存在设备、网络和维护带来的单点风险；
  \item 数据可以如何使用与传输？若某类内容不能出域，其云端推理与工具路径也必须相应限制；
  \item 能投入多少资金和维护精力？按实际负载比较总体成本，并优先使用可管理的方案。
\end{enumerate}

\subsubsection{演进路径与迁移代价}

形态迁移的代价取决于抽象层的稳定程度。本方案希望降低 Agent 定义、工具和流程对特定位置的依赖，但迁移仍可能涉及：

\begin{itemize}[leftmargin=2em,itemsep=2pt]
  \item 服务容器的重新部署（配置即代码时成本较低）；
  \item 数据的迁移与校验（需要提前规划可导出的存储格式）；
  \item 网络与访问方式的调整；
  \item 监控告警的重新配置。
\end{itemize}

接口和职责保持稳定时，可以减少 Agent 逻辑与业务流程的改写；是否需要调整，应以迁移后的兼容性和行为验证为准。

演进不必单向递进。当条件变化时（如算力需求下降、协作需求出现），形态可以回调或转向混合。

\subsubsection{本项目的建议}

结合当前以个人主导、先厘清框架再实施的阶段，建议：

\begin{enumerate}[leftmargin=2em,itemsep=4pt]
  \item 以 Mac 自建控制闭环与 GitHub 私有版本核心作为起步基线，使用云端模型 API；dot 云电脑补充轻量研发与多 Agent 协作，现有阿里云只承担一般 Web 服务和轻量接口验证；
  \item 当训练负载、在线要求或隔离需要超出单机能力时，再依据实际数据评估多机分工；
  \item 全云不作为当前默认方向，过程数据主要本地保存，各云端服务必要运行状态及日志仍在所属节点；后续扩展重新检查数据传输、权限、归档和运行责任；
  \item 实施时明确验证目标与进入下一阶段的条件，先按第二章验证单任务闭环，并结合第三章检查协作与交接，并在选定方案支持且配置恢复能力时验证中断恢复，再决定资源投入与扩展组件。
\end{enumerate}

% ════════════════════════════════════════════════

\endgroup

\newpage
\section*{参考资料}
\phantomsection\addcontentsline{toc}{section}{参考资料}
以下包括模型研究的原始论文以及软件能力与限制的官方资料。原有软件资料核验日期为 2026-10-11，Hadoop、Agent 软件生态、模型及原第五章实施资料核验日期为 2026-10-12；本次 Codex 接入、Agent 编排与订阅使用资料核验日期为 2026-10-13。本文采用解释性抽象与示意流程，不代表各框架具有完全一致的功能；实施时需固定版本并再次验证。\par
\begingroup
\renewcommand{\section}[2]{}
\begin{thebibliography}{99}
\renewcommand{\baselinestretch}{1.0}
\fontsize{8.5pt}{10pt}\selectfont
\setlength{\itemsep}{0pt}
\setlength{\parskip}{0pt}
\bibitem{hadoop-yarn} Apache Hadoop. \href{https://hadoop.apache.org/docs/current/hadoop-yarn/hadoop-yarn-site/YARN.html}{Apache Hadoop YARN}.\par\url{https://hadoop.apache.org/docs/current/hadoop-yarn/hadoop-yarn-site/YARN.html}
\bibitem{hadoop-hdfs} Apache Hadoop. \href{https://hadoop.apache.org/docs/r3.4.0/hadoop-project-dist/hadoop-hdfs/HdfsDesign.html}{HDFS Architecture}.\par\url{https://hadoop.apache.org/docs/r3.4.0/hadoop-project-dist/hadoop-hdfs/HdfsDesign.html}
\bibitem{langgraph-persist} LangGraph. \href{https://docs.langchain.com/oss/python/langgraph/persistence}{Persistence}.\par\url{https://docs.langchain.com/oss/python/langgraph/persistence}
\bibitem{langgraph-interrupt} LangGraph. \href{https://docs.langchain.com/oss/python/langgraph/interrupts}{Interrupts}.\par\url{https://docs.langchain.com/oss/python/langgraph/interrupts}
\bibitem{codex-sdk} OpenAI. \href{https://learn.chatgpt.com/docs/codex-sdk}{Codex SDK}.\par\url{https://learn.chatgpt.com/docs/codex-sdk} 访问日期：2026-10-13。
\bibitem{codex-app} OpenAI. \href{https://learn.chatgpt.com/docs/app-server}{Codex app-server}.\par\url{https://learn.chatgpt.com/docs/app-server} 访问日期：2026-10-13。
\bibitem{openai-function} OpenAI. \href{https://developers.openai.com/api/docs/guides/function-calling}{Function calling}.\par\url{https://developers.openai.com/api/docs/guides/function-calling}
\bibitem{mlflow-tracking} MLflow. \href{https://mlflow.org/docs/latest/ml/tracking/}{ML Experiment Tracking}.\par\url{https://mlflow.org/docs/latest/ml/tracking/}
\bibitem{pytorch-repro} PyTorch. \href{https://docs.pytorch.org/docs/2.14/notes/randomness.html}{Reproducibility}.\par\url{https://docs.pytorch.org/docs/2.14/notes/randomness.html}
\bibitem{fastapi-background} FastAPI. \href{https://fastapi.tiangolo.com/tutorial/background-tasks/}{Background Tasks}.\par\url{https://fastapi.tiangolo.com/tutorial/background-tasks/}
\bibitem{mcp-transports} Model Context Protocol. \href{https://modelcontextprotocol.io/specification/2026-07-28/basic/transports}{Transports，2026-07-28 规范}.\par\url{https://modelcontextprotocol.io/specification/2026-07-28/basic/transports}
\bibitem{a2a-spec} A2A. \href{https://a2a-protocol.org/v1.0.0/specification/}{Protocol Specification v1.0.0}.\par\url{https://a2a-protocol.org/v1.0.0/specification/}
\bibitem{sqlite-wal} SQLite. \href{https://www.sqlite.org/wal.html}{Write-Ahead Logging}.\par\url{https://www.sqlite.org/wal.html}
\bibitem{sqlite-use} SQLite. \href{https://www.sqlite.org/whentouse.html}{Appropriate Uses For SQLite}.\par\url{https://www.sqlite.org/whentouse.html}
\bibitem{git-worktree} Git. \href{https://git-scm.com/docs/git-worktree}{git-worktree}.\par\url{https://git-scm.com/docs/git-worktree}
\bibitem{pytest-usage} pytest. \href{https://docs.pytest.org/en/stable/how-to/output.html}{Managing output and JUnit XML}.\par\url{https://docs.pytest.org/en/stable/how-to/output.html}
\bibitem{sqlite-backup} SQLite. \href{https://www.sqlite.org/backup.html}{Online Backup API}.\par\url{https://www.sqlite.org/backup.html}
\bibitem{postgres-mvcc} PostgreSQL. \href{https://www.postgresql.org/docs/current/mvcc-intro.html}{MVCC Introduction}.\par\url{https://www.postgresql.org/docs/current/mvcc-intro.html}
\newpage
\bibitem{langgraph-overview} LangChain. \href{https://docs.langchain.com/oss/python/langgraph/overview}{LangGraph overview}.\par\url{https://docs.langchain.com/oss/python/langgraph/overview}
\bibitem{deepagents} LangChain. \href{https://docs.langchain.com/oss/python/deepagents/overview}{Deep Agents overview}.\par\url{https://docs.langchain.com/oss/python/deepagents/overview}
\bibitem{crewai-flows} CrewAI. \href{https://docs.crewai.com/en/concepts/flows}{Flows}.\par\url{https://docs.crewai.com/en/concepts/flows}
\bibitem{agents-sdk} OpenAI. \href{https://openai.github.io/openai-agents-python/}{Agents SDK}.\par\url{https://openai.github.io/openai-agents-python/}
\bibitem{agents-models} OpenAI. \href{https://openai.github.io/openai-agents-python/models/}{Agents SDK Models}.\par\url{https://openai.github.io/openai-agents-python/models/}
\bibitem{microsoft-agent} Microsoft. \href{https://learn.microsoft.com/en-us/agent-framework/overview/}{Agent Framework Overview}.\par\url{https://learn.microsoft.com/en-us/agent-framework/overview/}
\bibitem{microsoft-providers} Microsoft. \href{https://learn.microsoft.com/en-us/agent-framework/integrations/by-component/model-providers/}{Agent Framework Model providers}.\par\url{https://learn.microsoft.com/en-us/agent-framework/integrations/by-component/model-providers/}
\bibitem{autogen-status} Microsoft. \href{https://github.com/microsoft/autogen}{AutoGen maintenance notice}.\par\url{https://github.com/microsoft/autogen}
\bibitem{google-adk} Google. \href{https://github.com/google/adk-python}{Agent Development Kit 2.0}.\par\url{https://github.com/google/adk-python}
\bibitem{adk-models} Google. \href{https://adk.dev/tutorials/agent-team/}{ADK Agent team tutorial}.\par\url{https://adk.dev/tutorials/agent-team/}
\bibitem{agentscope} AgentScope. \href{https://github.com/agentscope-ai/agentscope}{AgentScope 2.0}.\par\url{https://github.com/agentscope-ai/agentscope}
\bibitem{agentscope-runtime} AgentScope. \href{https://github.com/agentscope-ai/agentscope-runtime}{Runtime archive notice}.\par\url{https://github.com/agentscope-ai/agentscope-runtime}
\bibitem{qwen-agent} Qwen. \href{https://github.com/QwenLM/Qwen-Agent}{Qwen-Agent}.\par\url{https://github.com/QwenLM/Qwen-Agent}
\bibitem{dify-agent} Dify. \href{https://docs.dify.ai/en/cloud/use-dify/nodes/agent}{Agent node}.\par\url{https://docs.dify.ai/en/cloud/use-dify/nodes/agent}
\bibitem{n8n-ai} n8n. \href{https://docs.n8n.io/advanced-ai/}{Integrate AI}.\par\url{https://docs.n8n.io/advanced-ai/}
\bibitem{bailian-apps} 阿里云百炼. \href{https://help.aliyun.com/zh/model-studio/application-introduction}{应用类型介绍}.\par\url{https://help.aliyun.com/zh/model-studio/application-introduction}
\bibitem{agentrun} 阿里云. \href{https://help.aliyun.com/zh/agentrun/what-is-agentrun}{什么是 AgentRun}.\par\url{https://help.aliyun.com/zh/agentrun/what-is-agentrun}

% New references used only by Chapter 2. Insert inside the existing bibliography.
\bibitem{gpt3-fewshot} Brown, T. B., et al. \href{https://arxiv.org/abs/2005.14165}{Language Models are Few-Shot Learners}. 2020.\par\url{https://arxiv.org/abs/2005.14165}
\bibitem{instructgpt-training} Ouyang, L., et al. \href{https://arxiv.org/abs/2203.02155}{Training language models to follow instructions with human feedback}. 2022.\par\url{https://arxiv.org/abs/2203.02155}

% Official sources verified 2026-10-12; insert inside the bibliography.
\bibitem{ops-docker-security} Docker. \href{https://docs.docker.com/engine/security/}{Docker Engine security}.\par\url{https://docs.docker.com/engine/security/}
\bibitem{ops-pytorch-install} PyTorch. \href{https://docs.pytorch.org/get-started/locally/}{Start Locally}.\par\url{https://docs.pytorch.org/get-started/locally/}
\bibitem{ops-nvidia-container} NVIDIA. \href{https://docs.nvidia.com/datacenter/cloud-native/container-toolkit/latest/install-guide.html}{Installing the NVIDIA Container Toolkit}.\par\url{https://docs.nvidia.com/datacenter/cloud-native/container-toolkit/latest/install-guide.html}
\bibitem{ops-nvidia-legacy} NVIDIA. \href{https://developer.nvidia.com/cuda/gpus/legacy}{Legacy CUDA GPU Compute Capability}.\par\url{https://developer.nvidia.com/cuda/gpus/legacy}
\bibitem{ops-vllm-gpu} vLLM. \href{https://docs.vllm.ai/en/latest/getting_started/installation/gpu/}{GPU installation requirements}.\par\url{https://docs.vllm.ai/en/latest/getting_started/installation/gpu/}
\bibitem{ops-cuda13} NVIDIA. \href{https://docs.nvidia.com/cuda/archive/13.0.0/cuda-toolkit-release-notes/}{CUDA Toolkit 13.0 Release Notes，Deprecated Architectures}.\par\url{https://docs.nvidia.com/cuda/archive/13.0.0/cuda-toolkit-release-notes/}
\bibitem{ops-llama-cmake} ggml-org. \href{https://github.com/ggml-org/llama.cpp/blob/master/ggml/src/ggml-cuda/CMakeLists.txt}{llama.cpp CUDA build configuration}.\par\url{https://github.com/ggml-org/llama.cpp/blob/master/ggml/src/ggml-cuda/CMakeLists.txt}
\bibitem{ops-llama-build} ggml-org. \href{https://github.com/ggml-org/llama.cpp/blob/master/docs/build.md}{llama.cpp build instructions}.\par\url{https://github.com/ggml-org/llama.cpp/blob/master/docs/build.md}
\bibitem{ops-llama-server} ggml-org. \href{https://github.com/ggml-org/llama.cpp/blob/master/tools/server/README.md}{llama.cpp HTTP server documentation}.\par\url{https://github.com/ggml-org/llama.cpp/blob/master/tools/server/README.md}
\bibitem{ops-agents-tracing} OpenAI. \href{https://openai.github.io/openai-agents-python/tracing/}{Agents SDK Tracing}.\par\url{https://openai.github.io/openai-agents-python/tracing/}
\bibitem{ops-openai-data} OpenAI. \href{https://developers.openai.com/api/docs/guides/your-data}{Data controls in the OpenAI platform}.\par\url{https://developers.openai.com/api/docs/guides/your-data}
\bibitem{ops-postgres-backup} PostgreSQL. \href{https://www.postgresql.org/docs/17/backup-dump.html}{PostgreSQL 17: SQL Dump}.\par\url{https://www.postgresql.org/docs/17/backup-dump.html}

\bibitem{ch5-openai-install} OpenAI. \href{https://developers.openai.com/api/docs/guides/agents/quickstart}{Agents SDK Quickstart}；\href{https://developers.openai.com/api/docs/guides/agents/sdk}{Agents SDK Overview}；\href{https://github.com/openai/openai-agents-python/blob/main/pyproject.toml}{Python SDK project metadata}. 访问日期：2026-10-12。\par\url{https://developers.openai.com/api/docs/guides/agents/quickstart}\par\url{https://github.com/openai/openai-agents-python/blob/main/pyproject.toml}
\bibitem{ch5-uv-project} Astral. \href{https://docs.astral.sh/uv/getting-started/installation/}{Installing uv}；\href{https://docs.astral.sh/uv/guides/projects/}{Working on projects}；\href{https://docs.astral.sh/uv/concepts/projects/sync/}{Locking and syncing}. 访问日期：2026-10-12。\par\url{https://docs.astral.sh/uv/getting-started/installation/}\par\url{https://docs.astral.sh/uv/guides/projects/}\par\url{https://docs.astral.sh/uv/concepts/projects/sync/}
\bibitem{ch5-docker-mac} Docker. \href{https://docs.docker.com/desktop/setup/install/mac-install/}{Install Docker Desktop on Mac}. 访问日期：2026-10-12。\par\url{https://docs.docker.com/desktop/setup/install/mac-install/}
\bibitem{ch5-postgres-image} Docker Official Images. \href{https://hub.docker.com/_/postgres}{PostgreSQL Official Image：architectures, environment variables, secrets and data volume}. 访问日期：2026-10-12。\par\url{https://hub.docker.com/_/postgres}

\bibitem{rdRuntimeSdk} OpenAI. Agents SDK. \url{https://developers.openai.com/api/docs/guides/agents/sdk}（访问：2026-10-13）.
\bibitem{rdRuntimePg} PostgreSQL Global Development Group. PostgreSQL 17: SELECT, The Locking Clause. \url{https://www.postgresql.org/docs/17/sql-select.html}（访问：2026-10-12）.
\bibitem{rdRuntimeModels} OpenAI. Models and providers. \url{https://developers.openai.com/api/docs/guides/agents/models}（访问：2026-10-12）.
\bibitem{rdRuntimeAgent} OpenAI. Agent definitions. \url{https://developers.openai.com/api/docs/guides/agents/define-agents}（访问：2026-10-12）.
\bibitem{rdRuntimeOrchestration} OpenAI. Orchestration and handoffs. \url{https://developers.openai.com/api/docs/guides/agents/orchestration}（访问：2026-10-12）.
\bibitem{rdRuntimeResults} OpenAI. Results and state. \url{https://developers.openai.com/api/docs/guides/agents/results}（访问：2026-10-12）.
\bibitem{rdRuntimeRun} OpenAI. Running agents. \url{https://developers.openai.com/api/docs/guides/agents/running-agents}（访问：2026-10-12）.
\bibitem{rdRuntimeApproval} OpenAI. Guardrails and human review. \url{https://developers.openai.com/api/docs/guides/agents/guardrails-approvals}（访问：2026-10-12）.
\bibitem{rdRuntimePersistent} OpenAI. Python Agents SDK: Human-in-the-loop, Long-running approvals. \url{https://openai.github.io/openai-agents-python/human_in_the_loop/}（访问：2026-10-12）.

\bibitem{ch5-github-git} GitHub Docs. \href{https://docs.github.com/en/get-started/using-git/about-git}{About Git}. 访问日期：2026-10-12。\par\url{https://docs.github.com/en/get-started/using-git/about-git}
\bibitem{ch5-aliyun-monitoring} Alibaba Cloud. \href{https://www.alibabacloud.com/help/en/cms/cloudmonitor-1-0/user-guide/host-monitoring/}{CloudMonitor: Host monitoring}. 访问日期：2026-10-12。\par\url{https://www.alibabacloud.com/help/en/cms/cloudmonitor-1-0/user-guide/host-monitoring/}

\bibitem{budget-openai-pricing} OpenAI. \href{https://developers.openai.com/api/docs/pricing}{API Pricing：标准模型与工具费用}. 访问日期：2026-10-12。\par\url{https://developers.openai.com/api/docs/pricing}
\bibitem{budget-mini} OpenAI. \href{https://developers.openai.com/api/docs/models/gpt-5.4-mini}{GPT-5.4 mini：模型与价格}. 访问日期：2026-10-12。\par\url{https://developers.openai.com/api/docs/models/gpt-5.4-mini}
\bibitem{budget-gpt54} OpenAI. \href{https://developers.openai.com/api/docs/models/gpt-5.4}{GPT-5.4：价格、长上下文与区域处理条件}. 访问日期：2026-10-12。\par\url{https://developers.openai.com/api/docs/models/gpt-5.4}
\bibitem{budget-chatgpt-api} OpenAI Help Center. \href{https://help.openai.com/en/articles/9039756-managing-billing-for-chatgpt-and-the-api-platform}{Managing billing for ChatGPT and the API platform}. 访问日期：2026-10-12。\par{\def\UrlBreaks{\do\/\do-\do.}\url{https://help.openai.com/en/articles/9039756-managing-billing-for-chatgpt-and-the-api-platform}}

\bibitem{budget-github-plans} GitHub Docs. \href{https://docs.github.com/en/get-started/learning-about-github/githubs-plans}{GitHub's plans}. 访问日期：2026-10-12。\par\url{https://docs.github.com/en/get-started/learning-about-github/githubs-plans}
\bibitem{budget-github-actions} GitHub Docs. \href{https://docs.github.com/en/billing/concepts/product-billing/github-actions}{GitHub Actions billing}. 访问日期：2026-10-12。\par\url{https://docs.github.com/en/billing/concepts/product-billing/github-actions}
\bibitem{budget-github-lfs} GitHub Docs. \href{https://docs.github.com/en/billing/concepts/product-billing/git-lfs}{Git Large File Storage billing}. 访问日期：2026-10-12。\par\url{https://docs.github.com/en/billing/concepts/product-billing/git-lfs}
\bibitem{budget-docker-license} Docker. \href{https://docs.docker.com/subscription-billing/desktop-license/}{Docker Desktop license agreement}. 访问日期：2026-10-12。\par\url{https://docs.docker.com/subscription-billing/desktop-license/}
\bibitem{budget-sdk-license} OpenAI. \href{https://github.com/openai/openai-agents-python/blob/main/LICENSE}{OpenAI Agents SDK Python MIT License}. 访问日期：2026-10-12。\par\url{https://github.com/openai/openai-agents-python/blob/main/LICENSE}

\bibitem{budget-reasoning} OpenAI. Reasoning models. 访问日期：2026-10-12。\par\url{https://developers.openai.com/api/docs/guides/reasoning}
\bibitem{v121-codex-auth} OpenAI. Codex Authentication. 访问日期：2026-10-13。\par\url{https://learn.chatgpt.com/docs/auth}
\bibitem{v121-codex-exec} OpenAI. Codex Non-interactive mode. 访问日期：2026-10-13。\par\url{https://learn.chatgpt.com/docs/non-interactive-mode}
\bibitem{v121-plan-usage} OpenAI. Sign in with ChatGPT and plan usage. 访问日期：2026-10-13。\par\url{https://learn.chatgpt.com/docs/sign-in-with-chatgpt}
\bibitem{v121-plan-quickstart} OpenAI. Sign in with ChatGPT Quickstart. 访问日期：2026-10-13。\par\url{https://developers.openai.com/siwc/quickstart}
\end{thebibliography}
\endgroup


\end{document}
